% !Mode:: "TeX:UTF-8"
% ============================================================
% ANSWERS TO THE EXERCISES
% (checked automatically: all answers were computed with SymPy)
% ============================================================
\chapter*{Answers to the Exercises}
\addcontentsline{toc}{chapter}{Answers to the Exercises}
\markboth{ANSWERS TO THE EXERCISES}{ANSWERS TO THE EXERCISES}

\small
The answers are arranged by section, exercise set and exercise number.
Inverse matrices are written with the common denominator taken outside
the matrix. Approximate values are rounded to four decimal places.

\section*{Chapter 1. Matrices and Determinants}
% !Mode:: "TeX:UTF-8"
\javobmashgulot{1}{The concept of a matrix}
\par\noindent\emph{Exercise set 1.}\nopagebreak
\begin{multicols}{2}\raggedright
\begin{enumerate}[label=\arabic*., leftmargin=*, itemsep=1pt]
\item $\left(\begin{smallmatrix}-13 & -51\\ 68 & 4\end{smallmatrix}\right)$
\item $\left(\begin{smallmatrix}-21 & -15\\ 10 & 9\end{smallmatrix}\right)$
\item $\left(\begin{smallmatrix}-15 & -2\\ 4 & -8\end{smallmatrix}\right)$
\item $\left(\begin{smallmatrix}-21 & -2\\ -8 & -5\end{smallmatrix}\right)$
\item $\left(\begin{smallmatrix}28 & 32\\ -24 & 44\end{smallmatrix}\right)$
\item $\left(\begin{smallmatrix}-3 & -1\\ 3 & -8\end{smallmatrix}\right)$
\item $\left(\begin{smallmatrix}-13 & -12\\ -6 & -7\end{smallmatrix}\right)$
\item $\left(\begin{smallmatrix}-2 & -8\\ -16 & -10\end{smallmatrix}\right)$
\item $\left(\begin{smallmatrix}8 & 21\\ 0 & -20\end{smallmatrix}\right)$
\item $\left(\begin{smallmatrix}21 & 22\\ -33 & 10\end{smallmatrix}\right)$
\item $\left(\begin{smallmatrix}2 & -3\\ -9 & -4\end{smallmatrix}\right)$
\item $\left(\begin{smallmatrix}-15 & 6\\ 8 & -9\end{smallmatrix}\right)$
\item $\left(\begin{smallmatrix}14 & -15\\ 60 & -1\end{smallmatrix}\right)$
\item $\left(\begin{smallmatrix}-19 & 12\\ 3 & -10\end{smallmatrix}\right)$
\item $\left(\begin{smallmatrix}21 & 6\\ -6 & 13\end{smallmatrix}\right)$
\item $\left(\begin{smallmatrix}-53 & 24\\ 48 & -29\end{smallmatrix}\right)$
\item $\left(\begin{smallmatrix}-18 & -26\\ -39 & -31\end{smallmatrix}\right)$
\item $\left(\begin{smallmatrix}-26 & 30\\ -20 & 14\end{smallmatrix}\right)$
\item $\left(\begin{smallmatrix}-3 & 0\\ 0 & -3\end{smallmatrix}\right)$
\item $\left(\begin{smallmatrix}-3 & -8\\ -8 & 5\end{smallmatrix}\right)$
\item $\left(\begin{smallmatrix}10 & -20\\ -5 & -15\end{smallmatrix}\right)$
\item $\left(\begin{smallmatrix}-26 & 26\\ -104 & -26\end{smallmatrix}\right)$
\item $\left(\begin{smallmatrix}-27 & -4\\ -3 & -29\end{smallmatrix}\right)$
\item $\left(\begin{smallmatrix}-5 & -7\\ -7 & -26\end{smallmatrix}\right)$
\item $\left(\begin{smallmatrix}-5 & -6\\ -6 & 1\end{smallmatrix}\right)$
\item $\left(\begin{smallmatrix}-3 & 0\\ 9 & -12\end{smallmatrix}\right)$
\item $\left(\begin{smallmatrix}7 & -3\\ 6 & 7\end{smallmatrix}\right)$
\item $\left(\begin{smallmatrix}-15 & 28\\ 7 & -29\end{smallmatrix}\right)$
\item $\left(\begin{smallmatrix}-2 & -1\\ -1 & 1\end{smallmatrix}\right)$
\item $\left(\begin{smallmatrix}-22 & 0\\ 0 & -22\end{smallmatrix}\right)$
\end{enumerate}
\end{multicols}
\par\noindent\emph{Exercise set 2.}\nopagebreak
\begin{enumerate}[label=\arabic*., leftmargin=*, itemsep=1pt]
\item $A+A^T=\left(\begin{smallmatrix}-4 & 0 & 11\\ 0 & -2 & 6\\ 11 & 6 & -2\end{smallmatrix}\right)$
\item $A+A^T=\left(\begin{smallmatrix}-10 & -8\\ -8 & -6\end{smallmatrix}\right)$
\item $A+A^T=\left(\begin{smallmatrix}-6 & 0 & -9\\ 0 & 6 & 4\\ -9 & 4 & 8\end{smallmatrix}\right)$
\item $A+A^T=\left(\begin{smallmatrix}12 & -1\\ -1 & 8\end{smallmatrix}\right)$
\item $A+A^T=\left(\begin{smallmatrix}0 & 11\\ 11 & 12\end{smallmatrix}\right)$
\item $A+A^T=\left(\begin{smallmatrix}2 & -4\\ -4 & 12\end{smallmatrix}\right)$
\item $A+A^T=\left(\begin{smallmatrix}-10 & 11\\ 11 & 0\end{smallmatrix}\right)$
\item $A+A^T=\left(\begin{smallmatrix}0 & 10 & -9\\ 10 & -8 & -10\\ -9 & -10 & -8\end{smallmatrix}\right)$
\item $A+A^T=\left(\begin{smallmatrix}2 & 2 & 5\\ 2 & 6 & 7\\ 5 & 7 & -2\end{smallmatrix}\right)$
\item $A+A^T=\left(\begin{smallmatrix}4 & -2\\ -2 & -12\end{smallmatrix}\right)$
\item $A+A^T=\left(\begin{smallmatrix}12 & 9\\ 9 & -10\end{smallmatrix}\right)$
\item $A+A^T=\left(\begin{smallmatrix}10 & -7 & -2\\ -7 & -6 & -9\\ -2 & -9 & -4\end{smallmatrix}\right)$
\item $A+A^T=\left(\begin{smallmatrix}-6 & 5 & 3\\ 5 & -4 & -2\\ 3 & -2 & -12\end{smallmatrix}\right)$
\item $A+A^T=\left(\begin{smallmatrix}2 & 7\\ 7 & 4\end{smallmatrix}\right)$
\item $A+A^T=\left(\begin{smallmatrix}4 & -8 & -4\\ -8 & 4 & 3\\ -4 & 3 & 12\end{smallmatrix}\right)$
\item $A+A^T=\left(\begin{smallmatrix}6 & 0\\ 0 & 12\end{smallmatrix}\right)$
\item $A+A^T=\left(\begin{smallmatrix}-8 & -3\\ -3 & 6\end{smallmatrix}\right)$
\item $A+A^T=\left(\begin{smallmatrix}4 & -7\\ -7 & 8\end{smallmatrix}\right)$
\item $A+A^T=\left(\begin{smallmatrix}4 & 8 & 3\\ 8 & 12 & -11\\ 3 & -11 & -6\end{smallmatrix}\right)$
\item $A+A^T=\left(\begin{smallmatrix}-4 & 0\\ 0 & -10\end{smallmatrix}\right)$
\item $A+A^T=\left(\begin{smallmatrix}2 & 8 & -7\\ 8 & -10 & 4\\ -7 & 4 & 4\end{smallmatrix}\right)$
\item $A+A^T=\left(\begin{smallmatrix}4 & -5 & 7\\ -5 & 2 & 8\\ 7 & 8 & 2\end{smallmatrix}\right)$
\item $A+A^T=\left(\begin{smallmatrix}-6 & 3 & 3\\ 3 & 4 & -7\\ 3 & -7 & 0\end{smallmatrix}\right)$
\item $A+A^T=\left(\begin{smallmatrix}0 & 0\\ 0 & -10\end{smallmatrix}\right)$
\item $A+A^T=\left(\begin{smallmatrix}0 & -8\\ -8 & 8\end{smallmatrix}\right)$
\item $A+A^T=\left(\begin{smallmatrix}12 & 1\\ 1 & -8\end{smallmatrix}\right)$
\item $A+A^T=\left(\begin{smallmatrix}-8 & -6\\ -6 & 2\end{smallmatrix}\right)$
\item $A+A^T=\left(\begin{smallmatrix}10 & -5\\ -5 & 2\end{smallmatrix}\right)$
\item $A+A^T=\left(\begin{smallmatrix}8 & -7\\ -7 & 10\end{smallmatrix}\right)$
\item $A+A^T=\left(\begin{smallmatrix}4 & 0 & -2\\ 0 & -6 & -6\\ -2 & -6 & 10\end{smallmatrix}\right)$
\end{enumerate}
\par\noindent\emph{Exercise set 3.}\nopagebreak
\begin{multicols}{4}\raggedright
\begin{enumerate}[label=\arabic*., leftmargin=*, itemsep=1pt]
\item $Z=-4 X+2 Y$
\item $Z=-3 X- Y$
\item $Z=-2 X+2 Y$
\item $Z=3 X+ Y$
\item $Z=-4 X-3 Y$
\item $Z=-3 X+3 Y$
\item $Z=-2 X-4 Y$
\item $Z=-2 X- Y$
\item $Z=3 X+4 Y$
\item $Z=-4 X-4 Y$
\item $Z=- X+3 Y$
\item $Z=4 X+2 Y$
\item $Z=-2 X+2 Y$
\item $Z=2 X-2 Y$
\item $Z=- X-4 Y$
\item $Z= X+ Y$
\item $Z= X-2 Y$
\item $Z=4 X- Y$
\item $Z=-2 X+2 Y$
\item $Z=- X-3 Y$
\item $Z= X+3 Y$
\item $Z=4 X+2 Y$
\item $Z=- X-3 Y$
\item $Z=2 X+3 Y$
\item $Z=- X+3 Y$
\item $Z=-3 X+4 Y$
\item $Z=- X- Y$
\item $Z=-4 X- Y$
\item $Z=-2 X-4 Y$
\item $Z=- X+3 Y$
\end{enumerate}
\end{multicols}
\par\noindent\emph{Exercise set 4.}\nopagebreak
\begin{multicols}{2}\raggedright
\begin{enumerate}[label=\arabic*., leftmargin=*, itemsep=1pt]
\item $a=9,\ b=9$
\item $a=3,\ b=9$
\item $a=9,\ b=6$
\item $a=3,\ b=4$
\item $a=6,\ b=3$
\item $a=8,\ b=2$
\item $a=9,\ b=8$
\item $a=8,\ b=7$
\item $a=7,\ b=8$
\item $a=6,\ b=7$
\item $a=3,\ b=7$
\item $a=3,\ b=2$
\item $a=5,\ b=6$
\item $a=7,\ b=8$
\item $a=5,\ b=3$
\item $a=4,\ b=6$
\item $a=4,\ b=9$
\item $a=6,\ b=6$
\item $a=9,\ b=8$
\item $a=3,\ b=5$
\item $a=9,\ b=7$
\item $a=5,\ b=5$
\item $a=3,\ b=7$
\item $a=5,\ b=2$
\item $a=5,\ b=8$
\item $a=6,\ b=7$
\item $a=5,\ b=3$
\item $a=9,\ b=8$
\item $a=8,\ b=9$
\item $a=8,\ b=9$
\end{enumerate}
\end{multicols}
\par\noindent\emph{Exercise set 5.}\nopagebreak
\begin{enumerate}[label=\arabic*., leftmargin=*, itemsep=1pt]
\item 12 sizes: $1\times72,\ 2\times36,\ 3\times24,\ 4\times18,\ 6\times12,\ 8\times9,\ 9\times8,\ 12\times6,\ 18\times4,\ 24\times3,\ 36\times2,\ 72\times1$
\item 9 sizes: $1\times36,\ 2\times18,\ 3\times12,\ 4\times9,\ 6\times6,\ 9\times4,\ 12\times3,\ 18\times2,\ 36\times1$
\item 8 sizes: $1\times30,\ 2\times15,\ 3\times10,\ 5\times6,\ 6\times5,\ 10\times3,\ 15\times2,\ 30\times1$
\item 8 sizes: $1\times42,\ 2\times21,\ 3\times14,\ 6\times7,\ 7\times6,\ 14\times3,\ 21\times2,\ 42\times1$
\item 7 sizes: $1\times64,\ 2\times32,\ 4\times16,\ 8\times8,\ 16\times4,\ 32\times2,\ 64\times1$
\item 6 sizes: $1\times28,\ 2\times14,\ 4\times7,\ 7\times4,\ 14\times2,\ 28\times1$
\item 8 sizes: $1\times88,\ 2\times44,\ 4\times22,\ 8\times11,\ 11\times8,\ 22\times4,\ 44\times2,\ 88\times1$
\item 6 sizes: $1\times44,\ 2\times22,\ 4\times11,\ 11\times4,\ 22\times2,\ 44\times1$
\item 10 sizes: $1\times48,\ 2\times24,\ 3\times16,\ 4\times12,\ 6\times8,\ 8\times6,\ 12\times4,\ 16\times3,\ 24\times2,\ 48\times1$
\item 6 sizes: $1\times75,\ 3\times25,\ 5\times15,\ 15\times5,\ 25\times3,\ 75\times1$
\item 8 sizes: $1\times54,\ 2\times27,\ 3\times18,\ 6\times9,\ 9\times6,\ 18\times3,\ 27\times2,\ 54\times1$
\item 8 sizes: $1\times40,\ 2\times20,\ 4\times10,\ 5\times8,\ 8\times5,\ 10\times4,\ 20\times2,\ 40\times1$
\item 8 sizes: $1\times66,\ 2\times33,\ 3\times22,\ 6\times11,\ 11\times6,\ 22\times3,\ 33\times2,\ 66\times1$
\item 6 sizes: $1\times63,\ 3\times21,\ 7\times9,\ 9\times7,\ 21\times3,\ 63\times1$
\item 6 sizes: $1\times45,\ 3\times15,\ 5\times9,\ 9\times5,\ 15\times3,\ 45\times1$
\item 6 sizes: $1\times12,\ 2\times6,\ 3\times4,\ 4\times3,\ 6\times2,\ 12\times1$
\item 5 sizes: $1\times16,\ 2\times8,\ 4\times4,\ 8\times2,\ 16\times1$
\item 12 sizes: $1\times60,\ 2\times30,\ 3\times20,\ 4\times15,\ 5\times12,\ 6\times10,\ 10\times6,\ 12\times5,\ 15\times4,\ 20\times3,\ 30\times2,\ 60\times1$
\item 6 sizes: $1\times18,\ 2\times9,\ 3\times6,\ 6\times3,\ 9\times2,\ 18\times1$
\item 12 sizes: $1\times96,\ 2\times48,\ 3\times32,\ 4\times24,\ 6\times16,\ 8\times12,\ 12\times8,\ 16\times6,\ 24\times4,\ 32\times3,\ 48\times2,\ 96\times1$
\item 12 sizes: $1\times84,\ 2\times42,\ 3\times28,\ 4\times21,\ 6\times14,\ 7\times12,\ 12\times7,\ 14\times6,\ 21\times4,\ 28\times3,\ 42\times2,\ 84\times1$
\item 8 sizes: $1\times70,\ 2\times35,\ 5\times14,\ 7\times10,\ 10\times7,\ 14\times5,\ 35\times2,\ 70\times1$
\item 8 sizes: $1\times56,\ 2\times28,\ 4\times14,\ 7\times8,\ 8\times7,\ 14\times4,\ 28\times2,\ 56\times1$
\item 10 sizes: $1\times80,\ 2\times40,\ 4\times20,\ 5\times16,\ 8\times10,\ 10\times8,\ 16\times5,\ 20\times4,\ 40\times2,\ 80\times1$
\item 8 sizes: $1\times24,\ 2\times12,\ 3\times8,\ 4\times6,\ 6\times4,\ 8\times3,\ 12\times2,\ 24\times1$
\item 12 sizes: $1\times90,\ 2\times45,\ 3\times30,\ 5\times18,\ 6\times15,\ 9\times10,\ 10\times9,\ 15\times6,\ 18\times5,\ 30\times3,\ 45\times2,\ 90\times1$
\item 6 sizes: $1\times20,\ 2\times10,\ 4\times5,\ 5\times4,\ 10\times2,\ 20\times1$
\item 5 sizes: $1\times81,\ 3\times27,\ 9\times9,\ 27\times3,\ 81\times1$
\item 6 sizes: $1\times32,\ 2\times16,\ 4\times8,\ 8\times4,\ 16\times2,\ 32\times1$
\item 6 sizes: $1\times50,\ 2\times25,\ 5\times10,\ 10\times5,\ 25\times2,\ 50\times1$
\end{enumerate}
\par\noindent\emph{Exercise set 6.}\nopagebreak
\begin{enumerate}[label=\arabic*., leftmargin=*, itemsep=1pt]
\item $\left(\begin{smallmatrix}-1 & 9\\ 3 & -2\\ 0 & -4\end{smallmatrix}\right)$
\item $\left(\begin{smallmatrix}1 & -6 & 8\\ 1 & -5 & 3\\ 6 & -3 & -12\end{smallmatrix}\right)$
\item $\left(\begin{smallmatrix}1 & -1 & -15\\ -4 & 9 & -17\end{smallmatrix}\right)$
\item $\left(\begin{smallmatrix}-2 & -10 & 14\\ 4 & 14 & -8\end{smallmatrix}\right)$
\item $\left(\begin{smallmatrix}-9 & 13 & -2\\ 19 & 3 & 13\\ 2 & 15 & -4\end{smallmatrix}\right)$
\item $\left(\begin{smallmatrix}7 & -9\\ -13 & -6\\ -4 & 4\end{smallmatrix}\right)$
\item $\left(\begin{smallmatrix}-7 & -6\\ -5 & -2\end{smallmatrix}\right)$
\item $\left(\begin{smallmatrix}2 & 20\\ 16 & -8\\ 8 & 4\end{smallmatrix}\right)$
\item $\left(\begin{smallmatrix}1 & -4 & -6\\ -1 & -9 & 6\end{smallmatrix}\right)$
\item $\left(\begin{smallmatrix}3 & 11\\ -7 & -14\\ 12 & 16\end{smallmatrix}\right)$
\item $\left(\begin{smallmatrix}-10 & -15 & 19\\ -11 & -6 & 9\\ 2 & 9 & 31\end{smallmatrix}\right)$
\item $\left(\begin{smallmatrix}17 & -12 & -13\\ -7 & -8 & 19\end{smallmatrix}\right)$
\item $\left(\begin{smallmatrix}6 & -9\\ 6 & 2\end{smallmatrix}\right)$
\item $\left(\begin{smallmatrix}-4 & 6 & 5\\ 6 & -1 & -1\end{smallmatrix}\right)$
\item $\left(\begin{smallmatrix}10 & -13 & -2\\ 2 & -21 & -15\\ 20 & -26 & 16\end{smallmatrix}\right)$
\item $\left(\begin{smallmatrix}15 & 14\\ 10 & 2\end{smallmatrix}\right)$
\item $\left(\begin{smallmatrix}-15 & -10 & 10\\ 11 & 5 & -14\end{smallmatrix}\right)$
\item $\left(\begin{smallmatrix}-7 & -10 & 8\\ 7 & -9 & -12\end{smallmatrix}\right)$
\item $\left(\begin{smallmatrix}2 & 20\\ 6 & -24\\ 0 & -26\end{smallmatrix}\right)$
\item $\left(\begin{smallmatrix}-21 & -11 & 15\\ -15 & 11 & -3\end{smallmatrix}\right)$
\item $\left(\begin{smallmatrix}-4 & 16 & -2\\ -2 & 12 & -4\\ -10 & 4 & 6\end{smallmatrix}\right)$
\item $\left(\begin{smallmatrix}-1 & -13\\ 13 & 21\end{smallmatrix}\right)$
\item $\left(\begin{smallmatrix}22 & 16 & -2\\ -2 & 2 & 10\end{smallmatrix}\right)$
\item $\left(\begin{smallmatrix}-4 & 28 & -14\\ -12 & -8 & 20\\ -22 & 20 & -24\end{smallmatrix}\right)$
\item $\left(\begin{smallmatrix}-13 & 6 & 22\\ -1 & 16 & 6\end{smallmatrix}\right)$
\item $\left(\begin{smallmatrix}12 & 30 & -3\\ 21 & 12 & 3\\ 3 & -21 & -12\end{smallmatrix}\right)$
\item $\left(\begin{smallmatrix}-18 & -24 & -15\\ 9 & -3 & -3\end{smallmatrix}\right)$
\item $\left(\begin{smallmatrix}-20 & 26\\ -20 & -26\end{smallmatrix}\right)$
\item $\left(\begin{smallmatrix}-9 & -2 & 10\\ -14 & 2 & 5\end{smallmatrix}\right)$
\item $\left(\begin{smallmatrix}-2 & -1\\ 2 & 2\\ 5 & 0\end{smallmatrix}\right)$
\end{enumerate}
\par\noindent\emph{Exercise set 7.}\nopagebreak
\begin{enumerate}[label=\arabic*., leftmargin=*, itemsep=1pt]
\item $c=\left(\begin{smallmatrix}0\\25\end{smallmatrix}\right)$ ($Kp=(78;103)^T$)
\item $c=\left(\begin{smallmatrix}16\\3\end{smallmatrix}\right)$
\item $c=\left(\begin{smallmatrix}0\\19\end{smallmatrix}\right)$
\item $c=\left(\begin{smallmatrix}3\\15\end{smallmatrix}\right)$
\item $\det K=-2$, $\gcd(-2,26)=2$; since $\gcd(\det K,26)\ne1$, $K$ is not invertible modulo $26$ and cannot serve as a key
\item $\det K=0$; $K$ is singular (has no inverse), so it cannot serve as a key
\item $c=(1;0;1;1;0;1;0)$
\item $s=(0;0;1)^T\ne0$; the word contains an error. The syndrome equals the 4th column of $H$, so the 4th bit is wrong; the corrected word is $(1;0;1;0;1;0;1)$
\item $A^2=\left(\begin{smallmatrix}2&0&0&2&1\\0&2&2&0&1\\0&2&3&1&1\\2&0&1&3&1\\1&1&1&1&2\end{smallmatrix}\right)$; the most connected nodes are $S_3$ and $S_4$ (degree 3, the diagonal entries of $A^2$)
\item $K^2-3K+E=O$, hence $K(3E-K)=E$ and $K^{-1}=3E-K=\left(\begin{smallmatrix}1&-1\\-1&2\end{smallmatrix}\right)$
\end{enumerate}

% !Mode:: "TeX:UTF-8"
\javobmashgulot{2}{The concept of a determinant. Determinants of the second and third order}
\par\noindent\emph{Exercise set 1.}\nopagebreak
\begin{multicols}{5}\raggedright
\begin{enumerate}[label=\arabic*., leftmargin=*, itemsep=1pt]
\item $22$
\item $-11$
\item $20$
\item $42$
\item $32$
\item $16$
\item $35$
\item $-14$
\item $26$
\item $18$
\item $58$
\item $3$
\item $27$
\item $18$
\item $13$
\item $0$
\item $-41$
\item $84$
\item $28$
\item $-3$
\item $19$
\item $-43$
\item $-26$
\item $31$
\item $42$
\item $7$
\item $-37$
\item $35$
\item $-46$
\item $1$
\end{enumerate}
\end{multicols}
\par\noindent\emph{Exercise set 2.}\nopagebreak
\begin{multicols}{5}\raggedright
\begin{enumerate}[label=\arabic*., leftmargin=*, itemsep=1pt]
\item $6$
\item $83$
\item $17$
\item $38$
\item $65$
\item $90$
\item $-89$
\item $77$
\item $89$
\item $75$
\item $-32$
\item $-47$
\item $22$
\item $17$
\item $-52$
\item $-17$
\item $16$
\item $22$
\item $117$
\item $-129$
\item $39$
\item $51$
\item $71$
\item $7$
\item $20$
\item $49$
\item $87$
\item $24$
\item $66$
\item $-55$
\end{enumerate}
\end{multicols}
\par\noindent\emph{Exercise set 3.}\nopagebreak
\begin{multicols}{5}\raggedright
\begin{enumerate}[label=\arabic*., leftmargin=*, itemsep=1pt]
\item $-77$
\item $23$
\item $100$
\item $51$
\item $90$
\item $106$
\item $102$
\item $50$
\item $-79$
\item $13$
\item $-75$
\item $-70$
\item $83$
\item $93$
\item $5$
\item $-18$
\item $-31$
\item $63$
\item $73$
\item $-86$
\item $9$
\item $1$
\item $-36$
\item $24$
\item $63$
\item $27$
\item $-20$
\item $42$
\item $20$
\item $66$
\end{enumerate}
\end{multicols}
\par\noindent\emph{Exercise set 4.}\nopagebreak
\begin{multicols}{5}\raggedright
\begin{enumerate}[label=\arabic*., leftmargin=*, itemsep=1pt]
\item $0$
\item $0$
\item $-12$
\item $0$
\item $0$
\item $0$
\item $0$
\item $210$
\item $0$
\item $0$
\item $-15$
\item $0$
\item $0$
\item $0$
\item $0$
\item $0$
\item $0$
\item $0$
\item $-21$
\item $0$
\item $140$
\item $-42$
\item $0$
\item $0$
\item $-5$
\item $0$
\item $0$
\item $0$
\item $0$
\item $396$
\end{enumerate}
\end{multicols}
\par\noindent\emph{Exercise set 5.}\nopagebreak
\begin{multicols}{2}\raggedright
\begin{enumerate}[label=\arabic*., leftmargin=*, itemsep=1pt]
\item $\det K=11,\ \gcd(11,26)=1$ – admissible
\item $\det K=0,\ \gcd(0,26)=26$ – not admissible
\item $\det K=1,\ \gcd(1,26)=1$ – admissible
\item $\det K=19,\ \gcd(19,26)=1$ – admissible
\item $\det K=43,\ \gcd(43,26)=1$ – admissible
\item $\det K=0,\ \gcd(0,26)=26$ – not admissible
\item $\det K=-1,\ \gcd(-1,26)=1$ – admissible
\item $\det K=19,\ \gcd(19,26)=1$ – admissible
\item $\det K=51,\ \gcd(51,26)=1$ – admissible
\item $\det K=1,\ \gcd(1,26)=1$ – admissible
\item $\det K=89,\ \gcd(89,26)=1$ – admissible
\item $\det K=-2,\ \gcd(-2,26)=2$ – not admissible
\item $\det K=9,\ \gcd(9,26)=1$ – admissible
\item $\det K=-11,\ \gcd(-11,26)=1$ – admissible
\item $\det K=1,\ \gcd(1,26)=1$ – admissible
\item $\det K=2,\ \gcd(2,26)=2$ – not admissible
\item $\det K=24,\ \gcd(24,26)=2$ – not admissible
\item $\det K=31,\ \gcd(31,26)=1$ – admissible
\item $\det K=0,\ \gcd(0,26)=26$ – not admissible
\item $\det K=172,\ \gcd(172,26)=2$ – not admissible
\item $\det K=1,\ \gcd(1,26)=1$ – admissible
\item $\det K=26,\ \gcd(26,26)=26$ – not admissible
\item $\det K=17,\ \gcd(17,26)=1$ – admissible
\item $\det K=22,\ \gcd(22,26)=2$ – not admissible
\item $\det K=1,\ \gcd(1,26)=1$ – admissible
\item $\det K=-3,\ \gcd(-3,26)=1$ – admissible
\item $\det K=1,\ \gcd(1,26)=1$ – admissible
\item $\det K=-1,\ \gcd(-1,26)=1$ – admissible
\item $\det K=-29,\ \gcd(-29,26)=1$ – admissible
\item $\det K=53,\ \gcd(53,26)=1$ – admissible
\end{enumerate}
\end{multicols}

% !Mode:: "TeX:UTF-8"
\javobmashgulot{3}{Minors and cofactors. Evaluation of determinants of order \texorpdfstring{$n$}{n}}
\par\noindent\emph{Exercise set 1.}\nopagebreak
\begin{multicols}{2}\raggedright
\begin{enumerate}[label=\arabic*., leftmargin=*, itemsep=1pt]
\item $M_{11}=22,\ A_{11}=22;\,M_{23}=7,\ A_{23}=-7$
\item $M_{12}=-11,\ A_{12}=11;\,M_{31}=-1,\ A_{31}=-1$
\item $M_{22}=11,\ A_{22}=11;\,M_{13}=-2,\ A_{13}=-2$
\item $M_{32}=-20,\ A_{32}=20;\,M_{11}=33,\ A_{11}=33$
\item $M_{21}=-4,\ A_{21}=4;\,M_{33}=-6,\ A_{33}=-6$
\item $M_{13}=-6,\ A_{13}=-6;\,M_{22}=-8,\ A_{22}=-8$
\item $M_{31}=15,\ A_{31}=15;\,M_{12}=-7,\ A_{12}=7$
\item $M_{23}=-7,\ A_{23}=7;\,M_{11}=5,\ A_{11}=5$
\item $M_{32}=-12,\ A_{32}=12;\,M_{21}=-4,\ A_{21}=4$
\item $M_{11}=23,\ A_{11}=23;\,M_{33}=24,\ A_{33}=24$
\item $M_{12}=18,\ A_{12}=-18;\,M_{31}=-12,\ A_{31}=-12$
\item $M_{22}=-9,\ A_{22}=-9;\,M_{13}=-10,\ A_{13}=-10$
\item $M_{21}=4,\ A_{21}=-4;\,M_{32}=2,\ A_{32}=-2$
\item $M_{13}=6,\ A_{13}=6;\,M_{23}=-4,\ A_{23}=4$
\item $M_{11}=-18,\ A_{11}=-18;\,M_{32}=-26,\ A_{32}=26$
\item $M_{31}=10,\ A_{31}=10;\,M_{12}=-8,\ A_{12}=8$
\item $M_{22}=19,\ A_{22}=19;\,M_{33}=3,\ A_{33}=3$
\item $M_{13}=2,\ A_{13}=2;\,M_{21}=-4,\ A_{21}=4$
\item $M_{23}=-11,\ A_{23}=11;\,M_{11}=12,\ A_{11}=12$
\item $M_{32}=-12,\ A_{32}=12;\,M_{12}=3,\ A_{12}=-3$
\item $M_{11}=6,\ A_{11}=6;\,M_{23}=-10,\ A_{23}=10$
\item $M_{21}=-4,\ A_{21}=4;\,M_{33}=12,\ A_{33}=12$
\item $M_{31}=22,\ A_{31}=22;\,M_{22}=-20,\ A_{22}=-20$
\item $M_{12}=17,\ A_{12}=-17;\,M_{33}=20,\ A_{33}=20$
\item $M_{13}=4,\ A_{13}=4;\,M_{21}=4,\ A_{21}=-4$
\item $M_{22}=5,\ A_{22}=5;\,M_{31}=8,\ A_{31}=8$
\item $M_{32}=-6,\ A_{32}=6;\,M_{11}=-8,\ A_{11}=-8$
\item $M_{23}=-9,\ A_{23}=9;\,M_{12}=-21,\ A_{12}=21$
\item $M_{11}=-6,\ A_{11}=-6;\,M_{13}=4,\ A_{13}=4$
\item $M_{21}=4,\ A_{21}=-4;\,M_{32}=8,\ A_{32}=-8$
\end{enumerate}
\end{multicols}
\par\noindent\emph{Exercise set 2.}\nopagebreak
\begin{multicols}{5}\raggedright
\begin{enumerate}[label=\arabic*., leftmargin=*, itemsep=1pt]
\item $12$
\item $-79$
\item $24$
\item $-62$
\item $38$
\item $-36$
\item $66$
\item $-62$
\item $46$
\item $-49$
\item $70$
\item $62$
\item $-46$
\item $-54$
\item $-58$
\item $80$
\item $-64$
\item $-58$
\item $-30$
\item $-64$
\item $32$
\item $-62$
\item $85$
\item $62$
\item $53$
\item $56$
\item $-68$
\item $-39$
\item $14$
\item $32$
\end{enumerate}
\end{multicols}
\par\noindent\emph{Exercise set 3.}\nopagebreak
\begin{multicols}{5}\raggedright
\begin{enumerate}[label=\arabic*., leftmargin=*, itemsep=1pt]
\item $-10$
\item $-65$
\item $21$
\item $-8$
\item $-9$
\item $-134$
\item $92$
\item $114$
\item $110$
\item $83$
\item $1$
\item $212$
\item $-198$
\item $35$
\item $-160$
\item $-18$
\item $-52$
\item $28$
\item $-56$
\item $6$
\item $152$
\item $49$
\item $23$
\item $63$
\item $154$
\item $140$
\item $110$
\item $91$
\item $17$
\item $0$
\end{enumerate}
\end{multicols}
\par\noindent\emph{Exercise set 4.}\nopagebreak
\begin{multicols}{5}\raggedright
\begin{enumerate}[label=\arabic*., leftmargin=*, itemsep=1pt]
\item $79$ and $0$
\item $-70$ and $0$
\item $-60$ and $0$
\item $-90$ and $0$
\item $-54$ and $0$
\item $60$ and $0$
\item $-38$ and $0$
\item $-1$ and $0$
\item $50$ and $0$
\item $60$ and $0$
\item $7$ and $0$
\item $35$ and $0$
\item $64$ and $0$
\item $-41$ and $0$
\item $46$ and $0$
\item $75$ and $0$
\item $10$ and $0$
\item $48$ and $0$
\item $3$ and $0$
\item $45$ and $0$
\item $84$ and $0$
\item $23$ and $0$
\item $102$ and $0$
\item $49$ and $0$
\item $54$ and $0$
\item $21$ and $0$
\item $-40$ and $0$
\item $-5$ and $0$
\item $-55$ and $0$
\item $10$ and $0$
\end{enumerate}
\end{multicols}

% !Mode:: "TeX:UTF-8"
\javobmashgulot{4}{Computing the rank of a matrix. Singular and nonsingular matrices}
\par\noindent\emph{Exercise set 1.}\nopagebreak
\begin{multicols}{5}\raggedright
\begin{enumerate}[label=\arabic*., leftmargin=*, itemsep=1pt]
\item $\operatorname{rg}=2$
\item $\operatorname{rg}=2$
\item $\operatorname{rg}=1$
\item $\operatorname{rg}=2$
\item $\operatorname{rg}=2$
\item $\operatorname{rg}=3$
\item $\operatorname{rg}=1$
\item $\operatorname{rg}=2$
\item $\operatorname{rg}=1$
\item $\operatorname{rg}=2$
\item $\operatorname{rg}=1$
\item $\operatorname{rg}=2$
\item $\operatorname{rg}=2$
\item $\operatorname{rg}=2$
\item $\operatorname{rg}=3$
\item $\operatorname{rg}=1$
\item $\operatorname{rg}=2$
\item $\operatorname{rg}=3$
\item $\operatorname{rg}=1$
\item $\operatorname{rg}=2$
\item $\operatorname{rg}=1$
\item $\operatorname{rg}=2$
\item $\operatorname{rg}=2$
\item $\operatorname{rg}=1$
\item $\operatorname{rg}=2$
\item $\operatorname{rg}=2$
\item $\operatorname{rg}=1$
\item $\operatorname{rg}=3$
\item $\operatorname{rg}=1$
\item $\operatorname{rg}=2$
\end{enumerate}
\end{multicols}
\par\noindent\emph{Exercise set 2.}\nopagebreak
\begin{multicols}{5}\raggedright
\begin{enumerate}[label=\arabic*., leftmargin=*, itemsep=1pt]
\item $\operatorname{rg}=2$
\item $\operatorname{rg}=2$
\item $\operatorname{rg}=1$
\item $\operatorname{rg}=2$
\item $\operatorname{rg}=2$
\item $\operatorname{rg}=1$
\item $\operatorname{rg}=2$
\item $\operatorname{rg}=2$
\item $\operatorname{rg}=2$
\item $\operatorname{rg}=1$
\item $\operatorname{rg}=4$
\item $\operatorname{rg}=2$
\item $\operatorname{rg}=1$
\item $\operatorname{rg}=2$
\item $\operatorname{rg}=2$
\item $\operatorname{rg}=2$
\item $\operatorname{rg}=2$
\item $\operatorname{rg}=1$
\item $\operatorname{rg}=3$
\item $\operatorname{rg}=4$
\item $\operatorname{rg}=1$
\item $\operatorname{rg}=1$
\item $\operatorname{rg}=2$
\item $\operatorname{rg}=3$
\item $\operatorname{rg}=2$
\item $\operatorname{rg}=1$
\item $\operatorname{rg}=2$
\item $\operatorname{rg}=3$
\item $\operatorname{rg}=2$
\item $\operatorname{rg}=1$
\end{enumerate}
\end{multicols}
\par\noindent\emph{Exercise set 3.}\nopagebreak
\begin{multicols}{4}\raggedright
\begin{enumerate}[label=\arabic*., leftmargin=*, itemsep=1pt]
\item $\det=-2$, nonsingular
\item $\det=0$, singular
\item $\det=13$, nonsingular
\item $\det=0$, singular
\item $\det=-8$, nonsingular
\item $\det=0$, singular
\item $\det=0$, singular
\item $\det=0$, singular
\item $\det=0$, singular
\item $\det=60$, nonsingular
\item $\det=1$, nonsingular
\item $\det=0$, singular
\item $\det=0$, singular
\item $\det=1$, nonsingular
\item $\det=0$, singular
\item $\det=0$, singular
\item $\det=-30$, nonsingular
\item $\det=0$, singular
\item $\det=1$, nonsingular
\item $\det=0$, singular
\item $\det=0$, singular
\item $\det=-6$, nonsingular
\item $\det=0$, singular
\item $\det=0$, singular
\item $\det=0$, singular
\item $\det=4$, nonsingular
\item $\det=0$, singular
\item $\det=0$, singular
\item $\det=0$, singular
\item $\det=0$, singular
\end{enumerate}
\end{multicols}
\par\noindent\emph{Exercise set 4.}\nopagebreak
\begin{multicols}{2}\raggedright
\begin{enumerate}[label=\arabic*., leftmargin=*, itemsep=1pt]
\item $\operatorname{rg}A=\operatorname{rg}\bar A=2$; consistent (infinitely many solutions)
\item $\operatorname{rg}A=1,\ \operatorname{rg}\bar A=2$; inconsistent
\item $\operatorname{rg}A=\operatorname{rg}\bar A=2$; consistent (infinitely many solutions)
\item $\operatorname{rg}A=\operatorname{rg}\bar A=1$; consistent (infinitely many solutions)
\item $\operatorname{rg}A=\operatorname{rg}\bar A=2$; consistent (infinitely many solutions)
\item $\operatorname{rg}A=1,\ \operatorname{rg}\bar A=2$; inconsistent
\item $\operatorname{rg}A=\operatorname{rg}\bar A=2$; consistent (infinitely many solutions)
\item $\operatorname{rg}A=1,\ \operatorname{rg}\bar A=2$; inconsistent
\item $\operatorname{rg}A=\operatorname{rg}\bar A=2$; consistent (unique solution)
\item $\operatorname{rg}A=\operatorname{rg}\bar A=2$; consistent (infinitely many solutions)
\item $\operatorname{rg}A=\operatorname{rg}\bar A=1$; consistent (infinitely many solutions)
\item $\operatorname{rg}A=1,\ \operatorname{rg}\bar A=2$; inconsistent
\item $\operatorname{rg}A=\operatorname{rg}\bar A=1$; consistent (infinitely many solutions)
\item $\operatorname{rg}A=\operatorname{rg}\bar A=2$; consistent (infinitely many solutions)
\item $\operatorname{rg}A=\operatorname{rg}\bar A=1$; consistent (infinitely many solutions)
\item $\operatorname{rg}A=1,\ \operatorname{rg}\bar A=2$; inconsistent
\item $\operatorname{rg}A=\operatorname{rg}\bar A=3$; consistent (unique solution)
\item $\operatorname{rg}A=\operatorname{rg}\bar A=1$; consistent (infinitely many solutions)
\item $\operatorname{rg}A=\operatorname{rg}\bar A=2$; consistent (unique solution)
\item $\operatorname{rg}A=\operatorname{rg}\bar A=3$; consistent (infinitely many solutions)
\item $\operatorname{rg}A=\operatorname{rg}\bar A=1$; consistent (infinitely many solutions)
\item $\operatorname{rg}A=\operatorname{rg}\bar A=1$; consistent (infinitely many solutions)
\item $\operatorname{rg}A=\operatorname{rg}\bar A=2$; consistent (infinitely many solutions)
\item $\operatorname{rg}A=\operatorname{rg}\bar A=2$; consistent (unique solution)
\item $\operatorname{rg}A=\operatorname{rg}\bar A=1$; consistent (infinitely many solutions)
\item $\operatorname{rg}A=1,\ \operatorname{rg}\bar A=2$; inconsistent
\item $\operatorname{rg}A=2,\ \operatorname{rg}\bar A=3$; inconsistent
\item $\operatorname{rg}A=\operatorname{rg}\bar A=1$; consistent (infinitely many solutions)
\item $\operatorname{rg}A=1,\ \operatorname{rg}\bar A=2$; inconsistent
\item $\operatorname{rg}A=2,\ \operatorname{rg}\bar A=3$; inconsistent
\end{enumerate}
\end{multicols}

% !Mode:: "TeX:UTF-8"
\javobmashgulot{5}{The inverse matrix}
\par\noindent\emph{Exercise set 1.}\nopagebreak
\begin{multicols}{2}\raggedright
\begin{enumerate}[label=\arabic*., leftmargin=*, itemsep=1pt]
\item $\left(\begin{smallmatrix}2 & -1\\ -5 & 3\end{smallmatrix}\right)$
\item does not exist ($\det=0$)
\item $\left(\begin{smallmatrix}-5 & 2\\ 3 & -1\end{smallmatrix}\right)$
\item $\left(\begin{smallmatrix}1 & 3\\ 1 & 2\end{smallmatrix}\right)$
\item $\frac{1}{15}\left(\begin{smallmatrix}3 & 0\\ 0 & -5\end{smallmatrix}\right)$
\item does not exist ($\det=0$)
\item $\left(\begin{smallmatrix}5 & 3\\ 2 & 1\end{smallmatrix}\right)$
\item $\left(\begin{smallmatrix}1 & -2\\ -3 & 7\end{smallmatrix}\right)$
\item does not exist ($\det=0$)
\item $\left(\begin{smallmatrix}2 & 1\\ 1 & 1\end{smallmatrix}\right)$
\item $\left(\begin{smallmatrix}0 & 1\\ 1 & 0\end{smallmatrix}\right)$
\item does not exist ($\det=0$)
\item $\frac{1}{14}\left(\begin{smallmatrix}-3 & 1\\ 2 & 4\end{smallmatrix}\right)$
\item $\left(\begin{smallmatrix}1 & 0\\ 0 & 1\end{smallmatrix}\right)$
\item does not exist ($\det=0$)
\item $\frac{1}{17}\left(\begin{smallmatrix}3 & -2\\ 1 & 5\end{smallmatrix}\right)$
\item does not exist ($\det=0$)
\item $\left(\begin{smallmatrix}2 & -5\\ -3 & 8\end{smallmatrix}\right)$
\item $\left(\begin{smallmatrix}-7 & 4\\ 2 & -1\end{smallmatrix}\right)$
\item $\frac{1}{10}\left(\begin{smallmatrix}-5 & 0\\ 0 & -2\end{smallmatrix}\right)$
\item does not exist ($\det=0$)
\item $\frac{1}{13}\left(\begin{smallmatrix}1 & 1\\ -4 & 9\end{smallmatrix}\right)$
\item does not exist ($\det=0$)
\item $\left(\begin{smallmatrix}3 & -1\\ -5 & 2\end{smallmatrix}\right)$
\item $\left(\begin{smallmatrix}2 & -5\\ 1 & -3\end{smallmatrix}\right)$
\item $\frac{1}{6}\left(\begin{smallmatrix}0 & 2\\ -3 & 0\end{smallmatrix}\right)$
\item $\left(\begin{smallmatrix}11 & 4\\ -3 & -1\end{smallmatrix}\right)$
\item $\frac{1}{20}\left(\begin{smallmatrix}3 & -1\\ 2 & 6\end{smallmatrix}\right)$
\item does not exist ($\det=0$)
\item $\frac{1}{13}\left(\begin{smallmatrix}1 & 3\\ -2 & 7\end{smallmatrix}\right)$
\end{enumerate}
\end{multicols}
\par\noindent\emph{Exercise set 2.}\nopagebreak
\begin{enumerate}[label=\arabic*., leftmargin=*, itemsep=1pt]
\item $\left(\begin{smallmatrix}1 & -1 & 1\\ -38 & 41 & -34\\ 27 & -29 & 24\end{smallmatrix}\right)$
\item $\frac{1}{9}\left(\begin{smallmatrix}-1 & -2 & 5\\ 2 & 4 & -1\\ 6 & 3 & -3\end{smallmatrix}\right)$
\item $\frac{1}{73}\left(\begin{smallmatrix}26 & 5 & -2\\ 3 & 9 & 11\\ -5 & -15 & 6\end{smallmatrix}\right)$
\item $\frac{1}{41}\left(\begin{smallmatrix}2 & 11 & 1\\ 21 & 13 & -10\\ 8 & 3 & 4\end{smallmatrix}\right)$
\item $\frac{1}{10}\left(\begin{smallmatrix}4 & -8 & -6\\ 3 & 4 & 3\\ 1 & -2 & 1\end{smallmatrix}\right)$
\item $\frac{1}{43}\left(\begin{smallmatrix}5 & 1 & -10\\ -3 & 8 & 6\\ 10 & 2 & 23\end{smallmatrix}\right)$
\item $\frac{1}{16}\left(\begin{smallmatrix}8 & 4 & -4\\ 4 & 2 & 6\\ -6 & 1 & 3\end{smallmatrix}\right)$
\item $\frac{1}{8}\left(\begin{smallmatrix}2 & -9 & -3\\ 4 & -10 & -6\\ 2 & 3 & 1\end{smallmatrix}\right)$
\item $\frac{1}{16}\left(\begin{smallmatrix}4 & 2 & 14\\ 10 & 1 & 15\\ 8 & 4 & 12\end{smallmatrix}\right)$
\item $\frac{1}{54}\left(\begin{smallmatrix}10 & 4 & 2\\ -2 & 10 & 5\\ 4 & -20 & 17\end{smallmatrix}\right)$
\item $\frac{1}{20}\left(\begin{smallmatrix}15 & 5 & -10\\ -2 & 6 & 4\\ -3 & -1 & 6\end{smallmatrix}\right)$
\item $\frac{1}{48}\left(\begin{smallmatrix}-4 & 8 & 10\\ 2 & -4 & 7\\ 22 & 4 & 5\end{smallmatrix}\right)$
\item $\frac{1}{2}\left(\begin{smallmatrix}-20 & -11 & 12\\ 4 & 2 & -2\\ -10 & -5 & 6\end{smallmatrix}\right)$
\item $\frac{1}{49}\left(\begin{smallmatrix}10 & 6 & 3\\ -5 & -3 & -26\\ -3 & 8 & 4\end{smallmatrix}\right)$
\item $\frac{1}{53}\left(\begin{smallmatrix}6 & 9 & -2\\ -15 & 4 & 5\\ 8 & 12 & 15\end{smallmatrix}\right)$
\item $\frac{1}{61}\left(\begin{smallmatrix}17 & -5 & -1\\ 10 & 15 & 3\\ 6 & 9 & 14\end{smallmatrix}\right)$
\item $\frac{1}{39}\left(\begin{smallmatrix}5 & 11 & 2\\ -8 & -2 & -11\\ -20 & -5 & -8\end{smallmatrix}\right)$
\item $\frac{1}{17}\left(\begin{smallmatrix}4 & 1 & -3\\ 12 & 3 & 8\\ -9 & 2 & -6\end{smallmatrix}\right)$
\item $\frac{1}{26}\left(\begin{smallmatrix}-2 & -4 & -20\\ -7 & -1 & -5\\ -4 & -8 & -14\end{smallmatrix}\right)$
\item $\frac{1}{54}\left(\begin{smallmatrix}9 & 6 & 15\\ 12 & 8 & 2\\ 3 & 20 & 5\end{smallmatrix}\right)$
\item $\frac{1}{14}\left(\begin{smallmatrix}3 & 20 & 12\\ 1 & 2 & 4\\ 5 & 10 & 6\end{smallmatrix}\right)$
\item $\frac{1}{46}\left(\begin{smallmatrix}2 & 4 & 8\\ -20 & 6 & 12\\ -5 & -10 & 3\end{smallmatrix}\right)$
\item $\frac{1}{37}\left(\begin{smallmatrix}12 & 5 & 20\\ 9 & 13 & 15\\ -5 & 1 & 4\end{smallmatrix}\right)$
\item $\frac{1}{62}\left(\begin{smallmatrix}5 & 1 & -15\\ -2 & 12 & 6\\ 20 & 4 & 2\end{smallmatrix}\right)$
\item $\frac{1}{13}\left(\begin{smallmatrix}0 & 0 & 13\\ 5 & 1 & -20\\ 2 & 3 & -8\end{smallmatrix}\right)$
\item $\frac{1}{62}\left(\begin{smallmatrix}-20 & 2 & 8\\ 5 & 15 & -2\\ 1 & 3 & 12\end{smallmatrix}\right)$
\item $\frac{1}{27}\left(\begin{smallmatrix}12 & 20 & -3\\ 0 & 9 & 0\\ 3 & 5 & 6\end{smallmatrix}\right)$
\item $\frac{1}{2}\left(\begin{smallmatrix}-1 & -4 & -2\\ -5 & -20 & -12\\ -3 & -10 & -6\end{smallmatrix}\right)$
\item $\frac{1}{18}\left(\begin{smallmatrix}-2 & 1 & 7\\ 8 & 5 & -1\\ -6 & 3 & 3\end{smallmatrix}\right)$
\item $\frac{1}{38}\left(\begin{smallmatrix}-5 & 9 & 15\\ 6 & 12 & 20\\ -2 & -4 & 6\end{smallmatrix}\right)$
\end{enumerate}
\par\noindent\emph{Exercise set 3.}\nopagebreak
\begin{multicols}{2}\raggedright
\begin{enumerate}[label=\arabic*., leftmargin=*, itemsep=1pt]
\item $X=\left(\begin{smallmatrix}3\\ -8\end{smallmatrix}\right)$
\item $X=\left(\begin{smallmatrix}-16 & -11\\ 10 & 6\end{smallmatrix}\right)$
\item $X=\left(\begin{smallmatrix}-1\\ 1\end{smallmatrix}\right)$
\item $X=\frac{1}{17}\left(\begin{smallmatrix}7 & 2\\ 8 & -5\end{smallmatrix}\right)$
\item $X=\left(\begin{smallmatrix}7\\ 5\end{smallmatrix}\right)$
\item $X=\left(\begin{smallmatrix}-5 & 10\\ 17 & -33\end{smallmatrix}\right)$
\item $X=\frac{1}{14}\left(\begin{smallmatrix}5\\ 20\end{smallmatrix}\right)$
\item $X=\left(\begin{smallmatrix}0 & -10\\ 1 & 18\end{smallmatrix}\right)$
\item $X=\left(\begin{smallmatrix}9\\ -14\end{smallmatrix}\right)$
\item $X=\frac{1}{13}\left(\begin{smallmatrix}3 & 3\\ -12 & -38\end{smallmatrix}\right)$
\item $X=\left(\begin{smallmatrix}-22\\ -13\end{smallmatrix}\right)$
\item $X=\frac{1}{10}\left(\begin{smallmatrix}3 & -4\\ 2 & 14\end{smallmatrix}\right)$
\item $X=\left(\begin{smallmatrix}-29\\ 8\end{smallmatrix}\right)$
\item $X=\left(\begin{smallmatrix}1 & 2\\ 5 & -3\end{smallmatrix}\right)$
\item $X=\left(\begin{smallmatrix}-1\\ 1\end{smallmatrix}\right)$
\item $X=\left(\begin{smallmatrix}5 & -9\\ -6 & 12\end{smallmatrix}\right)$
\item $X=\left(\begin{smallmatrix}1\\ 2\end{smallmatrix}\right)$
\item $X=\left(\begin{smallmatrix}2 & -1\\ -2 & 3\end{smallmatrix}\right)$
\item $X=\left(\begin{smallmatrix}7\\ -3\end{smallmatrix}\right)$
\item $X=\left(\begin{smallmatrix}3 & 1\\ -2 & 3\end{smallmatrix}\right)$
\item $X=\frac{1}{13}\left(\begin{smallmatrix}1\\ -2\end{smallmatrix}\right)$
\item $\det A=0$; infinitely many solutions $X=\left(\begin{smallmatrix}t\\ -2-3t\end{smallmatrix}\right)$, $t\in\mathbb{R}$
\item $X=\left(\begin{smallmatrix}2 & -8\\ -1 & 5\end{smallmatrix}\right)$
\item $X=\left(\begin{smallmatrix}-13\\ -8\end{smallmatrix}\right)$
\item $X=\left(\begin{smallmatrix}-4 & 42\\ 1 & -11\end{smallmatrix}\right)$
\item $X=\frac{1}{2}\left(\begin{smallmatrix}3\\ 0\end{smallmatrix}\right)$
\item $X=\left(\begin{smallmatrix}19 & -5\\ -26 & 8\end{smallmatrix}\right)$
\item $\det A=0$; no solution (the system is inconsistent)
\item $X=\left(\begin{smallmatrix}-5 & 17\\ 4 & -13\end{smallmatrix}\right)$
\item $X=\frac{1}{3}\left(\begin{smallmatrix}1 & 1\\ -7 & 2\end{smallmatrix}\right)$
\end{enumerate}
\end{multicols}
\par\noindent\emph{Exercise set 4.}\nopagebreak
\begin{multicols}{2}\raggedright
\begin{enumerate}[label=\arabic*., leftmargin=*, itemsep=1pt]
\item $X=\left(\begin{smallmatrix}-23 & 9\\ -1 & 1\end{smallmatrix}\right)$
\item $X=\left(\begin{smallmatrix}-9 & 5\\ -15 & 9\end{smallmatrix}\right)$
\item $X=\left(\begin{smallmatrix}5 & 3\\ -11 & -7\end{smallmatrix}\right)$
\item $X=\left(\begin{smallmatrix}5 & -2\\ 2 & 0\end{smallmatrix}\right)$
\item $X=\left(\begin{smallmatrix}2 & -3\\ -2 & 4\end{smallmatrix}\right)$
\item $X=\left(\begin{smallmatrix}2 & -1\\ 2 & 1\end{smallmatrix}\right)$
\item $X=\frac{1}{17}\left(\begin{smallmatrix}-1 & 12\\ 4 & -31\end{smallmatrix}\right)$
\item $X=\left(\begin{smallmatrix}1 & 2\\ 2 & 2\end{smallmatrix}\right)$
\item $X=\left(\begin{smallmatrix}5 & -13\\ -1 & 4\end{smallmatrix}\right)$
\item $X=\left(\begin{smallmatrix}-3 & 2\\ 1 & 5\end{smallmatrix}\right)$
\item $X=\frac{1}{14}\left(\begin{smallmatrix}18 & 6\\ -13 & 19\end{smallmatrix}\right)$
\item $X=\frac{1}{13}\left(\begin{smallmatrix}10 & -16\\ -15 & 37\end{smallmatrix}\right)$
\item $X=\left(\begin{smallmatrix}3 & -9\\ 0 & -1\end{smallmatrix}\right)$
\item $X=\frac{1}{20}\left(\begin{smallmatrix}-7 & 9\\ 6 & -2\end{smallmatrix}\right)$
\item $X=\left(\begin{smallmatrix}-23 & 13\\ 4 & -2\end{smallmatrix}\right)$
\item $X=\frac{1}{14}\left(\begin{smallmatrix}-5 & -3\\ -11 & 13\end{smallmatrix}\right)$
\item $X=\left(\begin{smallmatrix}8 & 4\\ 5 & 4\end{smallmatrix}\right)$
\item $X=\left(\begin{smallmatrix}2 & 0\\ 0 & -3\end{smallmatrix}\right)$
\item $X=\left(\begin{smallmatrix}11 & -8\\ 9 & -6\end{smallmatrix}\right)$
\item $X=\left(\begin{smallmatrix}-2 & 2\\ -3 & -1\end{smallmatrix}\right)$
\item $X=\frac{1}{13}\left(\begin{smallmatrix}2 & -7\\ -8 & 41\end{smallmatrix}\right)$
\item $X=\left(\begin{smallmatrix}-10 & 4\\ 12 & -5\end{smallmatrix}\right)$
\item $X=\left(\begin{smallmatrix}-5 & 9\\ 10 & -16\end{smallmatrix}\right)$
\item $X=\left(\begin{smallmatrix}55 & 20\\ -20 & -7\end{smallmatrix}\right)$
\item $X=\left(\begin{smallmatrix}-1 & 4\\ -6 & 22\end{smallmatrix}\right)$
\item $X=\left(\begin{smallmatrix}18 & -7\\ 4 & -1\end{smallmatrix}\right)$
\item $X=\left(\begin{smallmatrix}-1 & 2\\ -1 & 1\end{smallmatrix}\right)$
\item $X=\frac{1}{2}\left(\begin{smallmatrix}-4 & 0\\ 1 & 9\end{smallmatrix}\right)$
\item $X=\left(\begin{smallmatrix}-4 & -7\\ 2 & 4\end{smallmatrix}\right)$
\item $X=\left(\begin{smallmatrix}-2 & -5\\ 5 & 10\end{smallmatrix}\right)$
\end{enumerate}
\end{multicols}

% !Mode:: "TeX:UTF-8"
\javobmashgulot{6}{Computing the inverse matrix by elementary row operations}
\par\noindent\emph{Exercise set 1.}\nopagebreak
\begin{multicols}{2}\raggedright
\begin{enumerate}[label=\arabic*., leftmargin=*, itemsep=1pt]
\item $\left(\begin{smallmatrix}3 & -1\\ -5 & 2\end{smallmatrix}\right)$
\item $\left(\begin{smallmatrix}-5 & 3\\ 2 & -1\end{smallmatrix}\right)$
\item $\left(\begin{smallmatrix}2 & 1\\ 5 & 3\end{smallmatrix}\right)$
\item $\frac{1}{2}\left(\begin{smallmatrix}4 & -2\\ -7 & 4\end{smallmatrix}\right)$
\item $\left(\begin{smallmatrix}5 & 2\\ 3 & 1\end{smallmatrix}\right)$
\item $\frac{1}{20}\left(\begin{smallmatrix}3 & -1\\ 2 & 6\end{smallmatrix}\right)$
\item $\left(\begin{smallmatrix}-7 & -4\\ -2 & -1\end{smallmatrix}\right)$
\item $\left(\begin{smallmatrix}-1 & 3\\ 2 & -5\end{smallmatrix}\right)$
\item $\left(\begin{smallmatrix}-3 & -2\\ -4 & -3\end{smallmatrix}\right)$
\item $\left(\begin{smallmatrix}1 & 1\\ 1 & 2\end{smallmatrix}\right)$
\item $\left(\begin{smallmatrix}1 & -2\\ 3 & -7\end{smallmatrix}\right)$
\item $\left(\begin{smallmatrix}2 & -1\\ -1 & 1\end{smallmatrix}\right)$
\item $\left(\begin{smallmatrix}2 & -5\\ -3 & 8\end{smallmatrix}\right)$
\item $\left(\begin{smallmatrix}-3 & -1\\ 5 & 2\end{smallmatrix}\right)$
\item $\left(\begin{smallmatrix}1 & -3\\ 1 & -4\end{smallmatrix}\right)$
\item $\left(\begin{smallmatrix}3 & -4\\ -2 & 3\end{smallmatrix}\right)$
\item $\frac{1}{13}\left(\begin{smallmatrix}1 & 1\\ -4 & 9\end{smallmatrix}\right)$
\item $\left(\begin{smallmatrix}5 & 3\\ -2 & -1\end{smallmatrix}\right)$
\item $\left(\begin{smallmatrix}-1 & -2\\ -3 & -5\end{smallmatrix}\right)$
\item $\left(\begin{smallmatrix}3 & -1\\ 1 & 0\end{smallmatrix}\right)$
\item $\left(\begin{smallmatrix}4 & -7\\ -1 & 2\end{smallmatrix}\right)$
\item $\left(\begin{smallmatrix}-1 & -1\\ -3 & -4\end{smallmatrix}\right)$
\item $\left(\begin{smallmatrix}-5 & 2\\ -3 & 1\end{smallmatrix}\right)$
\item $\left(\begin{smallmatrix}1 & 5\\ 1 & 6\end{smallmatrix}\right)$
\item $\left(\begin{smallmatrix}3 & 2\\ -4 & -3\end{smallmatrix}\right)$
\item $\frac{1}{6}\left(\begin{smallmatrix}1 & 2\\ -3 & 0\end{smallmatrix}\right)$
\item $\left(\begin{smallmatrix}1 & -3\\ -2 & 7\end{smallmatrix}\right)$
\item $\left(\begin{smallmatrix}2 & 1\\ -7 & -3\end{smallmatrix}\right)$
\item $\frac{1}{5}\left(\begin{smallmatrix}5 & 0\\ 0 & -1\end{smallmatrix}\right)$
\item $\left(\begin{smallmatrix}-1 & 3\\ -1 & 2\end{smallmatrix}\right)$
\end{enumerate}
\end{multicols}
\par\noindent\emph{Exercise set 2.}\nopagebreak
\begin{enumerate}[label=\arabic*., leftmargin=*, itemsep=1pt]
\item $\left(\begin{smallmatrix}3 & 3 & -1\\ -2 & -2 & 1\\ -4 & -5 & 2\end{smallmatrix}\right)$
\item $\frac{1}{8}\left(\begin{smallmatrix}2 & 1 & 1\\ -4 & 2 & 2\\ 2 & -7 & 1\end{smallmatrix}\right)$
\item $\frac{1}{5}\left(\begin{smallmatrix}-2 & 4 & -7\\ 3 & -1 & 3\\ -1 & 2 & -1\end{smallmatrix}\right)$
\item $\frac{1}{22}\left(\begin{smallmatrix}6 & -4 & -2\\ 3 & 9 & -1\\ 4 & 12 & 6\end{smallmatrix}\right)$
\item $\frac{1}{41}\left(\begin{smallmatrix}2 & 11 & 1\\ 21 & 13 & -10\\ 8 & 3 & 4\end{smallmatrix}\right)$
\item $\frac{1}{13}\left(\begin{smallmatrix}1 & 2 & 6\\ 2 & 4 & -1\\ -6 & 1 & 3\end{smallmatrix}\right)$
\item $\frac{1}{18}\left(\begin{smallmatrix}8 & 7 & -4\\ 6 & 3 & 6\\ -4 & 1 & 2\end{smallmatrix}\right)$
\item $\frac{1}{43}\left(\begin{smallmatrix}5 & 1 & -10\\ -3 & 8 & 6\\ 10 & 2 & 23\end{smallmatrix}\right)$
\item $\frac{1}{8}\left(\begin{smallmatrix}2 & -9 & -3\\ 4 & -10 & -6\\ 2 & 3 & 1\end{smallmatrix}\right)$
\item $\frac{1}{16}\left(\begin{smallmatrix}4 & 2 & 14\\ 10 & 1 & 15\\ 8 & 4 & 12\end{smallmatrix}\right)$
\item $\frac{1}{54}\left(\begin{smallmatrix}10 & 4 & 2\\ -2 & 10 & 5\\ 4 & -20 & 17\end{smallmatrix}\right)$
\item $\frac{1}{20}\left(\begin{smallmatrix}15 & 5 & -10\\ -2 & 6 & 4\\ -3 & -1 & 6\end{smallmatrix}\right)$
\item $\frac{1}{48}\left(\begin{smallmatrix}-4 & 8 & 10\\ 2 & -4 & 7\\ 22 & 4 & 5\end{smallmatrix}\right)$
\item $\frac{1}{2}\left(\begin{smallmatrix}-20 & -11 & 12\\ 4 & 2 & -2\\ -10 & -5 & 6\end{smallmatrix}\right)$
\item $\frac{1}{49}\left(\begin{smallmatrix}10 & 6 & 3\\ -5 & -3 & -26\\ -3 & 8 & 4\end{smallmatrix}\right)$
\item $\frac{1}{53}\left(\begin{smallmatrix}6 & 9 & -2\\ -15 & 4 & 5\\ 8 & 12 & 15\end{smallmatrix}\right)$
\item $\frac{1}{61}\left(\begin{smallmatrix}17 & -5 & -1\\ 10 & 15 & 3\\ 6 & 9 & 14\end{smallmatrix}\right)$
\item $\frac{1}{39}\left(\begin{smallmatrix}5 & 11 & 2\\ -8 & -2 & -11\\ -20 & -5 & -8\end{smallmatrix}\right)$
\item $\frac{1}{17}\left(\begin{smallmatrix}4 & 1 & -3\\ 12 & 3 & 8\\ -9 & 2 & -6\end{smallmatrix}\right)$
\item $\frac{1}{26}\left(\begin{smallmatrix}-2 & -4 & -20\\ -7 & -1 & -5\\ -4 & -8 & -14\end{smallmatrix}\right)$
\item $\frac{1}{54}\left(\begin{smallmatrix}9 & 6 & 15\\ 12 & 8 & 2\\ 3 & 20 & 5\end{smallmatrix}\right)$
\item $\frac{1}{14}\left(\begin{smallmatrix}3 & 20 & 12\\ 1 & 2 & 4\\ 5 & 10 & 6\end{smallmatrix}\right)$
\item $\frac{1}{46}\left(\begin{smallmatrix}2 & 4 & 8\\ -20 & 6 & 12\\ -5 & -10 & 3\end{smallmatrix}\right)$
\item $\frac{1}{37}\left(\begin{smallmatrix}12 & 5 & 20\\ 9 & 13 & 15\\ -5 & 1 & 4\end{smallmatrix}\right)$
\item $\frac{1}{62}\left(\begin{smallmatrix}5 & 1 & -15\\ -2 & 12 & 6\\ 20 & 4 & 2\end{smallmatrix}\right)$
\item $\frac{1}{13}\left(\begin{smallmatrix}0 & 0 & 13\\ 5 & 1 & -20\\ 2 & 3 & -8\end{smallmatrix}\right)$
\item $\frac{1}{62}\left(\begin{smallmatrix}-20 & 2 & 8\\ 5 & 15 & -2\\ 1 & 3 & 12\end{smallmatrix}\right)$
\item $\frac{1}{27}\left(\begin{smallmatrix}12 & 20 & -3\\ 0 & 9 & 0\\ 3 & 5 & 6\end{smallmatrix}\right)$
\item $\frac{1}{18}\left(\begin{smallmatrix}-2 & 1 & 7\\ 8 & 5 & -1\\ -6 & 3 & 3\end{smallmatrix}\right)$
\item $\frac{1}{38}\left(\begin{smallmatrix}-5 & 9 & 15\\ 6 & 12 & 20\\ -2 & -4 & 6\end{smallmatrix}\right)$
\end{enumerate}
\par\noindent\emph{Exercise set 3.}\nopagebreak
\begin{multicols}{2}\raggedright
\begin{enumerate}[label=\arabic*., leftmargin=*, itemsep=1pt]
\item $X=\frac{1}{5}\left(\begin{smallmatrix}13\\ -7\end{smallmatrix}\right)$
\item $X=\left(\begin{smallmatrix}1 & -7\\ 1 & 10\end{smallmatrix}\right)$
\item $X=\left(\begin{smallmatrix}10\\ 7\end{smallmatrix}\right)$
\item $X=\frac{1}{5}\left(\begin{smallmatrix}-24\\ 16\\ -7\end{smallmatrix}\right)$
\item $X=\frac{1}{17}\left(\begin{smallmatrix}1 & 10\\ -5 & 18\end{smallmatrix}\right)$
\item $X=\frac{1}{13}\left(\begin{smallmatrix}-17\\ 5\\ -15\end{smallmatrix}\right)$
\item $X=\left(\begin{smallmatrix}1\\ 1\end{smallmatrix}\right)$
\item $X=\frac{1}{22}\left(\begin{smallmatrix}16\\ 19\\ -4\end{smallmatrix}\right)$
\item $X=\left(\begin{smallmatrix}-1 & 3\\ -2 & 4\end{smallmatrix}\right)$
\item $X=\frac{1}{8}\left(\begin{smallmatrix}3\\ 6\\ -37\end{smallmatrix}\right)$
\item $X=\left(\begin{smallmatrix}23\\ 7\end{smallmatrix}\right)$
\item $X=\frac{1}{41}\left(\begin{smallmatrix}-29\\ -38\\ 7\end{smallmatrix}\right)$
\item $X=\left(\begin{smallmatrix}3 & -9\\ 3 & -11\end{smallmatrix}\right)$
\item $X=\frac{1}{43}\left(\begin{smallmatrix}-1\\ 35\\ -2\end{smallmatrix}\right)$
\item $X=\left(\begin{smallmatrix}25\\ -7\end{smallmatrix}\right)$
\item $X=\frac{1}{8}\left(\begin{smallmatrix}25\\ 34\\ 5\end{smallmatrix}\right)$
\item $X=\left(\begin{smallmatrix}-10 & 9\\ 7 & -6\end{smallmatrix}\right)$
\item $X=\frac{1}{4}\left(\begin{smallmatrix}-2\\ -3\\ 0\end{smallmatrix}\right)$
\item $X=\frac{1}{13}\left(\begin{smallmatrix}7\\ -2\end{smallmatrix}\right)$
\item $X=\frac{1}{27}\left(\begin{smallmatrix}17\\ 2\\ 50\end{smallmatrix}\right)$
\item $X=\left(\begin{smallmatrix}6 & 7\\ -2 & -2\end{smallmatrix}\right)$
\item $X=\frac{1}{9}\left(\begin{smallmatrix}11\\ 24\\ -10\end{smallmatrix}\right)$
\item $X=\left(\begin{smallmatrix}-13\\ 30\end{smallmatrix}\right)$
\item $X=\frac{1}{4}\left(\begin{smallmatrix}-1\\ -6\\ 1\end{smallmatrix}\right)$
\item $X=\left(\begin{smallmatrix}-1 & 1\\ 0 & -3\end{smallmatrix}\right)$
\item $X=\frac{1}{48}\left(\begin{smallmatrix}86\\ 41\\ -53\end{smallmatrix}\right)$
\item $X=\left(\begin{smallmatrix}-7 & -3\\ -18 & -7\end{smallmatrix}\right)$
\item $X=\frac{1}{2}\left(\begin{smallmatrix}-47\\ 8\\ -23\end{smallmatrix}\right)$
\item $X=\left(\begin{smallmatrix}-7\\ -6\end{smallmatrix}\right)$
\item $X=\frac{1}{46}\left(\begin{smallmatrix}38\\ 80\\ -49\end{smallmatrix}\right)$
\end{enumerate}
\end{multicols}
\par\noindent\emph{Exercise set 4.}\nopagebreak
\begin{enumerate}[label=\arabic*., leftmargin=*, itemsep=1pt]
\item $K^{-1}\equiv\left(\begin{smallmatrix}5&12\\15&25\end{smallmatrix}\right)\pmod{26}$; $p=\left(\begin{smallmatrix}25\\9\end{smallmatrix}\right)$
\item $\det K=441\equiv25\pmod{26}$, $\gcd(25,26)=1$, so $K^{-1}$ exists; $K^{-1}\equiv\left(\begin{smallmatrix}8&5&10\\21&8&21\\21&12&8\end{smallmatrix}\right)\pmod{26}$
\item $X=A^{-1}(C-B)=\frac15\left(\begin{smallmatrix}10&-6\\-5&9\end{smallmatrix}\right)$
\item $\det A=-5$, $A^{-1}=\frac15\left(\begin{smallmatrix}-2&4&-7\\3&-1&3\\-1&2&-1\end{smallmatrix}\right)$
\item $A^{-1}=\frac15\left(\begin{smallmatrix}1&1\\-3&2\end{smallmatrix}\right)$ (the same by both methods)
\item $G'=\left(\begin{smallmatrix}1&0&0&0&1&1\\0&1&0&1&1&1\\0&0&1&1&0&1\end{smallmatrix}\right)$
\item every row sum is $0$, so the sum of the columns is the zero vector; hence $\det L=0$ and $L^{-1}$ does not exist
\item $\det A=1+a^3\ne0$, i.e. $a\ne-1$; for $a=1$, $A^{-1}=\frac12\left(\begin{smallmatrix}1&-1&1\\1&1&-1\\-1&1&1\end{smallmatrix}\right)$
\item $A^{-1}=\frac15\left(\begin{smallmatrix}3&1&1\\-6&3&-2\\-2&1&1\end{smallmatrix}\right)$, $B^{-1}=\left(\begin{smallmatrix}1&-1&-1\\0&1&1\\0&0&1\end{smallmatrix}\right)$, $(AB)^{-1}=\frac15\left(\begin{smallmatrix}11&-3&2\\-8&4&-1\\-2&1&1\end{smallmatrix}\right)$
\item $\det M=0$ (mod 2), so $M$ is not invertible
\end{enumerate}


\section*{Chapter 2. Systems of Linear Algebraic Equations, Vectors and Lines in the Plane}
% !Mode:: "TeX:UTF-8"
\javobmashgulot{7}{Cramer’s rule for systems of linear algebraic equations}
\par\noindent\emph{Exercise set 1.}\nopagebreak
\begin{multicols}{2}\raggedright
\begin{enumerate}[label=\arabic*., leftmargin=*, itemsep=1pt]
\item $\Delta=-4,\ \Delta_1=24,\ \Delta_2=20;\,x_1=-6,\ x_2=-5$
\item $\Delta=1,\ \Delta_1=-6,\ \Delta_2=2;\,x_1=-6,\ x_2=2$
\item $\Delta=6,\ \Delta_1=-30,\ \Delta_2=-18;\,x_1=-5,\ x_2=-3$
\item $\Delta=-4,\ \Delta_1=20,\ \Delta_2=-12;\,x_1=-5,\ x_2=3$
\item $\Delta=-11,\ \Delta_1=-55,\ \Delta_2=55;\,x_1=5,\ x_2=-5$
\item $\Delta=9,\ \Delta_1=9,\ \Delta_2=36;\,x_1=1,\ x_2=4$
\item $\Delta=-2,\ \Delta_1=-2,\ \Delta_2=-6;\,x_1=1,\ x_2=3$
\item $\Delta=-5,\ \Delta_1=-30,\ \Delta_2=20;\,x_1=6,\ x_2=-4$
\item $\Delta=-7,\ \Delta_1=-21,\ \Delta_2=14;\,x_1=3,\ x_2=-2$
\item $\Delta=11,\ \Delta_1=11,\ \Delta_2=-22;\,x_1=1,\ x_2=-2$
\item $\Delta=3,\ \Delta_1=15,\ \Delta_2=-3;\,x_1=5,\ x_2=-1$
\item $\Delta=-8,\ \Delta_1=16,\ \Delta_2=-40;\,x_1=-2,\ x_2=5$
\item $\Delta=4,\ \Delta_1=4,\ \Delta_2=-4;\,x_1=1,\ x_2=-1$
\item $\Delta=11,\ \Delta_1=-66,\ \Delta_2=-33;\,x_1=-6,\ x_2=-3$
\item $\Delta=-3,\ \Delta_1=15,\ \Delta_2=12;\,x_1=-5,\ x_2=-4$
\item $\Delta=-2,\ \Delta_1=8,\ \Delta_2=0;\,x_1=-4,\ x_2=0$
\item $\Delta=10,\ \Delta_1=-10,\ \Delta_2=40;\,x_1=-1,\ x_2=4$
\item $\Delta=-12,\ \Delta_1=48,\ \Delta_2=48;\,x_1=-4,\ x_2=-4$
\item $\Delta=10,\ \Delta_1=-20,\ \Delta_2=-60;\,x_1=-2,\ x_2=-6$
\item $\Delta=4,\ \Delta_1=12,\ \Delta_2=12;\,x_1=3,\ x_2=3$
\item $\Delta=-1,\ \Delta_1=6,\ \Delta_2=-1;\,x_1=-6,\ x_2=1$
\item $\Delta=-12,\ \Delta_1=0,\ \Delta_2=0;\,x_1=0,\ x_2=0$
\item $\Delta=-12,\ \Delta_1=60,\ \Delta_2=36;\,x_1=-5,\ x_2=-3$
\item $\Delta=-7,\ \Delta_1=7,\ \Delta_2=-21;\,x_1=-1,\ x_2=3$
\item $\Delta=2,\ \Delta_1=-10,\ \Delta_2=-10;\,x_1=-5,\ x_2=-5$
\item $\Delta=-10,\ \Delta_1=-50,\ \Delta_2=10;\,x_1=5,\ x_2=-1$
\item $\Delta=6,\ \Delta_1=12,\ \Delta_2=-36;\,x_1=2,\ x_2=-6$
\item $\Delta=-10,\ \Delta_1=30,\ \Delta_2=-20;\,x_1=-3,\ x_2=2$
\item $\Delta=3,\ \Delta_1=-3,\ \Delta_2=15;\,x_1=-1,\ x_2=5$
\item $\Delta=-1,\ \Delta_1=-3,\ \Delta_2=1;\,x_1=3,\ x_2=-1$
\end{enumerate}
\end{multicols}
\par\noindent\emph{Exercise set 2.}\nopagebreak
\begin{multicols}{2}\raggedright
\begin{enumerate}[label=\arabic*., leftmargin=*, itemsep=1pt]
\item $\Delta=20;\,x_1=1,\ x_2=-1,\ x_3=3$
\item $\Delta=-36;\,x_1=2,\ x_2=1,\ x_3=-3$
\item $\Delta=2;\,x_1=3,\ x_2=-2,\ x_3=3$
\item $\Delta=-12;\,x_1=4,\ x_2=2,\ x_3=4$
\item $\Delta=21;\,x_1=-4,\ x_2=-3,\ x_3=4$
\item $\Delta=18;\,x_1=-3,\ x_2=-1,\ x_3=2$
\item $\Delta=-4;\,x_1=3,\ x_2=-1,\ x_3=-3$
\item $\Delta=3;\,x_1=2,\ x_2=-1,\ x_3=1$
\item $\Delta=-19;\,x_1=2,\ x_2=1,\ x_3=4$
\item $\Delta=23;\,x_1=0,\ x_2=-4,\ x_3=-3$
\item $\Delta=6;\,x_1=4,\ x_2=2,\ x_3=0$
\item $\Delta=28;\,x_1=-1,\ x_2=0,\ x_3=0$
\item $\Delta=16;\,x_1=-1,\ x_2=3,\ x_3=-3$
\item $\Delta=10;\,x_1=-1,\ x_2=-2,\ x_3=2$
\item $\Delta=-22;\,x_1=-2,\ x_2=0,\ x_3=3$
\item $\Delta=-16;\,x_1=-1,\ x_2=1,\ x_3=-2$
\item $\Delta=-9;\,x_1=4,\ x_2=-4,\ x_3=-3$
\item $\Delta=-19;\,x_1=-2,\ x_2=-4,\ x_3=4$
\item $\Delta=40;\,x_1=-4,\ x_2=-2,\ x_3=1$
\item $\Delta=28;\,x_1=0,\ x_2=-4,\ x_3=3$
\item $\Delta=-6;\,x_1=0,\ x_2=-4,\ x_3=-1$
\item $\Delta=-34;\,x_1=3,\ x_2=0,\ x_3=-3$
\item $\Delta=21;\,x_1=4,\ x_2=3,\ x_3=0$
\item $\Delta=-13;\,x_1=-2,\ x_2=4,\ x_3=0$
\item $\Delta=22;\,x_1=3,\ x_2=3,\ x_3=2$
\item $\Delta=12;\,x_1=1,\ x_2=1,\ x_3=2$
\item $\Delta=9;\,x_1=-3,\ x_2=1,\ x_3=2$
\item $\Delta=-14;\,x_1=4,\ x_2=-1,\ x_3=-3$
\item $\Delta=-8;\,x_1=-2,\ x_2=3,\ x_3=2$
\item $\Delta=-11;\,x_1=-1,\ x_2=-3,\ x_3=-2$
\end{enumerate}
\end{multicols}
\par\noindent\emph{Exercise set 3.}\nopagebreak
\begin{enumerate}[label=\arabic*., leftmargin=*, itemsep=1pt]
\item $\Delta=11$, $p=(7;\,4)$
\item $\Delta=5$, $p=(11;\,5)$
\item $\Delta=7$, $p=(8;\,19)$
\item $\Delta=4-k^2\ne0$, i.e. $k\ne\pm2$; for $k=1$, $x_1=2,\ x_2=1$
\item $\Delta=a^2+a-6=0$, i.e. for $a=2$ or $a=-3$
\item $\Delta=30,\ \Delta_1=30,\ \Delta_2=-30,\ \Delta_3=60$; $(1;-1;2)$
\item $\Delta=7$; $x_1=\frac{25}{7},\ x_2=\frac{23}{7},\ x_3=\frac{36}{7}$
\item $\Delta=-23$; $x_1=\frac{41}{23},\ x_2=\frac{19}{23}$
\item $\Delta=0$; the first system has infinitely many solutions ($x_1=3+2x_2$), the second is inconsistent
\item $a=3,\ b=2,\ c=5,\ d=7$, i.e. $K=\left(\begin{smallmatrix}3&2\\5&7\end{smallmatrix}\right)$
\end{enumerate}

% !Mode:: "TeX:UTF-8"
\javobmashgulot{8}{The matrix method for systems of linear equations}
\par\noindent\emph{Exercise set 1.}\nopagebreak
\begin{multicols}{3}\raggedright
\begin{enumerate}[label=\arabic*., leftmargin=*, itemsep=1pt]
\item $x_1=3,\ x_2=-4$
\item $x_1=13,\ x_2=9$
\item $x_1=2,\ x_2=3$
\item $x_1=\frac{13}{7},\ x_2=\frac{10}{7}$
\item $x_1=\frac{22}{7},\ x_2=\frac{17}{7}$
\item $x_1=1,\ x_2=2$
\item $x_1=5,\ x_2=3$
\item $x_1=\frac{34}{11},\ x_2=\frac{13}{11}$
\item $x_1=1,\ x_2=5$
\item $x_1=\frac{7}{2},\ x_2=4$
\item $x_1=3,\ x_2=1$
\item $x_1=1,\ x_2=1$
\item $x_1=\frac{11}{3},\ x_2=\frac{10}{3}$
\item $x_1=1,\ x_2=2$
\item $x_1=\frac{13}{5},\ x_2=\frac{11}{5}$
\item $x_1=1,\ x_2=1$
\item $x_1=1,\ x_2=2$
\item $x_1=\frac{9}{7},\ x_2=\frac{17}{7}$
\item $x_1=3,\ x_2=2$
\item $x_1=1,\ x_2=1$
\item $x_1=2,\ x_2=4$
\item $x_1=\frac{11}{5},\ x_2=\frac{23}{5}$
\item $x_1=\frac{2}{17},\ x_2=\frac{31}{17}$
\item $x_1=\frac{13}{11},\ x_2=\frac{25}{11}$
\item $x_1=\frac{13}{7},\ x_2=\frac{19}{7}$
\item $x_1=-\frac{1}{7},\ x_2=\frac{17}{7}$
\item $x_1=\frac{26}{11},\ x_2=\frac{23}{11}$
\item $x_1=-\frac{5}{7},\ x_2=\frac{17}{7}$
\item $x_1=\frac{13}{7},\ x_2=\frac{23}{7}$
\item $x_1=1,\ x_2=2$
\end{enumerate}
\end{multicols}
\par\noindent\emph{Exercise set 2.}\nopagebreak
\begin{multicols}{2}\raggedright
\begin{enumerate}[label=\arabic*., leftmargin=*, itemsep=1pt]
\item $x_1=-2,\ x_2=3,\ x_3=5$
\item $x_1=\frac{11}{14},\ x_2=\frac{13}{14},\ x_3=\frac{39}{14}$
\item $x_1=-3,\ x_2=5,\ x_3=4$
\item $x_1=4,\ x_2=3,\ x_3=2$
\item $x_1=\frac{1}{10},\ x_2=\frac{3}{2},\ x_3=\frac{31}{10}$
\item $x_1=1,\ x_2=2,\ x_3=3$
\item $x_1=2,\ x_2=1,\ x_3=3$
\item $x_1=\frac{20}{11},\ x_2=\frac{78}{55},\ x_3=\frac{59}{55}$
\item $x_1=\frac{7}{6},\ x_2=\frac{1}{3},\ x_3=\frac{13}{6}$
\item $x_1=\frac{29}{22},\ x_2=\frac{3}{2},\ x_3=\frac{67}{22}$
\item $x_1=\frac{6}{11},\ x_2=\frac{8}{11},\ x_3=\frac{29}{11}$
\item $x_1=\frac{5}{3},\ x_2=\frac{5}{3},\ x_3=2$
\item $x_1=\frac{59}{78},\ x_2=\frac{103}{78},\ x_3=\frac{145}{78}$
\item $x_1=\frac{55}{32},\ x_2=\frac{79}{32},\ x_3=\frac{53}{32}$
\item $x_1=\frac{8}{7},\ x_2=1,\ x_3=\frac{18}{7}$
\item $x_1=\frac{1}{2},\ x_2=\frac{13}{10},\ x_3=\frac{29}{10}$
\item $x_1=\frac{15}{8},\ x_2=\frac{13}{8},\ x_3=\frac{11}{8}$
\item $x_1=\frac{36}{25},\ x_2=\frac{13}{5},\ x_3=\frac{34}{25}$
\item $x_1=7,\ x_2=-2,\ x_3=4$
\item $x_1=\frac{16}{11},\ x_2=\frac{23}{11},\ x_3=\frac{21}{11}$
\item $x_1=\frac{27}{11},\ x_2=\frac{20}{11},\ x_3=\frac{52}{11}$
\item $x_1=\frac{9}{5},\ x_2=\frac{46}{25},\ x_3=\frac{53}{25}$
\item $x_1=\frac{16}{13},\ x_2=\frac{11}{13},\ x_3=\frac{31}{13}$
\item $x_1=\frac{29}{16},\ x_2=\frac{25}{16},\ x_3=\frac{21}{16}$
\item $x_1=\frac{13}{5},\ x_2=\frac{6}{5},\ x_3=1$
\item $x_1=\frac{1}{2},\ x_2=\frac{3}{2},\ x_3=\frac{5}{2}$
\item $x_1=1,\ x_2=2,\ x_3=3$
\item $x_1=\frac{4}{3},\ x_2=\frac{5}{3},\ x_3=\frac{8}{3}$
\item $x_1=\frac{41}{25},\ x_2=\frac{53}{25},\ x_3=\frac{54}{25}$
\item $x_1=\frac{41}{15},\ x_2=\frac{5}{3},\ x_3=\frac{2}{5}$
\end{enumerate}
\end{multicols}
\par\noindent\emph{Exercise set 3.}\nopagebreak
\begin{enumerate}[label=\arabic*., leftmargin=*, itemsep=1pt]
\item (a) $x=(1;\,1)$; (b) $x=(1;\,3)$; (c) $x=(-4;\,7)$
\item (a) $x=(\frac{9}{5};\,-\frac{11}{5})$; (b) $x=(1;\,1)$; (c) $x=(\frac{1}{5};\,\frac{16}{5})$
\item (a) $x=(\frac{13}{17};\,\frac{27}{17})$; (b) $x=(\frac{37}{17};\,\frac{1}{17})$; (c) $x=(\frac{2}{17};\,-\frac{5}{17})$
\item (a) $x=(\frac{11}{7};\,\frac{19}{7})$; (b) $x=(\frac{2}{7};\,\frac{13}{7})$; (c) $x=(\frac{8}{7};\,-\frac{32}{7})$
\item (a) $x=(1;\,2)$; (b) $x=(-\frac{16}{7};\,\frac{3}{7})$; (c) $x=(2;\,0)$
\item (a) $x=(-2;\,3;\,5)$; (b) $x=(-4;\,3;\,9)$
\item (a) $x=(-\frac{1}{4};\,\frac{5}{8};\,-\frac{9}{8})$; (b) $x=(-3;\,3;\,1)$
\item (a) $x=(\frac{7}{11};\,-\frac{2}{11};\,\frac{23}{11})$; (b) $x=(-\frac{12}{11};\,\frac{21}{22};\,\frac{14}{11})$
\item (a) $x=(\frac{13}{7};\,\frac{10}{7})$; (b) $x=(\frac{1}{7};\,\frac{25}{7})$; (c) $x=(-\frac{1}{7};\,-\frac{11}{7})$
\item (a) $x=(\frac{7}{13};\,\frac{1}{13};\,-\frac{29}{13})$; (b) $x=(-\frac{10}{13};\,\frac{19}{13};\,-\frac{5}{13})$
\item (a) $x=(1;\,2)$; (b) $x=(-5;\,-15)$; (c) $x=(1;\,-1)$
\item (a) $x=(-3;\,-1)$; (b) $x=(35;\,21)$; (c) $x=(2;\,0)$
\item (a) $x=(\frac{5}{8};\,\frac{1}{4};\,\frac{13}{8})$; (b) $x=(\frac{1}{2};\,-2;\,\frac{1}{2})$
\item (a) $x=(2;\,-3)$; (b) $x=(1;\,1)$; (c) $x=(1;\,-3)$
\item (a) $x=(\frac{32}{41};\,\frac{49}{41};\,\frac{5}{41})$; (b) $x=(-\frac{8}{41};\,\frac{39}{41};\,\frac{50}{41})$
\item (a) $x=(1;\,2)$; (b) $x=(1;\,-1)$; (c) $x=(5;\,5)$
\item (a) $x=(\frac{4}{43};\,-\frac{11}{43};\,\frac{137}{43})$; (b) $x=(\frac{48}{43};\,\frac{40}{43};\,-\frac{76}{43})$
\item (a) $x=(3;\,-2)$; (b) $x=(1;\,1)$; (c) $x=(-18;\,13)$
\item (a) $x=(\frac{59}{53};\,-\frac{15}{53};\,\frac{61}{53})$; (b) $x=(-\frac{32}{53};\,\frac{80}{53};\,\frac{81}{53})$
\item (a) $x=(1;\,1)$; (b) $x=(\frac{11}{5};\,-\frac{1}{5})$; (c) $x=(-\frac{3}{10};\,\frac{9}{5})$
\item (a) $x=(11;\,-1;\,4)$; (b) $x=(41;\,-7;\,21)$
\item (a) $x=(-1;\,-3)$; (b) $x=(-2;\,-3)$; (c) $x=(3;\,2)$
\item (a) $x=(-\frac{11}{61};\,\frac{94}{61};\,\frac{32}{61})$; (b) $x=(-\frac{88}{61};\,-\frac{41}{61};\,\frac{12}{61})$
\item (a) $x=(2;\,2)$; (b) $x=(-3;\,-1)$; (c) $x=(18;\,5)$
\item (a) $x=(\frac{89}{39};\,\frac{37}{39};\,-\frac{5}{39})$; (b) $x=(\frac{8}{39};\,-\frac{44}{39};\,-\frac{71}{39})$
\item (a) $x=(1;\,1)$; (b) $x=(-\frac{3}{13};\,-\frac{27}{13})$; (c) $x=(2;\,1)$
\item (a) $x=(-\frac{10}{17};\,\frac{38}{17};\,-\frac{37}{17})$; (b) $x=(-\frac{7}{17};\,-\frac{21}{17};\,\frac{37}{17})$
\item (a) $x=(-1;\,-1)$; (b) $x=(-5;\,-4)$; (c) $x=(7;\,8)$
\item (a) $x=(-\frac{6}{13};\,-\frac{8}{13};\,\frac{1}{13})$; (b) $x=(-\frac{11}{13};\,\frac{1}{26};\,-\frac{22}{13})$
\item (a) $x=(3;\,3)$; (b) $x=(3;\,8)$; (c) $x=(3;\,-1)$
\end{enumerate}

% !Mode:: "TeX:UTF-8"
\javobmashgulot{9}{The Gauss and Gauss–Jordan methods for systems of linear equations}
\par\noindent\emph{Exercise set 1.}\nopagebreak
\begin{multicols}{3}\raggedright
\begin{enumerate}[label=\arabic*., leftmargin=*, itemsep=1pt]
\item $x_1=1,\ x_2=2$
\item $x_1=2,\ x_2=1$
\item $x_1=\frac{23}{9},\ x_2=\frac{17}{9}$
\item $x_1=\frac{31}{13},\ x_2=\frac{38}{13}$
\item $x_1=2,\ x_2=3$
\item $x_1=\frac{25}{11},\ x_2=\frac{18}{11}$
\item $x_1=\frac{8}{11},\ x_2=\frac{27}{11}$
\item $x_1=\frac{7}{5},\ x_2=\frac{8}{5}$
\item $x_1=\frac{13}{9},\ x_2=\frac{19}{9}$
\item $x_1=3,\ x_2=1$
\item $x_1=\frac{35}{19},\ x_2=\frac{37}{19}$
\item $x_1=-1,\ x_2=-1$
\item $x_1=3,\ x_2=2$
\item $x_1=\frac{9}{7},\ x_2=\frac{12}{7}$
\item $x_1=\frac{34}{19},\ x_2=\frac{51}{19}$
\item $x_1=\frac{19}{9},\ x_2=\frac{23}{9}$
\item $x_1=\frac{23}{7},\ x_2=\frac{16}{7}$
\item $x_1=\frac{20}{13},\ x_2=\frac{27}{13}$
\item $x_1=\frac{23}{11},\ x_2=\frac{15}{11}$
\item $x_1=3,\ x_2=2$
\item $x_1=\frac{15}{7},\ x_2=\frac{22}{7}$
\item $x_1=\frac{37}{13},\ x_2=\frac{19}{13}$
\item $x_1=2,\ x_2=1$
\item $x_1=1,\ x_2=-1$
\item $x_1=\frac{21}{19},\ x_2=\frac{37}{19}$
\item $x_1=\frac{13}{8},\ x_2=\frac{17}{8}$
\item $x_1=\frac{13}{5},\ x_2=\frac{11}{5}$
\item $x_1=\frac{25}{11},\ x_2=\frac{20}{11}$
\item $x_1=\frac{25}{13},\ x_2=\frac{30}{13}$
\item $x_1=\frac{17}{11},\ x_2=\frac{16}{11}$
\end{enumerate}
\end{multicols}
\par\noindent\emph{Exercise set 2.}\nopagebreak
\begin{multicols}{2}\raggedright
\begin{enumerate}[label=\arabic*., leftmargin=*, itemsep=1pt]
\item $x_1=1,\ x_2=2,\ x_3=3$
\item no solution ($\operatorname{rank}A=2<\operatorname{rank}\bar A=3$)
\item $x_1=-3,\ x_2=5,\ x_3=4$
\item $x_1=4,\ x_2=3,\ x_3=2$
\item $x_1=\frac{1}{10},\ x_2=\frac{3}{2},\ x_3=\frac{31}{10}$
\item $x_1=\frac{6}{11},\ x_2=\frac{8}{11},\ x_3=\frac{29}{11}$
\item $x_1=\frac{20}{11},\ x_2=\frac{78}{55},\ x_3=\frac{59}{55}$
\item $x_1=\frac{7}{6},\ x_2=\frac{1}{3},\ x_3=\frac{13}{6}$
\item $x_1=2,\ x_2=1,\ x_3=3$
\item $x_1=\frac{5}{3},\ x_2=\frac{5}{3},\ x_3=2$
\item $x_1=\frac{59}{78},\ x_2=\frac{103}{78},\ x_3=\frac{145}{78}$
\item $x_1=\frac{8}{7},\ x_2=1,\ x_3=\frac{18}{7}$
\item $x_1=\frac{55}{32},\ x_2=\frac{79}{32},\ x_3=\frac{53}{32}$
\item $x_1=\frac{1}{2},\ x_2=\frac{13}{10},\ x_3=\frac{29}{10}$
\item $x_1=\frac{15}{8},\ x_2=\frac{13}{8},\ x_3=\frac{11}{8}$
\item $x_1=\frac{36}{25},\ x_2=\frac{13}{5},\ x_3=\frac{34}{25}$
\item no solution ($\operatorname{rank}A=2<\operatorname{rank}\bar A=3$)
\item $x_1=\frac{16}{11},\ x_2=\frac{23}{11},\ x_3=\frac{21}{11}$
\item $x_1=\frac{3}{2},\ x_2=\frac{1}{2},\ x_3=\frac{5}{2}$
\item $x_1=\frac{27}{11},\ x_2=\frac{20}{11},\ x_3=\frac{52}{11}$
\item $x_1=\frac{9}{5},\ x_2=\frac{46}{25},\ x_3=\frac{53}{25}$
\item $x_1=\frac{16}{13},\ x_2=\frac{11}{13},\ x_3=\frac{31}{13}$
\item $x_1=\frac{13}{5},\ x_2=\frac{6}{5},\ x_3=1$
\item $x_1=\frac{29}{16},\ x_2=\frac{25}{16},\ x_3=\frac{21}{16}$
\item $x_1=\frac{1}{2},\ x_2=\frac{3}{2},\ x_3=\frac{5}{2}$
\item no solution ($\operatorname{rank}A=2<\operatorname{rank}\bar A=3$)
\item $x_1=\frac{4}{3},\ x_2=\frac{5}{3},\ x_3=\frac{8}{3}$
\item $x_1=\frac{41}{25},\ x_2=\frac{53}{25},\ x_3=\frac{54}{25}$
\item $x_1=\frac{41}{15},\ x_2=\frac{5}{3},\ x_3=\frac{2}{5}$
\item $x_1=\frac{11}{7},\ x_2=\frac{15}{7},\ x_3=\frac{16}{7}$
\end{enumerate}
\end{multicols}
\par\noindent\emph{Exercise set 3.}\nopagebreak
\begin{multicols}{2}\raggedright
\begin{enumerate}[label=\arabic*., leftmargin=*, itemsep=1pt]
\item $x_1=-\frac{11}{5},\ x_2=\frac{24}{5},\ x_3=\frac{17}{5}$
\item $x_1=2,\ x_2=1$
\item $x_1=\frac{43}{19},\ x_2=\frac{32}{19},\ x_3=\frac{23}{19}$
\item $x_1=\frac{17}{7},\ x_2=\frac{15}{7}$
\item $x_1=\frac{11}{5},\ x_2=\frac{6}{5},\ x_3=-\frac{1}{5}$
\item $x_1=2,\ x_2=1,\ x_3=2$
\item $x_1=1,\ x_2=2$
\item $x_1=1,\ x_2=2,\ x_3=3$
\item $x_1=\frac{41}{25},\ x_2=\frac{53}{25},\ x_3=\frac{54}{25}$
\item $x_1=\frac{18}{7},\ x_2=\frac{17}{7}$
\item no solution ($\operatorname{rank}A=2<\operatorname{rank}\bar A=3$)
\item $x_1=1,\ x_2=3$
\item $x_1=\frac{16}{13},\ x_2=\frac{11}{13},\ x_3=\frac{31}{13}$
\item $x_1=1,\ x_2=5$
\item $x_1=\frac{8}{7},\ x_2=1,\ x_3=\frac{18}{7}$
\item $x_1=\frac{5}{3},\ x_2=\frac{5}{3},\ x_3=2$
\item $x_1=\frac{15}{7},\ x_2=\frac{11}{7}$
\item $x_1=\frac{16}{11},\ x_2=\frac{23}{11},\ x_3=\frac{21}{11}$
\item $x_1=\frac{27}{11},\ x_2=\frac{20}{11},\ x_3=\frac{52}{11}$
\item $x_1=\frac{7}{3},\ x_2=\frac{5}{3}$
\item $x_1=\frac{1}{2},\ x_2=\frac{3}{2},\ x_3=\frac{5}{2}$
\item $x_1=\frac{3}{2},\ x_2=\frac{1}{2},\ x_3=\frac{5}{2}$
\item $x_1=\frac{9}{4},\ x_2=\frac{7}{4}$
\item $x_1=\frac{36}{25},\ x_2=\frac{13}{5},\ x_3=\frac{34}{25}$
\item $x_1=\frac{13}{5},\ x_2=\frac{6}{5},\ x_3=1$
\item $x_1=\frac{35}{19},\ x_2=\frac{37}{19}$
\item no solution ($\operatorname{rank}A=2<\operatorname{rank}\bar A=3$)
\item $x_1=\frac{29}{16},\ x_2=\frac{25}{16},\ x_3=\frac{21}{16}$
\item $x_1=\frac{13}{11},\ x_2=\frac{25}{11}$
\item $x_1=\frac{4}{3},\ x_2=\frac{5}{3},\ x_3=\frac{8}{3}$
\end{enumerate}
\end{multicols}
\par\noindent\emph{Exercise set 4.}\nopagebreak
\begin{enumerate}[label=\arabic*., leftmargin=*, itemsep=1pt]
\item infinitely many solutions: $x_1=- 3 x_{3} - 1,\ x_2=2 x_{3} + 2$ ($x_{3}$ arbitrary)
\item no solution ($\operatorname{rank}A=1<\operatorname{rank}\bar A=2$)
\item infinitely many solutions: $x_1=- x_{2} - x_{3} + 1$ ($x_{2}, x_{3}$ arbitrary)
\item infinitely many solutions: $x_1=5 - 2 x_{2}$ ($x_{2}$ arbitrary)
\item no solution ($\operatorname{rank}A=1<\operatorname{rank}\bar A=2$)
\item infinitely many solutions: $x_1=1,\ x_2=x_{3} - x_{4} - 1$ ($x_{3}, x_{4}$ arbitrary)
\item infinitely many solutions: $x_1=\frac{x_{2}}{3} + \frac{5}{3}$ ($x_{2}$ arbitrary)
\item infinitely many solutions: $x_1=x_{3} - \frac{1}{3},\ x_2=\frac{2}{3} - 2 x_{3}$ ($x_{3}$ arbitrary)
\item no solution ($\operatorname{rank}A=1<\operatorname{rank}\bar A=2$)
\item infinitely many solutions: $x_1=2 - x_{4},\ x_2=x_{3} - x_{4} + 1$ ($x_{3}, x_{4}$ arbitrary)
\item infinitely many solutions: $x_1=- 3 x_{2} + 2 x_{3}$ ($x_{2}, x_{3}$ arbitrary)
\item no solution ($\operatorname{rank}A=1<\operatorname{rank}\bar A=2$)
\item infinitely many solutions: $x_1=2 - \frac{x_{3}}{2},\ x_2=\frac{3 x_{3}}{2} - 1$ ($x_{3}$ arbitrary)
\item infinitely many solutions: $x_1=- \frac{x_{2}}{2} + \frac{x_{3}}{2} + 2$ ($x_{2}, x_{3}$ arbitrary)
\item no solution ($\operatorname{rank}A=1<\operatorname{rank}\bar A=2$)
\item infinitely many solutions: $x_1=- 2 x_{2} - x_{3} - x_{4} + 5$ ($x_{2}, x_{3}, x_{4}$ arbitrary)
\item $x_1=3,\ x_2=1$
\item infinitely many solutions: $x_1=- 2 x_{2} + x_{3} + 1$ ($x_{2}, x_{3}$ arbitrary)
\item no solution ($\operatorname{rank}A=1<\operatorname{rank}\bar A=2$)
\item infinitely many solutions: $x_1=- \frac{11 x_{4}}{5} - 1,\ x_2=- \frac{x_{4}}{5} - 1,\ x_3=\frac{7 x_{4}}{5} + 2$ ($x_{4}$ arbitrary)
\item infinitely many solutions: $x_1=- 2 x_{2} - 3 x_{3}$ ($x_{2}, x_{3}$ arbitrary)
\item infinitely many solutions: $x_1=x_{2} + 2$ ($x_{2}$ arbitrary)
\item no solution ($\operatorname{rank}A=1<\operatorname{rank}\bar A=2$)
\item infinitely many solutions: $x_1=- x_{2} + x_{3} + 3$ ($x_{2}, x_{3}$ arbitrary)
\item infinitely many solutions: $x_1=- 2 x_{2} + x_{3} + 1,\ x_4=1$ ($x_{2}, x_{3}$ arbitrary)
\item infinitely many solutions: $x_1=3 x_{2} + 5$ ($x_{2}$ arbitrary)
\item no solution ($\operatorname{rank}A=1<\operatorname{rank}\bar A=2$)
\item infinitely many solutions: $x_1=3 - 2 x_{2}$ ($x_{2}$ arbitrary)
\item no solution ($\operatorname{rank}A=1<\operatorname{rank}\bar A=2$)
\item infinitely many solutions: $x_1=- 4 x_{3} - 1,\ x_2=3 x_{3} + 2$ ($x_{3}$ arbitrary)
\end{enumerate}

% !Mode:: "TeX:UTF-8"
\javobmashgulot{10}{Arithmetic operations on vectors: multiplication of a vector by a number, addition and subtraction of vectors}
\par\noindent\emph{Exercise set 1.}\nopagebreak
\begin{multicols}{3}\raggedright
\begin{enumerate}[label=\arabic*., leftmargin=*, itemsep=1pt]
\item $(-2;\,22;\,-17)$
\item $(7;\,16;\,-14)$
\item $(1;\,3;\,1)$
\item $(13;\,-7;\,5)$
\item $(11;\,-2;\,7)$
\item $(-1;\,-1;\,12)$
\item $(12;\,4;\,-13)$
\item $(1;\,6;\,-9)$
\item $(7;\,17;\,-13)$
\item $(-4;\,6;\,-14)$
\item $(4;\,-4;\,8)$
\item $(-2;\,6;\,0)$
\item $(-6;\,18;\,13)$
\item $(5;\,-4;\,7)$
\item $(10;\,11;\,14)$
\item $(-7;\,8;\,-11)$
\item $(-1;\,3;\,2)$
\item $(6;\,-18;\,6)$
\item $(-4;\,17;\,13)$
\item $(-3;\,-6;\,3)$
\item $(-5;\,2;\,2)$
\item $(-13;\,-9;\,17)$
\item $(18;\,0;\,-6)$
\item $(10;\,2;\,-7)$
\item $(-3;\,3;\,0)$
\item $(4;\,7;\,6)$
\item $(13;\,-8;\,-8)$
\item $(2;\,2;\,-4)$
\item $(15;\,-19;\,-1)$
\item $(15;\,12;\,6)$
\end{enumerate}
\end{multicols}
\par\noindent\emph{Exercise set 2.}\nopagebreak
\begin{multicols}{2}\raggedright
\begin{enumerate}[label=\arabic*., leftmargin=*, itemsep=1pt]
\item $\overrightarrow{AB}=(2;\,2;\,7)$, $|\overrightarrow{AB}|=\sqrt{57}$, $M(0;\,-1;\,\frac{5}{2})$
\item $\overrightarrow{AB}=(1;\,-3;\,0)$, $|\overrightarrow{AB}|=\sqrt{10}$, $M(-\frac{5}{2};\,\frac{5}{2};\,5)$
\item $\overrightarrow{AB}=(4;\,-4;\,-1)$, $|\overrightarrow{AB}|=\sqrt{33}$, $M(0;\,0;\,\frac{1}{2})$
\item $\overrightarrow{AB}=(-2;\,1;\,-1)$, $|\overrightarrow{AB}|=\sqrt{6}$, $M(0;\,\frac{9}{2};\,-\frac{3}{2})$
\item $\overrightarrow{AB}=(2;\,0;\,-6)$, $|\overrightarrow{AB}|=2 \sqrt{10}$, $M(2;\,1;\,-2)$
\item $\overrightarrow{AB}=(-6;\,4;\,0)$, $|\overrightarrow{AB}|=2 \sqrt{13}$, $M(2;\,4;\,5)$
\item $\overrightarrow{AB}=(3;\,-6;\,9)$, $|\overrightarrow{AB}|=3 \sqrt{14}$, $M(\frac{3}{2};\,3;\,\frac{1}{2})$
\item $\overrightarrow{AB}=(2;\,3;\,-9)$, $|\overrightarrow{AB}|=\sqrt{94}$, $M(4;\,-\frac{7}{2};\,-\frac{3}{2})$
\item $\overrightarrow{AB}=(-4;\,-4;\,1)$, $|\overrightarrow{AB}|=\sqrt{33}$, $M(1;\,-3;\,\frac{9}{2})$
\item $\overrightarrow{AB}=(-2;\,0;\,-3)$, $|\overrightarrow{AB}|=\sqrt{13}$, $M(-1;\,3;\,-\frac{5}{2})$
\item $\overrightarrow{AB}=(-2;\,6;\,8)$, $|\overrightarrow{AB}|=2 \sqrt{26}$, $M(-4;\,-1;\,2)$
\item $\overrightarrow{AB}=(11;\,-2;\,-3)$, $|\overrightarrow{AB}|=\sqrt{134}$, $M(-\frac{1}{2};\,-4;\,\frac{3}{2})$
\item $\overrightarrow{AB}=(4;\,3;\,0)$, $|\overrightarrow{AB}|=5$, $M(-1;\,\frac{9}{2};\,3)$
\item $\overrightarrow{AB}=(5;\,3;\,0)$, $|\overrightarrow{AB}|=\sqrt{34}$, $M(-\frac{7}{2};\,\frac{3}{2};\,-3)$
\item $\overrightarrow{AB}=(4;\,0;\,-4)$, $|\overrightarrow{AB}|=4 \sqrt{2}$, $M(4;\,6;\,1)$
\item $\overrightarrow{AB}=(9;\,-10;\,-11)$, $|\overrightarrow{AB}|=\sqrt{302}$, $M(-\frac{3}{2};\,-1;\,-\frac{1}{2})$
\item $\overrightarrow{AB}=(-4;\,4;\,4)$, $|\overrightarrow{AB}|=4 \sqrt{3}$, $M(1;\,1;\,2)$
\item $\overrightarrow{AB}=(-1;\,-3;\,10)$, $|\overrightarrow{AB}|=\sqrt{110}$, $M(-\frac{9}{2};\,\frac{9}{2};\,1)$
\item $\overrightarrow{AB}=(6;\,6;\,8)$, $|\overrightarrow{AB}|=2 \sqrt{34}$, $M(2;\,-1;\,-2)$
\item $\overrightarrow{AB}=(-6;\,2;\,8)$, $|\overrightarrow{AB}|=2 \sqrt{26}$, $M(3;\,0;\,-1)$
\item $\overrightarrow{AB}=(2;\,11;\,5)$, $|\overrightarrow{AB}|=5 \sqrt{6}$, $M(0;\,-\frac{1}{2};\,-\frac{3}{2})$
\item $\overrightarrow{AB}=(-2;\,2;\,-9)$, $|\overrightarrow{AB}|=\sqrt{89}$, $M(2;\,2;\,\frac{1}{2})$
\item $\overrightarrow{AB}=(3;\,2;\,2)$, $|\overrightarrow{AB}|=\sqrt{17}$, $M(\frac{1}{2};\,1;\,0)$
\item $\overrightarrow{AB}=(0;\,7;\,1)$, $|\overrightarrow{AB}|=5 \sqrt{2}$, $M(0;\,\frac{1}{2};\,-\frac{9}{2})$
\item $\overrightarrow{AB}=(10;\,1;\,-1)$, $|\overrightarrow{AB}|=\sqrt{102}$, $M(-1;\,-\frac{7}{2};\,\frac{7}{2})$
\item $\overrightarrow{AB}=(2;\,-6;\,-5)$, $|\overrightarrow{AB}|=\sqrt{65}$, $M(1;\,3;\,-\frac{7}{2})$
\item $\overrightarrow{AB}=(-5;\,3;\,-2)$, $|\overrightarrow{AB}|=\sqrt{38}$, $M(\frac{1}{2};\,-\frac{9}{2};\,-5)$
\item $\overrightarrow{AB}=(1;\,7;\,4)$, $|\overrightarrow{AB}|=\sqrt{66}$, $M(\frac{7}{2};\,-\frac{1}{2};\,-3)$
\item $\overrightarrow{AB}=(4;\,5;\,12)$, $|\overrightarrow{AB}|=\sqrt{185}$, $M(-2;\,\frac{1}{2};\,0)$
\item $\overrightarrow{AB}=(6;\,2;\,2)$, $|\overrightarrow{AB}|=2 \sqrt{11}$, $M(3;\,-4;\,0)$
\end{enumerate}
\end{multicols}
\par\noindent\emph{Exercise set 3.}\nopagebreak
\begin{multicols}{2}\raggedright
\begin{enumerate}[label=\arabic*., leftmargin=*, itemsep=1pt]
\item $\vec a^{\,0}=(\frac23;-\frac13;\frac23)$
\item yes, $\vec b=-\frac32\vec a$; opposite directions
\item $m=2,\ n=6$
\item $D(2;-3;3)$
\item $G(2;1;1)$
\item $\cos\alpha=\frac13,\ \cos\beta=\frac23,\ \cos\gamma=-\frac23$
\item $\vec u-\vec v=(4;-2;-2;0)$, $\frac12(\vec u+\vec v)=(3;3;1;3)$
\item $\vec c=2\vec a+3\vec b$
\item $|\vec a+\vec b|=|\vec a-\vec b|=5$
\item $M(3;0;2)$
\end{enumerate}
\end{multicols}

% !Mode:: "TeX:UTF-8"
\javobmashgulot{11}{Scalar, vector and mixed products of vectors}
\par\noindent\emph{Exercise set 1.}\nopagebreak
\begin{multicols}{3}\raggedright
\begin{enumerate}[label=\arabic*., leftmargin=*, itemsep=1pt]
\item $\vec a\cdot\vec b=21$, $\cos\varphi=\frac{21 \sqrt{493}}{493}$
\item $\vec a\cdot\vec b=-7$, $\cos\varphi=- \frac{7 \sqrt{110}}{330}$
\item $\vec a\cdot\vec b=7$, $\cos\varphi=\frac{7 \sqrt{754}}{754}$
\item $\vec a\cdot\vec b=7$, $\cos\varphi=\frac{\sqrt{1239}}{177}$
\item $\vec a\cdot\vec b=-24$, $\cos\varphi=- \frac{4 \sqrt{33}}{55}$
\item $\vec a\cdot\vec b=3$, $\cos\varphi=\frac{\sqrt{66}}{110}$
\item $\vec a\cdot\vec b=-2$, $\cos\varphi=- \frac{1}{15}$
\item $\vec a\cdot\vec b=-17$, $\cos\varphi=- \frac{\sqrt{510}}{30}$
\item $\vec a\cdot\vec b=19$, $\cos\varphi=\frac{19 \sqrt{55}}{330}$
\item $\vec a\cdot\vec b=7$, $\cos\varphi=\frac{7 \sqrt{130}}{130}$
\item $\vec a\cdot\vec b=4$, $\cos\varphi=\frac{\sqrt{105}}{105}$
\item $\vec a\cdot\vec b=22$, $\cos\varphi=\frac{11 \sqrt{21}}{63}$
\item $\vec a\cdot\vec b=11$, $\cos\varphi=\frac{\sqrt{22}}{18}$
\item $\vec a\cdot\vec b=5$, $\cos\varphi=\frac{5 \sqrt{19}}{57}$
\item $\vec a\cdot\vec b=-19$, $\cos\varphi=- \frac{19 \sqrt{510}}{510}$
\item $\vec a\cdot\vec b=13$, $\cos\varphi=\frac{\sqrt{754}}{58}$
\item $\vec a\cdot\vec b=-3$, $\cos\varphi=- \frac{\sqrt{1122}}{374}$
\item $\vec a\cdot\vec b=23$, $\cos\varphi=\frac{23 \sqrt{646}}{646}$
\item $\vec a\cdot\vec b=-2$, $\cos\varphi=- \frac{\sqrt{70}}{105}$
\item $\vec a\cdot\vec b=-4$, $\cos\varphi=- \frac{2 \sqrt{15}}{45}$
\item $\vec a\cdot\vec b=-28$, $\cos\varphi=- \frac{2 \sqrt{399}}{57}$
\item $\vec a\cdot\vec b=-14$, $\cos\varphi=- \frac{14 \sqrt{969}}{969}$
\item $\vec a\cdot\vec b=9$, $\cos\varphi=\frac{3 \sqrt{21}}{14}$
\item $\vec a\cdot\vec b=-28$, $\cos\varphi=- \frac{7 \sqrt{15}}{45}$
\item $\vec a\cdot\vec b=16$, $\cos\varphi=\frac{4 \sqrt{5}}{15}$
\item $\vec a\cdot\vec b=-20$, $\cos\varphi=- \frac{20}{29}$
\item $\vec a\cdot\vec b=-10$, $\cos\varphi=- \frac{\sqrt{10}}{7}$
\item $\vec a\cdot\vec b=-28$, $\cos\varphi=- \frac{7 \sqrt{258}}{129}$
\item $\vec a\cdot\vec b=12$, $\cos\varphi=\frac{\sqrt{30}}{15}$
\item $\vec a\cdot\vec b=9$, $\cos\varphi=\frac{3 \sqrt{21}}{70}$
\end{enumerate}
\end{multicols}
\par\noindent\emph{Exercise set 2.}\nopagebreak
\begin{multicols}{2}\raggedright
\begin{enumerate}[label=\arabic*., leftmargin=*, itemsep=1pt]
\item $\vec a\times\vec b=(-7;\,0;\,0)$, $S=7$
\item $\vec a\times\vec b=(-8;\,6;\,12)$, $S=2 \sqrt{61}$
\item $\vec a\times\vec b=(-7;\,3;\,-3)$, $S=\sqrt{67}$
\item $\vec a\times\vec b=(3;\,-4;\,-6)$, $S=\sqrt{61}$
\item $\vec a\times\vec b=(-14;\,16;\,-4)$, $S=6 \sqrt{13}$
\item $\vec a\times\vec b=(-6;\,0;\,-12)$, $S=6 \sqrt{5}$
\item $\vec a\times\vec b=(10;\,12;\,4)$, $S=2 \sqrt{65}$
\item $\vec a\times\vec b=(24;\,12;\,-12)$, $S=12 \sqrt{6}$
\item $\vec a\times\vec b=(-3;\,6;\,3)$, $S=3 \sqrt{6}$
\item $\vec a\times\vec b=(5;\,4;\,6)$, $S=\sqrt{77}$
\item $\vec a\times\vec b=(-3;\,3;\,6)$, $S=3 \sqrt{6}$
\item $\vec a\times\vec b=(0;\,-20;\,20)$, $S=20 \sqrt{2}$
\item $\vec a\times\vec b=(-6;\,6;\,3)$, $S=9$
\item $\vec a\times\vec b=(21;\,0;\,-21)$, $S=21 \sqrt{2}$
\item $\vec a\times\vec b=(16;\,16;\,0)$, $S=16 \sqrt{2}$
\item $\vec a\times\vec b=(12;\,5;\,9)$, $S=5 \sqrt{10}$
\item $\vec a\times\vec b=(8;\,4;\,4)$, $S=4 \sqrt{6}$
\item $\vec a\times\vec b=(3;\,-5;\,4)$, $S=5 \sqrt{2}$
\item $\vec a\times\vec b=(0;\,-2;\,-1)$, $S=\sqrt{5}$
\item $\vec a\times\vec b=(1;\,5;\,4)$, $S=\sqrt{42}$
\item $\vec a\times\vec b=(6;\,3;\,-2)$, $S=7$
\item $\vec a\times\vec b=(6;\,-6;\,9)$, $S=3 \sqrt{17}$
\item $\vec a\times\vec b=(-10;\,15;\,-5)$, $S=5 \sqrt{14}$
\item $\vec a\times\vec b=(4;\,0;\,-4)$, $S=4 \sqrt{2}$
\item $\vec a\times\vec b=(10;\,-10;\,10)$, $S=10 \sqrt{3}$
\item $\vec a\times\vec b=(-9;\,8;\,-4)$, $S=\sqrt{161}$
\item $\vec a\times\vec b=(-3;\,-6;\,15)$, $S=3 \sqrt{30}$
\item $\vec a\times\vec b=(6;\,-4;\,9)$, $S=\sqrt{133}$
\item $\vec a\times\vec b=(-7;\,16;\,-4)$, $S=\sqrt{321}$
\item $\vec a\times\vec b=(-14;\,8;\,-12)$, $S=2 \sqrt{101}$
\end{enumerate}
\end{multicols}
\par\noindent\emph{Exercise set 3.}\nopagebreak
\begin{multicols}{2}\raggedright
\begin{enumerate}[label=\arabic*., leftmargin=*, itemsep=1pt]
\item $(\vec a,\vec b,\vec c)=-12$; not coplanar, $V=12$
\item $(\vec a,\vec b,\vec c)=4$; not coplanar, $V=4$
\item $(\vec a,\vec b,\vec c)=65$; not coplanar, $V=65$
\item $(\vec a,\vec b,\vec c)=-17$; not coplanar, $V=17$
\item $(\vec a,\vec b,\vec c)=-26$; not coplanar, $V=26$
\item $(\vec a,\vec b,\vec c)=0$; coplanar, $V=0$
\item $V=\frac16|-4|=\frac{2}{3}$
\item $\lambda=5$
\item $7m-10=0$, $m=\frac{10}{7}$
\item $S=\frac12|\overrightarrow{AB}\times\overrightarrow{AC}|=\frac72$
\end{enumerate}
\end{multicols}

% !Mode:: "TeX:UTF-8"
\javobmashgulot{12}{Equations of a straight line in the plane in various forms. The angle between two lines. The distance from a point to a line}
\par\noindent\emph{Exercise set 1.}\nopagebreak
\begin{multicols}{4}\raggedright
\begin{enumerate}[label=\arabic*., leftmargin=*, itemsep=1pt]
\item $4x-3y-10=0$
\item $x+8y+4=0$
\item $x+7y+13=0$
\item $x-3y+14=0$
\item $3x-y-12=0$
\item $x-4y+8=0$
\item $x+1=0$
\item $8x+3y+39=0$
\item $y-4=0$
\item $12x-5y+30=0$
\item $7x+8y+10=0$
\item $5x+6y+25=0$
\item $x+4y-4=0$
\item $7x+3y-2=0$
\item $2x+3y+8=0$
\item $4x+y-3=0$
\item $x+y-4=0$
\item $x-y+2=0$
\item $9x+11y+10=0$
\item $8x+7y-18=0$
\item $x-6y-5=0$
\item $y+2=0$
\item $y+1=0$
\item $3x-4y+6=0$
\item $x+2y+6=0$
\item $y+6=0$
\item $10x+3y+18=0$
\item $x-y-3=0$
\item $9x+7y+3=0$
\item $x-3y-11=0$
\end{enumerate}
\end{multicols}
\par\noindent\emph{Exercise set 2.}\nopagebreak
\begin{multicols}{4}\raggedright
\begin{enumerate}[label=\arabic*., leftmargin=*, itemsep=1pt]
\item $d=\frac{\sqrt{10}}{10}$
\item $d=\frac{14 \sqrt{41}}{41}$
\item $d=\frac{2 \sqrt{10}}{5}$
\item $d=\frac{11 \sqrt{26}}{26}$
\item $d=\sqrt{5}$
\item $d=\frac{3 \sqrt{61}}{61}$
\item $d=\frac{\sqrt{10}}{4}$
\item $d=\frac{11 \sqrt{2}}{12}$
\item $d=\frac{2 \sqrt{34}}{17}$
\item $d=\frac{\sqrt{13}}{2}$
\item $d=\frac{5 \sqrt{41}}{41}$
\item $d=3 \sqrt{2}$
\item $d=\frac{7 \sqrt{26}}{26}$
\item $d=\frac{19 \sqrt{5}}{15}$
\item $d=\frac{8 \sqrt{41}}{41}$
\item $d=\frac{11 \sqrt{2}}{4}$
\item $d=\frac{16 \sqrt{61}}{61}$
\item $d=\frac{13}{5}$
\item $d=\frac{17 \sqrt{10}}{20}$
\item $d=\frac{7 \sqrt{2}}{6}$
\item $d=\frac{5 \sqrt{61}}{61}$
\item $d=\sqrt{2}$
\item $d=0$ ($M$ lies on the line)
\item $d=\frac{16 \sqrt{29}}{29}$
\item $d=\frac{7 \sqrt{26}}{26}$
\item $d=\frac{40 \sqrt{61}}{61}$
\item $d=\frac{7 \sqrt{34}}{17}$
\item $d=\frac{15 \sqrt{41}}{41}$
\item $d=\frac{8 \sqrt{26}}{13}$
\item $d=\frac{2 \sqrt{17}}{17}$
\end{enumerate}
\end{multicols}
\par\noindent\emph{Exercise set 3.}\nopagebreak
\begin{multicols}{3}\raggedright
\begin{enumerate}[label=\arabic*., leftmargin=*, itemsep=1pt]
\item perpendicular ($\varphi=90^\circ$)
\item parallel ($\varphi=0$)
\item $\cos\varphi=\frac{\sqrt{2}}{10}$, $\varphi\approx81.9^\circ$
\item perpendicular ($\varphi=90^\circ$)
\item $\cos\varphi=\frac{5 \sqrt{221}}{221}$, $\varphi\approx70.3^\circ$
\item $3x-4y-18=0$
\item $x-2y+7=0$
\item $\frac{x}{-4}+\frac{y}{6}=1$
\item $d=5$
\item $AB:\ 4x-3y-1=0$, $h=\frac{24}{5}$
\end{enumerate}
\end{multicols}


\section*{Chapter 3. Lines and Planes in Space; Linear and Euclidean Spaces}
% !Mode:: "TeX:UTF-8"
\javobmashgulot{13}{Lines in space. Equations of a line in space, the angle between lines, the distance between lines}
\par\noindent\emph{Exercise set 1.}\nopagebreak
\begin{multicols}{3}\raggedright
\begin{enumerate}[label=\arabic*., leftmargin=*, itemsep=1pt]
\item $\frac{x+1}{6}=\frac{y+1}{-2}=\frac{z-5}{0}$
\item $\frac{x+3}{1}=\frac{y+4}{3}=\frac{z-3}{-8}$
\item $\frac{x-1}{3}=\frac{y+3}{-2}=\frac{z-5}{-6}$
\item $\frac{x+3}{5}=\frac{y+4}{5}=\frac{z+1}{-2}$
\item $\frac{x+1}{3}=\frac{y}{3}=\frac{z+2}{6}$
\item $\frac{x-5}{-1}=\frac{y+4}{9}=\frac{z}{-1}$
\item $\frac{x-2}{-5}=\frac{y-3}{-1}=\frac{z-4}{1}$
\item $\frac{x-2}{-3}=\frac{y-3}{-5}=\frac{z+3}{0}$
\item $\frac{x-3}{-3}=\frac{y}{2}=\frac{z+1}{0}$
\item $\frac{x-4}{-7}=\frac{y+1}{5}=\frac{z-1}{1}$
\item $\frac{x-3}{-5}=\frac{y+2}{0}=\frac{z-4}{-4}$
\item $\frac{x+3}{8}=\frac{y+4}{9}=\frac{z-1}{1}$
\item $\frac{x-1}{-6}=\frac{y+5}{3}=\frac{z-1}{-4}$
\item $\frac{x-2}{-3}=\frac{y-2}{-5}=\frac{z-5}{-4}$
\item $\frac{x+1}{6}=\frac{y}{0}=\frac{z+2}{5}$
\item $\frac{x-5}{-7}=\frac{y+3}{5}=\frac{z-4}{-6}$
\item $\frac{x-4}{0}=\frac{y+5}{0}=\frac{z-4}{-4}$
\item $\frac{x-4}{-7}=\frac{y-4}{-5}=\frac{z}{-2}$
\item $\frac{x+3}{-1}=\frac{y-5}{-7}=\frac{z-3}{-6}$
\item $\frac{x+4}{8}=\frac{y-1}{2}=\frac{z-5}{-2}$
\item $\frac{x+5}{10}=\frac{y-2}{-6}=\frac{z}{3}$
\item $\frac{x+5}{1}=\frac{y+3}{4}=\frac{z+1}{0}$
\item $\frac{x+1}{-3}=\frac{y-4}{-2}=\frac{z-4}{-1}$
\item $\frac{x}{0}=\frac{y+5}{5}=\frac{z+5}{2}$
\item $\frac{x-2}{1}=\frac{y-3}{-1}=\frac{z-3}{-5}$
\item $\frac{x-5}{-10}=\frac{y+2}{3}=\frac{z+1}{4}$
\item $\frac{x+3}{4}=\frac{y+3}{8}=\frac{z+1}{-1}$
\item $\frac{x+1}{-1}=\frac{y-1}{-1}=\frac{z-4}{-6}$
\item $\frac{x-4}{-6}=\frac{y+1}{6}=\frac{z-1}{-1}$
\item $\frac{x-1}{0}=\frac{y+5}{7}=\frac{z}{-4}$
\end{enumerate}
\end{multicols}
\par\noindent\emph{Exercise set 2.}\nopagebreak
\begin{multicols}{4}\raggedright
\begin{enumerate}[label=\arabic*., leftmargin=*, itemsep=1pt]
\item $\cos\varphi=\frac{4 \sqrt{410}}{205}$
\item $\cos\varphi=\frac{5 \sqrt{174}}{174}$
\item $\cos\varphi=\frac{\sqrt{70}}{10}$
\item $\cos\varphi=\frac{3 \sqrt{165}}{55}$
\item $\cos\varphi=\frac{2 \sqrt{266}}{133}$
\item $\cos\varphi=\frac{\sqrt{7}}{7}$
\item $\varphi=90^\circ$
\item $\cos\varphi=\frac{10 \sqrt{41}}{123}$
\item $\cos\varphi=\frac{\sqrt{17}}{34}$
\item $\cos\varphi=\frac{5 \sqrt{34}}{34}$
\item $\cos\varphi=\frac{3 \sqrt{102}}{34}$
\item $\cos\varphi=\frac{19 \sqrt{462}}{462}$
\item $\cos\varphi=\frac{3 \sqrt{170}}{170}$
\item $\varphi=90^\circ$
\item $\cos\varphi=\frac{\sqrt{130}}{130}$
\item $\cos\varphi=\frac{2 \sqrt{858}}{429}$
\item $\cos\varphi=\frac{2}{7}$
\item $\cos\varphi=\frac{12 \sqrt{221}}{221}$
\item $\cos\varphi=\frac{5 \sqrt{574}}{574}$
\item $\cos\varphi=\frac{11 \sqrt{14}}{70}$
\item $\cos\varphi=\frac{6 \sqrt{377}}{377}$
\item $\cos\varphi=\frac{8 \sqrt{357}}{357}$
\item $\cos\varphi=\frac{7 \sqrt{374}}{187}$
\item $\cos\varphi=\frac{7 \sqrt{6}}{18}$
\item $\cos\varphi=\frac{5 \sqrt{638}}{638}$
\item $\cos\varphi=\frac{\sqrt{58}}{58}$
\item $\cos\varphi=\frac{\sqrt{2}}{10}$
\item $\cos\varphi=\frac{9 \sqrt{130}}{130}$
\item $\cos\varphi=\frac{\sqrt{114}}{57}$
\item $\cos\varphi=\frac{\sqrt{3}}{15}$
\end{enumerate}
\end{multicols}
\par\noindent\emph{Exercise set 3.}\nopagebreak
\begin{multicols}{2}\raggedright
\begin{enumerate}[label=\arabic*., leftmargin=*, itemsep=1pt]
\item $d=\sqrt{13}$
\item $d=\frac{13 \sqrt{34}}{17}$
\item $d=\frac{\sqrt{1115}}{5}$
\item $d=\frac{\sqrt{70}}{5}$
\item $d=\frac{2 \sqrt{66}}{11}$
\item $d=\frac{2 \sqrt{874}}{19}$
\item $d=\frac{7 \sqrt{30}}{5}$
\item $x=1+2t,\ y=-3-t,\ z=4t$; $(5;-5;8)$
\item yes, at the point $(2;4;4)$
\item $d=2$
\end{enumerate}
\end{multicols}

% !Mode:: "TeX:UTF-8"
\javobmashgulot{14}{Planes in space. The general equation of a plane. Relative position of lines and planes in space}
\par\noindent\emph{Exercise set 1.}\nopagebreak
\begin{multicols}{3}\raggedright
\begin{enumerate}[label=\arabic*., leftmargin=*, itemsep=1pt]
\item $x+y+z=0$
\item $x-6y+2z+5=0$
\item $16x-9y+24z+68=0$
\item $2x-3y-z+5=0$
\item $4x-15y-2z-31=0$
\item $x-z=0$
\item $20x+15y+22z+66=0$
\item $5x+20y+2z-66=0$
\item $x-3y+z+7=0$
\item $7x+14y+13z+10=0$
\item $3x+8y-z-8=0$
\item $3x-6y-5z+6=0$
\item $61x-20y+28z+52=0$
\item $11x+13y-3z+23=0$
\item $7x+2y-z+5=0$
\item $3x+8y+7z-41=0$
\item $9x-4y+11z+10=0$
\item $7x+10y-8z+31=0$
\item $37x-2y+12z-46=0$
\item $4x-13y+7z-4=0$
\item $5x+4y-11z+18=0$
\item $z+3=0$
\item $4x-7y+3z+28=0$
\item $5x-y+2z-3=0$
\item $13x+2y-3z+41=0$
\item $4x+3y-20z+8=0$
\item $7x-3y+20z-41=0$
\item $y-2=0$
\item $2y+z+3=0$
\item $14x+y+11z-8=0$
\end{enumerate}
\end{multicols}
\par\noindent\emph{Exercise set 2.}\nopagebreak
\begin{multicols}{4}\raggedright
\begin{enumerate}[label=\arabic*., leftmargin=*, itemsep=1pt]
\item $d=\frac{4 \sqrt{66}}{11}$
\item $d=\frac{27 \sqrt{34}}{34}$
\item $d=\frac{11 \sqrt{3}}{3}$
\item $d=\frac{\sqrt{26}}{26}$
\item $d=\frac{2 \sqrt{14}}{7}$
\item $d=\frac{4 \sqrt{6}}{3}$
\item $d=\frac{21 \sqrt{22}}{44}$
\item $d=\frac{\sqrt{22}}{22}$
\item $d=\frac{33 \sqrt{53}}{53}$
\item $d=\frac{\sqrt{6}}{2}$
\item $d=2 \sqrt{6}$
\item $d=\frac{15 \sqrt{38}}{38}$
\item $d=\frac{40 \sqrt{61}}{61}$
\item $d=\frac{12 \sqrt{38}}{19}$
\item $d=\frac{25 \sqrt{11}}{22}$
\item $d=\frac{23 \sqrt{42}}{21}$
\item $d=\frac{25 \sqrt{22}}{22}$
\item $d=0$
\item $d=\frac{23 \sqrt{42}}{42}$
\item $d=\frac{26 \sqrt{3}}{9}$
\item $d=\frac{17 \sqrt{6}}{6}$
\item $d=\frac{21 \sqrt{22}}{22}$
\item $d=\frac{13 \sqrt{10}}{10}$
\item $d=\frac{5 \sqrt{11}}{11}$
\item $d=\frac{3 \sqrt{65}}{13}$
\item $d=2 \sqrt{2}$
\item $d=\frac{16 \sqrt{61}}{61}$
\item $d=\frac{29 \sqrt{43}}{43}$
\item $d=\frac{33 \sqrt{19}}{38}$
\item $d=\frac{10 \sqrt{26}}{13}$
\end{enumerate}
\end{multicols}
\par\noindent\emph{Exercise set 3.}\nopagebreak
\begin{multicols}{3}\raggedright
\begin{enumerate}[label=\arabic*., leftmargin=*, itemsep=1pt]
\item $\cos\varphi=\frac{9 \sqrt{95}}{95}$
\item $\cos\varphi=\frac{9}{13}$
\item $\cos\varphi=\frac{5 \sqrt{3}}{9}$
\item $\cos\varphi=\frac{\sqrt{14}}{14}$
\item perpendicular ($90^\circ$)
\item $t=\frac{7}{3}$: $\left(\frac{17}{3};\frac43;-\frac73\right)$
\item $x-2y+4z-16=0$
\item $3x-y+2z+1=0$
\item $\frac{x}{6}+\frac{y}{-4}+\frac{z}{2}=1$
\item $d=4$
\end{enumerate}
\end{multicols}

% !Mode:: "TeX:UTF-8"
\javobmashgulot{15}{Expansion of a vector of a linear space in a basis. Subspaces of a linear space}
\par\noindent\emph{Exercise set 1.}\nopagebreak
\begin{multicols}{3}\raggedright
\begin{enumerate}[label=\arabic*., leftmargin=*, itemsep=1pt]
\item $\vec x=-2\vec e_2+\vec e_3$
\item $\vec x=\vec e_1-4\vec e_2+\vec e_3$
\item $\vec x=3\vec e_1-3\vec e_2+2\vec e_3$
\item $\vec x=2\vec e_1-2\vec e_2+\vec e_3$
\item $\vec x=-3\vec e_1+2\vec e_2+\vec e_3$
\item $\vec x=-2\vec e_1+3\vec e_2+4\vec e_3$
\item $\vec x=3\vec e_1+2\vec e_2+3\vec e_3$
\item $\vec x=-\vec e_2-\vec e_3$
\item $\vec x=3\vec e_1-\vec e_2+2\vec e_3$
\item $\vec x=-2\vec e_1-\vec e_2-2\vec e_3$
\item $\vec x=\vec e_1+4\vec e_2-4\vec e_3$
\item $\vec x=3\vec e_2$
\item $\vec x=2\vec e_1-2\vec e_2-3\vec e_3$
\item $\vec x=-3\vec e_1$
\item $\vec x=-3\vec e_1-2\vec e_2-4\vec e_3$
\item $\vec x=2\vec e_2-\vec e_3$
\item $\vec x=-2\vec e_1-3\vec e_2+4\vec e_3$
\item $\vec x=2\vec e_1+4\vec e_2-2\vec e_3$
\item $\vec x=-3\vec e_1+2\vec e_2+3\vec e_3$
\item $\vec x=2\vec e_1+4\vec e_3$
\item $\vec x=3\vec e_1+4\vec e_2-2\vec e_3$
\item $\vec x=3\vec e_1+3\vec e_2+4\vec e_3$
\item $\vec x=2\vec e_1+3\vec e_2-3\vec e_3$
\item $\vec x=2\vec e_1-3\vec e_2+4\vec e_3$
\item $\vec x=-3\vec e_1+\vec e_2$
\item $\vec x=-3\vec e_1+3\vec e_2+4\vec e_3$
\item $\vec x=-2\vec e_2+4\vec e_3$
\item $\vec x=-\vec e_1-\vec e_2-3\vec e_3$
\item $\vec x=-\vec e_1+4\vec e_2+\vec e_3$
\item $\vec x=-4\vec e_2+\vec e_3$
\end{enumerate}
\end{multicols}
\par\noindent\emph{Exercise set 2.}\nopagebreak
\begin{multicols}{2}\raggedright
\begin{enumerate}[label=\arabic*., leftmargin=*, itemsep=1pt]
\item $\operatorname{rank}=2$; linearly dependent, $\dim L=2$
\item $\operatorname{rank}=3$; linearly independent, $\dim L=3$
\item $\operatorname{rank}=2$; linearly dependent, $\dim L=2$
\item $\operatorname{rank}=3$; linearly independent, $\dim L=3$
\item $\operatorname{rank}=2$; linearly dependent, $\dim L=2$
\item $\operatorname{rank}=3$; linearly independent, $\dim L=3$
\item $\operatorname{rank}=2$; linearly dependent, $\dim L=2$
\item $\operatorname{rank}=3$; linearly independent, $\dim L=3$
\item $\operatorname{rank}=2$; linearly dependent, $\dim L=2$
\item $\operatorname{rank}=3$; linearly independent, $\dim L=3$
\item $\operatorname{rank}=1$; linearly dependent, $\dim L=1$
\item $\operatorname{rank}=3$; linearly independent, $\dim L=3$
\item $\operatorname{rank}=2$; linearly dependent, $\dim L=2$
\item $\operatorname{rank}=3$; linearly independent, $\dim L=3$
\item $\operatorname{rank}=2$; linearly dependent, $\dim L=2$
\item $\operatorname{rank}=3$; linearly independent, $\dim L=3$
\item $\operatorname{rank}=2$; linearly dependent, $\dim L=2$
\item $\operatorname{rank}=3$; linearly independent, $\dim L=3$
\item $\operatorname{rank}=2$; linearly dependent, $\dim L=2$
\item $\operatorname{rank}=3$; linearly independent, $\dim L=3$
\item $\operatorname{rank}=2$; linearly dependent, $\dim L=2$
\item $\operatorname{rank}=3$; linearly independent, $\dim L=3$
\item $\operatorname{rank}=2$; linearly dependent, $\dim L=2$
\item $\operatorname{rank}=3$; linearly independent, $\dim L=3$
\item $\operatorname{rank}=2$; linearly dependent, $\dim L=2$
\item $\operatorname{rank}=3$; linearly independent, $\dim L=3$
\item $\operatorname{rank}=2$; linearly dependent, $\dim L=2$
\item $\operatorname{rank}=1$; linearly dependent, $\dim L=1$
\item $\operatorname{rank}=2$; linearly dependent, $\dim L=2$
\item $\operatorname{rank}=3$; linearly independent, $\dim L=3$
\end{enumerate}
\end{multicols}
\par\noindent\emph{Exercise set 3.}\nopagebreak
\begin{multicols}{2}\raggedright
\begin{enumerate}[label=\arabic*., leftmargin=*, itemsep=1pt]
\item basis $(1;-1;0),\ (1;0;-1)$; $\dim L=2$
\item no: $(0;0)\notin M$
\item yes: $\vec x=2\vec a+\vec b$
\item no (the system is inconsistent)
\item $p=6+5(x-1)+(x-1)^2$
\item $4$; $E_{11},E_{12},E_{21},E_{22}$
\item $\vec x=3\vec e_1+2\vec e_2$
\item dependent: $\vec a+\vec b+\vec c=\vec 0 \pmod 2$
\item basis $(-2;1;0),\ (1;0;1)$; $\dim=2$
\item $2^2=4$ codewords: $0000,\ 1011,\ 0110,\ 1101$
\end{enumerate}
\end{multicols}

% !Mode:: "TeX:UTF-8"
\javobmashgulot{16}{Definition of a Euclidean space. The norm of an element of a Euclidean space}
\par\noindent\emph{Exercise set 1.}\nopagebreak
\begin{multicols}{2}\raggedright
\begin{enumerate}[label=\arabic*., leftmargin=*, itemsep=1pt]
\item $(x,y)=5$, $\|x\|=3$, $\|y\|=\sqrt{34}$
\item $(x,y)=2$, $\|x\|=\sqrt{31}$, $\|y\|=\sqrt{55}$
\item $(x,y)=-23$, $\|x\|=\sqrt{35}$, $\|y\|=\sqrt{46}$
\item $(x,y)=1$, $\|x\|=\sqrt{51}$, $\|y\|=\sqrt{91}$
\item $(x,y)=-10$, $\|x\|=7$, $\|y\|=\sqrt{6}$
\item $(x,y)=-13$, $\|x\|=3 \sqrt{2}$, $\|y\|=\sqrt{39}$
\item $(x,y)=-5$, $\|x\|=\sqrt{42}$, $\|y\|=\sqrt{67}$
\item $(x,y)=-7$, $\|x\|=\sqrt{17}$, $\|y\|=\sqrt{11}$
\item $(x,y)=-20$, $\|x\|=\sqrt{33}$, $\|y\|=\sqrt{37}$
\item $(x,y)=-26$, $\|x\|=7$, $\|y\|=\sqrt{61}$
\item $(x,y)=4$, $\|x\|=\sqrt{39}$, $\|y\|=\sqrt{59}$
\item $(x,y)=33$, $\|x\|=\sqrt{67}$, $\|y\|=\sqrt{38}$
\item $(x,y)=-10$, $\|x\|=\sqrt{70}$, $\|y\|=\sqrt{41}$
\item $(x,y)=15$, $\|x\|=\sqrt{46}$, $\|y\|=\sqrt{91}$
\item $(x,y)=32$, $\|x\|=2 \sqrt{15}$, $\|y\|=\sqrt{30}$
\item $(x,y)=49$, $\|x\|=\sqrt{42}$, $\|y\|=\sqrt{66}$
\item $(x,y)=-7$, $\|x\|=2 \sqrt{7}$, $\|y\|=\sqrt{13}$
\item $(x,y)=-23$, $\|x\|=2 \sqrt{19}$, $\|y\|=\sqrt{31}$
\item $(x,y)=-23$, $\|x\|=\sqrt{30}$, $\|y\|=\sqrt{58}$
\item $(x,y)=-8$, $\|x\|=\sqrt{22}$, $\|y\|=\sqrt{23}$
\item $(x,y)=4$, $\|x\|=2 \sqrt{5}$, $\|y\|=\sqrt{14}$
\item $(x,y)=-8$, $\|x\|=\sqrt{39}$, $\|y\|=\sqrt{55}$
\item $(x,y)=-10$, $\|x\|=\sqrt{22}$, $\|y\|=\sqrt{37}$
\item $(x,y)=32$, $\|x\|=\sqrt{42}$, $\|y\|=\sqrt{33}$
\item $(x,y)=-37$, $\|x\|=\sqrt{38}$, $\|y\|=\sqrt{91}$
\item $(x,y)=31$, $\|x\|=3 \sqrt{5}$, $\|y\|=\sqrt{38}$
\item $(x,y)=2$, $\|x\|=\sqrt{34}$, $\|y\|=\sqrt{19}$
\item $(x,y)=-6$, $\|x\|=4$, $\|y\|=\sqrt{79}$
\item $(x,y)=11$, $\|x\|=\sqrt{33}$, $\|y\|=\sqrt{55}$
\item $(x,y)=45$, $\|x\|=2 \sqrt{19}$, $\|y\|=\sqrt{47}$
\end{enumerate}
\end{multicols}
\par\noindent\emph{Exercise set 2.}\nopagebreak
\begin{multicols}{3}\raggedright
\begin{enumerate}[label=\arabic*., leftmargin=*, itemsep=1pt]
\item $\cos\varphi=\frac{3 \sqrt{705}}{470}$, $\rho(x,y)=\sqrt{89}$
\item $\cos\varphi=\frac{10 \sqrt{462}}{231}$, $\rho(x,y)=\sqrt{14}$
\item $\cos\varphi=\frac{4 \sqrt{1342}}{671}$, $\rho(x,y)=\sqrt{67}$
\item $\cos\varphi=- \frac{37 \sqrt{3009}}{3009}$, $\rho(x,y)=2 \sqrt{46}$
\item $\cos\varphi=- \frac{13 \sqrt{37}}{111}$, $\rho(x,y)=5 \sqrt{5}$
\item $\cos\varphi=\frac{\sqrt{295}}{177}$, $\rho(x,y)=\sqrt{94}$
\item $\cos\varphi=- \frac{2 \sqrt{286}}{143}$, $\rho(x,y)=\sqrt{129}$
\item $\cos\varphi=\frac{9 \sqrt{703}}{703}$, $\rho(x,y)=\sqrt{38}$
\item $\cos\varphi=- \frac{4}{15}$, $\rho(x,y)=\sqrt{58}$
\item $\cos\varphi=- \frac{17 \sqrt{1426}}{713}$, $\rho(x,y)=\sqrt{145}$
\item $\cos\varphi=\frac{\sqrt{2193}}{51}$, $\rho(x,y)=2 \sqrt{2}$
\item $\cos\varphi=\frac{5 \sqrt{14}}{21}$, $\rho(x,y)=\sqrt{10}$
\item $\cos\varphi=- \frac{\sqrt{1155}}{42}$, $\rho(x,y)=\sqrt{249}$
\item $\cos\varphi=\frac{2 \sqrt{87}}{29}$, $\rho(x,y)=2 \sqrt{10}$
\item $\cos\varphi=\frac{25 \sqrt{871}}{1742}$, $\rho(x,y)=\sqrt{69}$
\item $\cos\varphi=0$, $\rho(x,y)=\sqrt{78}$
\item $\cos\varphi=\frac{\sqrt{14}}{63}$, $\rho(x,y)=\sqrt{65}$
\item $\cos\varphi=- \frac{\sqrt{462}}{42}$, $\rho(x,y)=\sqrt{138}$
\item $\cos\varphi=\frac{9 \sqrt{35}}{70}$, $\rho(x,y)=\sqrt{19}$
\item $\cos\varphi=- \frac{17 \sqrt{122}}{366}$, $\rho(x,y)=\sqrt{113}$
\item $\cos\varphi=0$, $\rho(x,y)=2 \sqrt{17}$
\item $\cos\varphi=\frac{\sqrt{2}}{2}$, $\rho(x,y)=\sqrt{21}$
\item $\cos\varphi=\frac{\sqrt{273}}{78}$, $\rho(x,y)=3 \sqrt{6}$
\item $\cos\varphi=\frac{2 \sqrt{2}}{21}$, $\rho(x,y)=\sqrt{69}$
\item $\cos\varphi=- \frac{6}{17}$, $\rho(x,y)=\sqrt{138}$
\item $\cos\varphi=\frac{\sqrt{15}}{30}$, $\rho(x,y)=\sqrt{82}$
\item $\cos\varphi=\frac{25 \sqrt{1419}}{1419}$, $\rho(x,y)=\sqrt{26}$
\item $\cos\varphi=- \frac{7 \sqrt{2}}{30}$, $\rho(x,y)=\sqrt{57}$
\item $\cos\varphi=\frac{4 \sqrt{22}}{55}$, $\rho(x,y)=\sqrt{62}$
\item $\cos\varphi=\frac{11 \sqrt{133}}{399}$, $\rho(x,y)=2 \sqrt{15}$
\end{enumerate}
\end{multicols}
\par\noindent\emph{Exercise set 3.}\nopagebreak
\begin{multicols}{2}\raggedright
\begin{enumerate}[label=\arabic*., leftmargin=*, itemsep=1pt]
\item yes: $(x,y)=0$
\item $\frac15(1;-2;2;4)$
\item $|8|\le\sqrt{14}\sqrt{21}\approx17.15$
\item $\|x+y\|=\sqrt{290}\le 5+13=18$
\item $\lambda=-\frac{1}{2}$
\item $(x,y)=0$, $\|x\|=\sqrt{14}$, $\|y\|=\sqrt{21}$
\item $(f,g)=\frac12$, $\|f\|=\frac{1}{\sqrt3}$
\item $d=2$
\item $\cos(u,v)=\frac{7}{\sqrt{14}\sqrt6}=\frac{7}{2\sqrt{21}}$
\item $\|x+y\|^2=10=2+8$, since $(x,y)=0$
\end{enumerate}
\end{multicols}

% !Mode:: "TeX:UTF-8"
\javobmashgulot{17}{Constructing an orthonormal basis in a Euclidean space}
\par\noindent\emph{Exercise set 1.}\nopagebreak
\begin{multicols}{2}\raggedright
\begin{enumerate}[label=\arabic*., leftmargin=*, itemsep=1pt]
\item $\vec b_1=(1;\,0;\,3)$, $\vec b_2=(-\frac{3}{5};\,-1;\,\frac{1}{5})$
\item $\vec b_1=(-2;\,3;\,2)$, $\vec b_2=(\frac{40}{17};\,\frac{8}{17};\,\frac{28}{17})$
\item $\vec b_1=(-1;\,-3;\,-3)$, $\vec b_2=(-\frac{45}{19};\,-\frac{21}{19};\,\frac{36}{19})$
\item $\vec b_1=(2;\,0;\,-1)$, $\vec b_2=(\frac{9}{5};\,1;\,\frac{18}{5})$
\item $\vec b_1=(-1;\,2;\,3)$, $\vec b_2=(\frac{15}{7};\,\frac{12}{7};\,-\frac{3}{7})$
\item $\vec b_1=(-2;\,1;\,-3)$, $\vec b_2=(-\frac{16}{7};\,\frac{22}{7};\,\frac{18}{7})$
\item $\vec b_1=(-2;\,-2;\,2)$, $\vec b_2=(\frac{4}{3};\,\frac{4}{3};\,\frac{8}{3})$
\item $\vec b_1=(2;\,-2;\,-1)$, $\vec b_2=(-\frac{1}{3};\,-\frac{5}{3};\,\frac{8}{3})$
\item $\vec b_1=(2;\,-3;\,-1)$, $\vec b_2=(-\frac{8}{7};\,-\frac{9}{7};\,\frac{11}{7})$
\item $\vec b_1=(3;\,0;\,2)$, $\vec b_2=(-\frac{8}{13};\,-2;\,\frac{12}{13})$
\item $\vec b_1=(3;\,0;\,3)$, $\vec b_2=(2;\,-1;\,-2)$
\item $\vec b_1=(0;\,-1;\,1)$, $\vec b_2=(-3;\,\frac{5}{2};\,\frac{5}{2})$
\item $\vec b_1=(-1;\,-3;\,0)$, $\vec b_2=(-\frac{6}{5};\,\frac{2}{5};\,3)$
\item $\vec b_1=(-3;\,3;\,0)$, $\vec b_2=(-\frac{5}{2};\,-\frac{5}{2};\,1)$
\item $\vec b_1=(-2;\,-2;\,1)$, $\vec b_2=(\frac{5}{9};\,-\frac{13}{9};\,-\frac{16}{9})$
\item $\vec b_1=(-3;\,3;\,-1)$, $\vec b_2=(-\frac{5}{19};\,\frac{5}{19};\,\frac{30}{19})$
\item $\vec b_1=(3;\,1;\,1)$, $\vec b_2=(1;\,-2;\,-1)$
\item $\vec b_1=(2;\,-1;\,-1)$, $\vec b_2=(-\frac{4}{3};\,-\frac{11}{6};\,-\frac{5}{6})$
\item $\vec b_1=(-3;\,-3;\,3)$, $\vec b_2=(\frac{1}{3};\,-\frac{5}{3};\,-\frac{4}{3})$
\item $\vec b_1=(2;\,-1;\,0)$, $\vec b_2=(-\frac{4}{5};\,-\frac{8}{5};\,2)$
\item $\vec b_1=(-2;\,-1;\,-1)$, $\vec b_2=(-\frac{7}{3};\,\frac{7}{3};\,\frac{7}{3})$
\item $\vec b_1=(0;\,-3;\,1)$, $\vec b_2=(1;\,\frac{7}{10};\,\frac{21}{10})$
\item $\vec b_1=(-2;\,0;\,0)$, $\vec b_2=(0;\,1;\,3)$
\item $\vec b_1=(0;\,-2;\,-2)$, $\vec b_2=(-1;\,-2;\,2)$
\item $\vec b_1=(-3;\,-2;\,-3)$, $\vec b_2=(\frac{21}{11};\,\frac{3}{11};\,-\frac{23}{11})$
\item $\vec b_1=(3;\,-1;\,3)$, $\vec b_2=(-\frac{1}{19};\,-\frac{63}{19};\,-\frac{20}{19})$
\item $\vec b_1=(0;\,-3;\,2)$, $\vec b_2=(3;\,\frac{14}{13};\,\frac{21}{13})$
\item $\vec b_1=(-1;\,2;\,0)$, $\vec b_2=(-\frac{6}{5};\,-\frac{3}{5};\,1)$
\item $\vec b_1=(-3;\,1;\,2)$, $\vec b_2=(-\frac{5}{7};\,-\frac{17}{7};\,\frac{1}{7})$
\item $\vec b_1=(-1;\,1;\,-2)$, $\vec b_2=(\frac{10}{3};\,-\frac{4}{3};\,-\frac{7}{3})$
\end{enumerate}
\end{multicols}
\par\noindent\emph{Exercise set 2.}\nopagebreak
\begin{enumerate}[label=\arabic*., leftmargin=*, itemsep=1pt]
\item $\vec b_1=(-2;\,-1;\,-1)$, $\vec b_2=(0;\,\frac{1}{2};\,-\frac{1}{2})$, $\vec b_3=(-1;\,1;\,1)$
\item $\vec b_1=(-2;\,2;\,0)$, $\vec b_2=(\frac{3}{2};\,\frac{3}{2};\,2)$, $\vec b_3=(\frac{6}{17};\,\frac{6}{17};\,-\frac{9}{17})$
\item $\vec b_1=(2;\,0;\,-1)$, $\vec b_2=(-\frac{1}{5};\,0;\,-\frac{2}{5})$, $\vec b_3=(0;\,-1;\,0)$
\item $\vec b_1=(0;\,-2;\,2)$, $\vec b_2=(1;\,0;\,0)$, $\vec b_3=(0;\,-\frac{3}{2};\,-\frac{3}{2})$
\item $\vec b_1=(2;\,1;\,2)$, $\vec b_2=(\frac{4}{3};\,\frac{2}{3};\,-\frac{5}{3})$, $\vec b_3=(-\frac{3}{5};\,\frac{6}{5};\,0)$
\item $\vec b_1=(2;\,-2;\,-1)$, $\vec b_2=(\frac{1}{9};\,\frac{8}{9};\,-\frac{14}{9})$, $\vec b_3=(\frac{24}{29};\,\frac{18}{29};\,\frac{12}{29})$
\item $\vec b_1=(0;\,-2;\,2)$, $\vec b_2=(0;\,-1;\,-1)$, $\vec b_3=(-1;\,0;\,0)$
\item $\vec b_1=(-1;\,-2;\,2)$, $\vec b_2=(-\frac{22}{9};\,\frac{1}{9};\,-\frac{10}{9})$, $\vec b_3=(\frac{38}{65};\,-\frac{114}{65};\,-\frac{19}{13})$
\item $\vec b_1=(2;\,0;\,2)$, $\vec b_2=(-\frac{1}{2};\,0;\,\frac{1}{2})$, $\vec b_3=(0;\,2;\,0)$
\item $\vec b_1=(-2;\,-2;\,0)$, $\vec b_2=(2;\,-2;\,0)$, $\vec b_3=(0;\,0;\,-2)$
\item $\vec b_1=(0;\,-1;\,1)$, $\vec b_2=(-1;\,-\frac{1}{2};\,-\frac{1}{2})$, $\vec b_3=(\frac{2}{3};\,-\frac{2}{3};\,-\frac{2}{3})$
\item $\vec b_1=(-1;\,0;\,0)$, $\vec b_2=(0;\,-1;\,-1)$, $\vec b_3=(0;\,1;\,-1)$
\item $\vec b_1=(2;\,0;\,1)$, $\vec b_2=(-\frac{2}{5};\,-2;\,\frac{4}{5})$, $\vec b_3=(-\frac{2}{3};\,\frac{2}{3};\,\frac{4}{3})$
\item $\vec b_1=(2;\,2;\,2)$, $\vec b_2=(1;\,-2;\,1)$, $\vec b_3=(-\frac{1}{2};\,0;\,\frac{1}{2})$
\item $\vec b_1=(1;\,-2;\,-1)$, $\vec b_2=(\frac{1}{2};\,1;\,-\frac{3}{2})$, $\vec b_3=(\frac{4}{21};\,\frac{1}{21};\,\frac{2}{21})$
\item $\vec b_1=(1;\,2;\,-1)$, $\vec b_2=(\frac{4}{3};\,-\frac{1}{3};\,\frac{2}{3})$, $\vec b_3=(-\frac{1}{14};\,\frac{1}{7};\,\frac{3}{14})$
\item $\vec b_1=(-2;\,2;\,-1)$, $\vec b_2=(\frac{13}{9};\,\frac{5}{9};\,-\frac{16}{9})$, $\vec b_3=(-\frac{33}{50};\,-\frac{11}{10};\,-\frac{22}{25})$
\item $\vec b_1=(1;\,2;\,-2)$, $\vec b_2=(\frac{2}{9};\,-\frac{5}{9};\,-\frac{4}{9})$, $\vec b_3=(\frac{6}{5};\,0;\,\frac{3}{5})$
\item $\vec b_1=(-1;\,0;\,0)$, $\vec b_2=(0;\,1;\,-2)$, $\vec b_3=(0;\,\frac{4}{5};\,\frac{2}{5})$
\item $\vec b_1=(1;\,2;\,1)$, $\vec b_2=(1;\,0;\,-1)$, $\vec b_3=(\frac{2}{3};\,-\frac{2}{3};\,\frac{2}{3})$
\item $\vec b_1=(-1;\,0;\,1)$, $\vec b_2=(-\frac{3}{2};\,-1;\,-\frac{3}{2})$, $\vec b_3=(\frac{3}{11};\,-\frac{9}{11};\,\frac{3}{11})$
\item $\vec b_1=(0;\,-1;\,1)$, $\vec b_2=(1;\,-\frac{1}{2};\,-\frac{1}{2})$, $\vec b_3=(-\frac{2}{3};\,-\frac{2}{3};\,-\frac{2}{3})$
\item $\vec b_1=(1;\,1;\,1)$, $\vec b_2=(-\frac{1}{3};\,-\frac{1}{3};\,\frac{2}{3})$, $\vec b_3=(-\frac{1}{2};\,\frac{1}{2};\,0)$
\item $\vec b_1=(-1;\,-2;\,1)$, $\vec b_2=(\frac{1}{3};\,-\frac{4}{3};\,-\frac{7}{3})$, $\vec b_3=(-\frac{15}{11};\,\frac{5}{11};\,-\frac{5}{11})$
\item $\vec b_1=(0;\,2;\,2)$, $\vec b_2=(2;\,0;\,0)$, $\vec b_3=(0;\,-\frac{1}{2};\,\frac{1}{2})$
\item $\vec b_1=(-2;\,0;\,-2)$, $\vec b_2=(0;\,2;\,0)$, $\vec b_3=(1;\,0;\,-1)$
\item $\vec b_1=(0;\,1;\,-2)$, $\vec b_2=(2;\,-\frac{8}{5};\,-\frac{4}{5})$, $\vec b_3=(\frac{8}{9};\,\frac{8}{9};\,\frac{4}{9})$
\item $\vec b_1=(0;\,0;\,-1)$, $\vec b_2=(1;\,1;\,0)$, $\vec b_3=(\frac{1}{2};\,-\frac{1}{2};\,0)$
\item $\vec b_1=(-2;\,-1;\,2)$, $\vec b_2=(\frac{2}{3};\,\frac{4}{3};\,\frac{4}{3})$, $\vec b_3=(\frac{4}{3};\,-\frac{4}{3};\,\frac{2}{3})$
\item $\vec b_1=(-2;\,1;\,2)$, $\vec b_2=(\frac{7}{9};\,-\frac{8}{9};\,\frac{11}{9})$, $\vec b_3=(-\frac{18}{13};\,-\frac{24}{13};\,-\frac{6}{13})$
\end{enumerate}
\par\noindent\emph{Exercise set 3.}\nopagebreak
\begin{enumerate}[label=\arabic*., leftmargin=*, itemsep=1pt]
\item $\mathrm{pr}=(\frac{27}{11};\,-\frac{27}{11};\,-\frac{18}{11})$, $\vec h=(-\frac{27}{11};\,-\frac{17}{11};\,-\frac{15}{11})$
\item $\mathrm{pr}=(\frac{3}{2};\,1;\,-\frac{1}{2})$, $\vec h=(\frac{3}{2};\,-2;\,\frac{1}{2})$
\item $\mathrm{pr}=(2;\,0;\,2)$, $\vec h=(-2;\,1;\,2)$
\item $\mathrm{pr}=(-\frac{3}{11};\,-\frac{1}{11};\,\frac{1}{11})$, $\vec h=(\frac{25}{11};\,-\frac{43}{11};\,\frac{32}{11})$
\item $\mathrm{pr}=(\frac{5}{2};\,0;\,-\frac{5}{2})$, $\vec h=(-\frac{3}{2};\,3;\,-\frac{3}{2})$
\item $\pm\frac{1}{\sqrt{11}}(3;\,-1;\,1)$
\item $(\vec e_i,\vec e_j)=\delta_{ij}$, so the basis is orthonormal
\item $\vec x=3\vec e_1+3\vec e_2+0\cdot\vec e_3$
\item $Q^TQ=E$, $Q^{-1}=Q^T=Q$
\item $\frac{1}{\sqrt2}(1;0;1)$, $\frac{1}{\sqrt6}(-1;2;1)$
\end{enumerate}


\section*{Chapter 4. Linear Operators}
% !Mode:: "TeX:UTF-8"
\javobmashgulot{18}{Operations on linear operators. The space of linear operators. The matrix of a linear operator under a change of basis}
\par\noindent\emph{Exercise set 1.}\nopagebreak
\begin{enumerate}[label=\arabic*., leftmargin=*, itemsep=1pt]
\item $\left(\begin{smallmatrix}-5 & -9\\ -4 & 0\end{smallmatrix}\right)$; $\left(\begin{smallmatrix}8 & 4\\ -2 & -4\end{smallmatrix}\right)$; $\left(\begin{smallmatrix}-2 & 3\\ 4 & 6\end{smallmatrix}\right)$
\item $\left(\begin{smallmatrix}0 & -5\\ -6 & -4\end{smallmatrix}\right)$; $\left(\begin{smallmatrix}4 & 2\\ 2 & 1\end{smallmatrix}\right)$; $\left(\begin{smallmatrix}0 & 1\\ 0 & 5\end{smallmatrix}\right)$
\item $\left(\begin{smallmatrix}-6 & -4\\ 3 & -1\end{smallmatrix}\right)$; $\left(\begin{smallmatrix}-7 & -12\\ -1 & 9\end{smallmatrix}\right)$; $\left(\begin{smallmatrix}-6 & 3\\ 9 & 8\end{smallmatrix}\right)$
\item $\left(\begin{smallmatrix}-5 & 1\\ 1 & -4\end{smallmatrix}\right)$; $\left(\begin{smallmatrix}0 & 2\\ 4 & -6\end{smallmatrix}\right)$; $\left(\begin{smallmatrix}0 & 4\\ 2 & -6\end{smallmatrix}\right)$
\item $\left(\begin{smallmatrix}7 & -1\\ 5 & 2\end{smallmatrix}\right)$; $\left(\begin{smallmatrix}-9 & 3\\ -6 & 6\end{smallmatrix}\right)$; $\left(\begin{smallmatrix}0 & -6\\ -6 & -3\end{smallmatrix}\right)$
\item $\left(\begin{smallmatrix}3 & 0\\ 0 & 1\end{smallmatrix}\right)$; $\left(\begin{smallmatrix}-2 & 0\\ -1 & -3\end{smallmatrix}\right)$; $\left(\begin{smallmatrix}-2 & 0\\ 1 & -3\end{smallmatrix}\right)$
\item $\left(\begin{smallmatrix}-1 & 3\\ 3 & 1\end{smallmatrix}\right)$; $\left(\begin{smallmatrix}1 & 0\\ 7 & -4\end{smallmatrix}\right)$; $\left(\begin{smallmatrix}-2 & 2\\ 3 & -1\end{smallmatrix}\right)$
\item $\left(\begin{smallmatrix}0 & -7\\ 0 & 7\end{smallmatrix}\right)$; $\left(\begin{smallmatrix}-3 & 4\\ 0 & 2\end{smallmatrix}\right)$; $\left(\begin{smallmatrix}-3 & -3\\ 0 & 2\end{smallmatrix}\right)$
\item $\left(\begin{smallmatrix}1 & -4\\ -5 & -2\end{smallmatrix}\right)$; $\left(\begin{smallmatrix}0 & 6\\ -10 & -4\end{smallmatrix}\right)$; $\left(\begin{smallmatrix}2 & -8\\ 9 & -6\end{smallmatrix}\right)$
\item $\left(\begin{smallmatrix}-6 & 5\\ -4 & 3\end{smallmatrix}\right)$; $\left(\begin{smallmatrix}-2 & 1\\ -2 & 1\end{smallmatrix}\right)$; $\left(\begin{smallmatrix}0 & -1\\ 0 & -1\end{smallmatrix}\right)$
\item $\left(\begin{smallmatrix}6 & 2\\ 5 & 4\end{smallmatrix}\right)$; $\left(\begin{smallmatrix}7 & 8\\ 15 & 12\end{smallmatrix}\right)$; $\left(\begin{smallmatrix}18 & 6\\ 9 & 1\end{smallmatrix}\right)$
\item $\left(\begin{smallmatrix}-9 & 0\\ -11 & -6\end{smallmatrix}\right)$; $\left(\begin{smallmatrix}6 & 0\\ 4 & 0\end{smallmatrix}\right)$; $\left(\begin{smallmatrix}6 & 0\\ 13 & 0\end{smallmatrix}\right)$
\item $\left(\begin{smallmatrix}3 & 5\\ -1 & 4\end{smallmatrix}\right)$; $\left(\begin{smallmatrix}-9 & -3\\ -3 & -8\end{smallmatrix}\right)$; $\left(\begin{smallmatrix}-11 & -1\\ -3 & -6\end{smallmatrix}\right)$
\item $\left(\begin{smallmatrix}6 & -8\\ -4 & 12\end{smallmatrix}\right)$; $\left(\begin{smallmatrix}1 & -6\\ -4 & 3\end{smallmatrix}\right)$; $\left(\begin{smallmatrix}-3 & -8\\ 0 & 7\end{smallmatrix}\right)$
\item $\left(\begin{smallmatrix}6 & -3\\ -2 & 0\end{smallmatrix}\right)$; $\left(\begin{smallmatrix}15 & 0\\ 0 & 0\end{smallmatrix}\right)$; $\left(\begin{smallmatrix}9 & -9\\ -6 & 6\end{smallmatrix}\right)$
\item $\left(\begin{smallmatrix}0 & -5 & -4\\ -2 & 4 & 6\\ -1 & -6 & -3\end{smallmatrix}\right)$; $\left(\begin{smallmatrix}1 & 1 & -2\\ 2 & 2 & -2\\ -1 & -2 & 0\end{smallmatrix}\right)$; $\left(\begin{smallmatrix}-1 & -1 & -1\\ 0 & 1 & 2\\ -1 & 2 & 3\end{smallmatrix}\right)$
\item $\left(\begin{smallmatrix}1 & 6 & -4\\ 1 & -4 & -5\\ 2 & 1 & 8\end{smallmatrix}\right)$; $\left(\begin{smallmatrix}0 & -5 & 7\\ -2 & -10 & -6\\ 3 & -1 & 5\end{smallmatrix}\right)$; $\left(\begin{smallmatrix}-7 & 9 & -1\\ 10 & 7 & 5\\ -3 & -11 & -5\end{smallmatrix}\right)$
\item $\left(\begin{smallmatrix}2 & 1 & 0\\ -1 & -3 & -1\\ -3 & 2 & 3\end{smallmatrix}\right)$; $\left(\begin{smallmatrix}12 & -3 & 6\\ 5 & -4 & -7\\ 3 & -2 & -2\end{smallmatrix}\right)$; $\left(\begin{smallmatrix}3 & 1 & 6\\ -3 & 0 & 9\\ 11 & 3 & 3\end{smallmatrix}\right)$
\item $\left(\begin{smallmatrix}9 & 9 & -8\\ -3 & 8 & 4\\ 3 & -4 & 3\end{smallmatrix}\right)$; $\left(\begin{smallmatrix}2 & 22 & -12\\ -1 & -1 & 11\\ -9 & 3 & -6\end{smallmatrix}\right)$; $\left(\begin{smallmatrix}9 & 15 & 9\\ -7 & 3 & 5\\ 5 & 2 & -17\end{smallmatrix}\right)$
\item $\left(\begin{smallmatrix}2 & 2 & 1\\ 3 & -4 & 5\\ 7 & 4 & 5\end{smallmatrix}\right)$; $\left(\begin{smallmatrix}13 & 3 & 5\\ 3 & 1 & 2\\ -17 & -3 & -6\end{smallmatrix}\right)$; $\left(\begin{smallmatrix}9 & 6 & -7\\ 2 & 2 & 3\\ 8 & 6 & -3\end{smallmatrix}\right)$
\item $\left(\begin{smallmatrix}4 & -1 & 5\\ 5 & 4 & -9\\ -1 & 0 & 12\end{smallmatrix}\right)$; $\left(\begin{smallmatrix}1 & -2 & 2\\ 0 & -4 & 10\\ -1 & 0 & 13\end{smallmatrix}\right)$; $\left(\begin{smallmatrix}5 & -1 & 5\\ -1 & -5 & -10\\ 5 & 1 & 10\end{smallmatrix}\right)$
\item $\left(\begin{smallmatrix}-6 & -5 & 5\\ 6 & 5 & -3\\ -11 & -10 & 7\end{smallmatrix}\right)$; $\left(\begin{smallmatrix}8 & 11 & -9\\ 7 & 7 & -9\\ 6 & 8 & -10\end{smallmatrix}\right)$; $\left(\begin{smallmatrix}-1 & -1 & 3\\ -6 & 0 & -8\\ 3 & -3 & 6\end{smallmatrix}\right)$
\item $\left(\begin{smallmatrix}-7 & 9 & -6\\ 10 & 3 & 9\\ -2 & -5 & 6\end{smallmatrix}\right)$; $\left(\begin{smallmatrix}8 & -8 & 21\\ 7 & 2 & 6\\ 4 & 8 & -9\end{smallmatrix}\right)$; $\left(\begin{smallmatrix}1 & -3 & 3\\ 0 & 12 & 0\\ 2 & -6 & -12\end{smallmatrix}\right)$
\item $\left(\begin{smallmatrix}-5 & -8 & -1\\ -3 & 1 & -2\\ 2 & 1 & -3\end{smallmatrix}\right)$; $\left(\begin{smallmatrix}-8 & -9 & -7\\ 4 & 9 & 8\\ 0 & 0 & 11\end{smallmatrix}\right)$; $\left(\begin{smallmatrix}6 & 1 & 5\\ -9 & -4 & 3\\ 3 & 5 & 10\end{smallmatrix}\right)$
\item $\left(\begin{smallmatrix}-1 & -1 & 4\\ -4 & 5 & -3\\ 0 & -2 & 2\end{smallmatrix}\right)$; $\left(\begin{smallmatrix}-3 & -6 & -5\\ 6 & -10 & 8\\ -12 & 0 & 0\end{smallmatrix}\right)$; $\left(\begin{smallmatrix}-1 & 13 & 0\\ 5 & -5 & 12\\ 5 & -2 & -7\end{smallmatrix}\right)$
\item $\left(\begin{smallmatrix}-1 & -2 & 3\\ 5 & -2 & -3\\ -4 & -4 & 1\end{smallmatrix}\right)$; $\left(\begin{smallmatrix}12 & 2 & 1\\ 1 & -7 & -5\\ 6 & 2 & 4\end{smallmatrix}\right)$; $\left(\begin{smallmatrix}4 & -2 & 2\\ -1 & 5 & -9\\ 4 & -8 & 0\end{smallmatrix}\right)$
\item $\left(\begin{smallmatrix}-4 & 1 & -1\\ 2 & 4 & -3\\ 2 & -4 & 1\end{smallmatrix}\right)$; $\left(\begin{smallmatrix}1 & 5 & -3\\ 0 & 9 & -3\\ -5 & -2 & 2\end{smallmatrix}\right)$; $\left(\begin{smallmatrix}4 & 0 & 4\\ -3 & 4 & -10\\ 0 & -4 & 4\end{smallmatrix}\right)$
\item $\left(\begin{smallmatrix}-3 & 0 & 1\\ -2 & 4 & 4\\ -1 & -1 & -1\end{smallmatrix}\right)$; $\left(\begin{smallmatrix}-3 & -11 & -4\\ -2 & -8 & -4\\ -12 & -7 & 9\end{smallmatrix}\right)$; $\left(\begin{smallmatrix}-8 & 10 & 9\\ 10 & 5 & -6\\ -6 & -14 & 1\end{smallmatrix}\right)$
\item $\left(\begin{smallmatrix}-1 & 2 & -5\\ 3 & -2 & 5\\ 0 & 0 & 0\end{smallmatrix}\right)$; $\left(\begin{smallmatrix}-8 & -13 & 4\\ 9 & 12 & -4\\ -2 & 5 & -2\end{smallmatrix}\right)$; $\left(\begin{smallmatrix}8 & 8 & 3\\ -13 & -9 & -9\\ 5 & 3 & 3\end{smallmatrix}\right)$
\item $\left(\begin{smallmatrix}-6 & -6 & 3\\ 6 & 3 & 0\\ 3 & 4 & 3\end{smallmatrix}\right)$; $\left(\begin{smallmatrix}4 & 6 & 8\\ 1 & 2 & 5\\ 0 & 0 & -2\end{smallmatrix}\right)$; $\left(\begin{smallmatrix}4 & 6 & -10\\ -5 & -8 & 14\\ -6 & -6 & 8\end{smallmatrix}\right)$
\end{enumerate}
\par\noindent\emph{Exercise set 2.}\nopagebreak
\begin{multicols}{2}\raggedright
\begin{enumerate}[label=\arabic*., leftmargin=*, itemsep=1pt]
\item $A^{\prime}=\left(\begin{smallmatrix}-1 & 9\\ 0 & 4\end{smallmatrix}\right)$
\item $A^{\prime}=\left(\begin{smallmatrix}1 & 1\\ -2 & -3\end{smallmatrix}\right)$
\item $A^{\prime}=\left(\begin{smallmatrix}-3 & -2\\ 2 & 3\end{smallmatrix}\right)$
\item $A^{\prime}=\left(\begin{smallmatrix}-3 & 4\\ -6 & 4\end{smallmatrix}\right)$
\item $A^{\prime}=\left(\begin{smallmatrix}0 & -1\\ 3 & -2\end{smallmatrix}\right)$
\item $A^{\prime}=\left(\begin{smallmatrix}5 & -4\\ 3 & -3\end{smallmatrix}\right)$
\item $A^{\prime}=\left(\begin{smallmatrix}6 & 3\\ -5 & 1\end{smallmatrix}\right)$
\item $A^{\prime}=\left(\begin{smallmatrix}12 & -10\\ 22 & -18\end{smallmatrix}\right)$
\item $A^{\prime}=\left(\begin{smallmatrix}2 & -4\\ 0 & -3\end{smallmatrix}\right)$
\item $A^{\prime}=\left(\begin{smallmatrix}-2 & -4\\ -2 & 2\end{smallmatrix}\right)$
\item $A^{\prime}=\left(\begin{smallmatrix}1 & -9\\ 1 & 6\end{smallmatrix}\right)$
\item $A^{\prime}=\left(\begin{smallmatrix}-8 & 4\\ -6 & 5\end{smallmatrix}\right)$
\item $A^{\prime}=\left(\begin{smallmatrix}-6 & 29\\ -2 & 9\end{smallmatrix}\right)$
\item $A^{\prime}=\left(\begin{smallmatrix}-9 & -7\\ 11 & 9\end{smallmatrix}\right)$
\item $A^{\prime}=\left(\begin{smallmatrix}-5 & 2\\ -6 & 6\end{smallmatrix}\right)$
\item $A^{\prime}=\left(\begin{smallmatrix}-4 & 2\\ -6 & 4\end{smallmatrix}\right)$
\item $A^{\prime}=\left(\begin{smallmatrix}2 & -2\\ 2 & 4\end{smallmatrix}\right)$
\item $A^{\prime}=\left(\begin{smallmatrix}1 & 0\\ -2 & 0\end{smallmatrix}\right)$
\item $A^{\prime}=\left(\begin{smallmatrix}-11 & 8\\ -8 & 7\end{smallmatrix}\right)$
\item $A^{\prime}=\left(\begin{smallmatrix}7 & -4\\ 4 & 0\end{smallmatrix}\right)$
\item $A^{\prime}=\left(\begin{smallmatrix}-3 & -7\\ 2 & 2\end{smallmatrix}\right)$
\item $A^{\prime}=\left(\begin{smallmatrix}-13 & 7\\ -12 & 6\end{smallmatrix}\right)$
\item $A^{\prime}=\left(\begin{smallmatrix}2 & -12\\ 0 & -2\end{smallmatrix}\right)$
\item $A^{\prime}=\left(\begin{smallmatrix}10 & 15\\ -3 & -5\end{smallmatrix}\right)$
\item $A^{\prime}=\left(\begin{smallmatrix}4 & -2\\ 14 & -5\end{smallmatrix}\right)$
\item $A^{\prime}=\left(\begin{smallmatrix}-3 & 1\\ 5 & -1\end{smallmatrix}\right)$
\item $A^{\prime}=\left(\begin{smallmatrix}3 & 4\\ 0 & -2\end{smallmatrix}\right)$
\item $A^{\prime}=\left(\begin{smallmatrix}-4 & -3\\ -3 & -4\end{smallmatrix}\right)$
\item $A^{\prime}=\left(\begin{smallmatrix}-21 & -10\\ 34 & 16\end{smallmatrix}\right)$
\item $A^{\prime}=\left(\begin{smallmatrix}1 & 6\\ 1 & 3\end{smallmatrix}\right)$
\end{enumerate}
\end{multicols}
\par\noindent\emph{Exercise set 3.}\nopagebreak
\begin{multicols}{2}\raggedright
\begin{enumerate}[label=\arabic*., leftmargin=*, itemsep=1pt]
\item yes; $\left(\begin{smallmatrix}1&2\\3&-1\end{smallmatrix}\right)$
\item no: $\mathcal A(\vec 0)\ne\vec 0$
\item $\left(\begin{smallmatrix}0&-1\\1&0\end{smallmatrix}\right)$; $(-1;3)$
\item $\mathrm{diag}(1,1,0)$
\item $\left(\begin{smallmatrix}1&0\\0&-1\end{smallmatrix}\right)$
\item $(-1;-3;8)$
\item $\ker=\{t(-2;1)\}$, $\mathrm{Im}=\{t(1;2)\}$
\item $\left(\begin{smallmatrix}0&-1\\-1&0\end{smallmatrix}\right)$
\item $\left(\begin{smallmatrix}3&-1\\-5&2\end{smallmatrix}\right)$
\item $\det K=11$, $\gcd(11,26)=1$, hence invertible
\end{enumerate}
\end{multicols}

% !Mode:: "TeX:UTF-8"
\javobmashgulot{19}{Eigenvalues and eigenvectors of a linear operator}
\par\noindent\emph{Exercise set 1.}\nopagebreak
\begin{multicols}{2}\raggedright
\begin{enumerate}[label=\arabic*., leftmargin=*, itemsep=1pt]
\item $\lambda=4$: $(1;\,0)$; $\lambda=5$: $(1;\,1)$
\item $\lambda=0$: $(1;\,0)$; $\lambda=2$: $(1;\,-1)$
\item $\lambda=-4$: $(1;\,0)$; $\lambda=0$: $(1;\,-1)$
\item $\lambda=-3$: $(1;\,-1)$; $\lambda=3$: $(1;\,-2)$
\item $\lambda=0$: $(1;\,-2)$; $\lambda=5$: $(1;\,-1)$
\item $\lambda=0$: $(1;\,0)$; $\lambda=4$: $(2;\,1)$
\item $\lambda=1$: $(1;\,-1)$; $\lambda=3$: $(2;\,-1)$
\item $\lambda=-2$: $(0;\,1)$; $\lambda=4$: $(1;\,1)$
\item $\lambda=2$: $(1;\,-2)$; $\lambda=4$: $(0;\,1)$
\item $\lambda=-3$: $(2;\,-1)$; $\lambda=-2$: $(1;\,0)$
\item $\lambda=0$: $(0;\,1)$; $\lambda=1$: $(1;\,2)$
\item $\lambda=3$: $(1;\,-1)$; $\lambda=4$: $(2;\,-1)$
\item $\lambda=-2$: $(0;\,1)$; $\lambda=1$: $(1;\,-1)$
\item $\lambda=-2$: $(2;\,1)$; $\lambda=0$: $(1;\,0)$
\item $\lambda=2$: $(0;\,1)$; $\lambda=5$: $(1;\,1)$
\item $\lambda=-4$: $(1;\,0)$; $\lambda=-2$: $(2;\,-1)$
\item $\lambda=-1$: $(1;\,0)$; $\lambda=5$: $(1;\,-1)$
\item $\lambda=0$: $(2;\,-1)$; $\lambda=3$: $(1;\,0)$
\item $\lambda=-2$: $(0;\,1)$; $\lambda=5$: $(1;\,-1)$
\item $\lambda=-1$: $(1;\,-2)$; $\lambda=5$: $(0;\,1)$
\item $\lambda=1$: $(1;\,2)$; $\lambda=3$: $(0;\,1)$
\item $\lambda=1$: $(2;\,1)$; $\lambda=2$: $(1;\,0)$
\item $\lambda=-3$: $(1;\,-2)$; $\lambda=0$: $(1;\,-1)$
\item $\lambda=1$: $(2;\,-1)$; $\lambda=5$: $(1;\,0)$
\item $\lambda=-3$: $(1;\,0)$; $\lambda=1$: $(2;\,1)$
\item $\lambda=-2$: $(1;\,0)$; $\lambda=-1$: $(1;\,1)$
\item $\lambda=-1$: $(1;\,1)$; $\lambda=1$: $(2;\,1)$
\item $\lambda=-1$: $(1;\,0)$; $\lambda=1$: $(1;\,1)$
\item $\lambda=0$: $(0;\,1)$; $\lambda=5$: $(1;\,2)$
\item $\lambda=-4$: $(1;\,0)$; $\lambda=0$: $(1;\,-1)$
\end{enumerate}
\end{multicols}
\par\noindent\emph{Exercise set 2.}\nopagebreak
\begin{enumerate}[label=\arabic*., leftmargin=*, itemsep=1pt]
\item $\lambda=-2$: $(1;\,-1;\,-1)$; $\lambda=-1$: $(0;\,1;\,0)$; $\lambda=4$: $(1;\,2;\,-2)$
\item $\lambda=-2$: $(1;\,0;\,0)$; $\lambda=2$: $(0;\,2;\,-1)$; $\lambda=3$: $(1;\,1;\,-1)$
\item $\lambda=-2$: $(0;\,1;\,-1)$; $\lambda=3$: $(1;\,0;\,-1)$; $\lambda=4$: $(0;\,0;\,1)$
\item $\lambda=-2$: $(0;\,1;\,0)$; $\lambda=0$: $(1;\,-2;\,-1)$; $\lambda=4$: $(0;\,1;\,1)$
\item $\lambda=-2$: $(0;\,1;\,-2)$; $\lambda=-1$: $(0;\,0;\,1)$; $\lambda=1$: $(1;\,2;\,-1)$
\item $\lambda=-3$: $(0;\,1;\,1)$; $\lambda=-1$: $(1;\,-1;\,0)$; $\lambda=1$: $(1;\,-1;\,1)$
\item $\lambda=-2$: $(2;\,-1;\,2)$; $\lambda=-1$: $(1;\,-2;\,2)$; $\lambda=0$: $(0;\,1;\,-1)$
\item $\lambda=-2$: $(1;\,0;\,2)$; $\lambda=-1$: $(0;\,1;\,0)$; $\lambda=3$: $(1;\,0;\,1)$
\item $\lambda=-2$: $(1;\,0;\,0)$; $\lambda=1$: $(0;\,1;\,-1)$; $\lambda=4$: $(2;\,-1;\,0)$
\item $\lambda=-3$: $(2;\,1;\,1)$; $\lambda=-2$: $(1;\,0;\,0)$; $\lambda=3$: $(2;\,1;\,0)$
\item $\lambda=0$: $(1;\,0;\,0)$; $\lambda=1$: $(1;\,0;\,1)$; $\lambda=4$: $(2;\,-1;\,-1)$
\item $\lambda=1$: $(1;\,-2;\,0)$; $\lambda=2$: $(0;\,1;\,-1)$; $\lambda=3$: $(1;\,-2;\,1)$
\item $\lambda=-3$: $(1;\,0;\,0)$; $\lambda=-2$: $(2;\,2;\,1)$; $\lambda=2$: $(1;\,1;\,0)$
\item $\lambda=-3$: $(2;\,1;\,-2)$; $\lambda=2$: $(0;\,0;\,1)$; $\lambda=3$: $(1;\,1;\,2)$
\item $\lambda=-3$: $(1;\,1;\,-1)$; $\lambda=-2$: $(1;\,2;\,-1)$; $\lambda=3$: $(1;\,1;\,0)$
\item $\lambda=0$: $(1;\,-1;\,-2)$; $\lambda=1$: $(0;\,1;\,0)$; $\lambda=4$: $(1;\,1;\,-1)$
\item $\lambda=-2$: $(1;\,1;\,1)$; $\lambda=-1$: $(2;\,0;\,-1)$; $\lambda=0$: $(2;\,1;\,0)$
\item $\lambda=-2$: $(1;\,-1;\,1)$; $\lambda=2$: $(0;\,1;\,-2)$; $\lambda=4$: $(0;\,0;\,1)$
\item $\lambda=-3$: $(1;\,-2;\,2)$; $\lambda=-2$: $(0;\,1;\,-1)$; $\lambda=2$: $(1;\,1;\,0)$
\item $\lambda=-3$: $(0;\,1;\,-1)$; $\lambda=0$: $(1;\,0;\,1)$; $\lambda=4$: $(2;\,1;\,0)$
\item $\lambda=-2$: $(1;\,0;\,0)$; $\lambda=2$: $(1;\,-1;\,-1)$; $\lambda=4$: $(0;\,1;\,2)$
\item $\lambda=1$: $(1;\,-2;\,2)$; $\lambda=2$: $(1;\,-1;\,1)$; $\lambda=3$: $(1;\,-1;\,0)$
\item $\lambda=-2$: $(1;\,1;\,-2)$; $\lambda=-1$: $(1;\,0;\,-1)$; $\lambda=4$: $(1;\,0;\,-2)$
\item $\lambda=1$: $(1;\,0;\,-1)$; $\lambda=3$: $(0;\,0;\,1)$; $\lambda=4$: $(2;\,-1;\,-1)$
\item $\lambda=2$: $(1;\,0;\,0)$; $\lambda=3$: $(1;\,-2;\,-1)$; $\lambda=4$: $(2;\,1;\,1)$
\item $\lambda=1$: $(0;\,1;\,1)$; $\lambda=2$: $(1;\,0;\,1)$; $\lambda=4$: $(1;\,-1;\,1)$
\item $\lambda=-3$: $(1;\,1;\,-1)$; $\lambda=-2$: $(1;\,0;\,0)$; $\lambda=2$: $(1;\,-1;\,2)$
\item $\lambda=-3$: $(1;\,0;\,-1)$; $\lambda=0$: $(1;\,1;\,-1)$; $\lambda=4$: $(1;\,-1;\,0)$
\item $\lambda=-1$: $(2;\,-1;\,0)$; $\lambda=0$: $(1;\,2;\,2)$; $\lambda=1$: $(1;\,1;\,1)$
\item $\lambda=-3$: $(1;\,0;\,0)$; $\lambda=-2$: $(1;\,0;\,1)$; $\lambda=1$: $(0;\,1;\,-2)$
\end{enumerate}
\par\noindent\emph{Exercise set 3.}\nopagebreak
\begin{multicols}{2}\raggedright
\begin{enumerate}[label=\arabic*., leftmargin=*, itemsep=1pt]
\item $- \lambda^{3} + \lambda^{2} + 7 \lambda + 1$
\item $- \lambda^{3} - 4 \lambda^{2} + 4 \lambda + 13$
\item $- \lambda^{3} + 2 \lambda^{2} + 17 \lambda - 33$
\item $- \lambda^{3} + 4 \lambda^{2} - 4 \lambda - 9$
\item $- \lambda^{3} - \lambda^{2} + 8 \lambda$
\item $- \lambda^{3} - \lambda^{2} + 3 \lambda - 1$
\item $A^2$: $1;4;9$; \ $A^{-1}$: $1;\frac12;-\frac13$; \ $A+2E$: $3;4;-1$
\item $2;-3;1$
\item $x=\frac13(2;1)$
\item $\lambda=1$: $(1;-1)$, $\lambda=3$: $(1;1)$; $1\cdot1+(-1)\cdot1=0$
\end{enumerate}
\end{multicols}

% !Mode:: "TeX:UTF-8"
\javobmashgulot{20}{Linear operators with a basis of eigenvectors}
\par\noindent\emph{Exercise set 1.}\nopagebreak
\begin{enumerate}[label=\arabic*., leftmargin=*, itemsep=1pt]
\item $T=\left(\begin{smallmatrix}1 & 0\\ -2 & 1\end{smallmatrix}\right)$, $D=\left(\begin{smallmatrix}-2 & 0\\ 0 & 0\end{smallmatrix}\right)$
\item $T=\left(\begin{smallmatrix}2 & 1\\ -1 & -1\end{smallmatrix}\right)$, $D=\left(\begin{smallmatrix}-2 & 0\\ 0 & 0\end{smallmatrix}\right)$
\item $T=\left(\begin{smallmatrix}1 & 2\\ 1 & 1\end{smallmatrix}\right)$, $D=\left(\begin{smallmatrix}-3 & 0\\ 0 & 2\end{smallmatrix}\right)$
\item $T=\left(\begin{smallmatrix}0 & 1\\ 1 & 1\end{smallmatrix}\right)$, $D=\left(\begin{smallmatrix}-3 & 0\\ 0 & -2\end{smallmatrix}\right)$
\item $T=\left(\begin{smallmatrix}2 & 1\\ -1 & -1\end{smallmatrix}\right)$, $D=\left(\begin{smallmatrix}-3 & 0\\ 0 & 0\end{smallmatrix}\right)$
\item $T=\left(\begin{smallmatrix}0 & 1\\ 1 & 0\end{smallmatrix}\right)$, $D=\left(\begin{smallmatrix}-3 & 0\\ 0 & 3\end{smallmatrix}\right)$
\item $T=\left(\begin{smallmatrix}0 & 1\\ 1 & 0\end{smallmatrix}\right)$, $D=\left(\begin{smallmatrix}1 & 0\\ 0 & 3\end{smallmatrix}\right)$
\item $T=\left(\begin{smallmatrix}1 & 0\\ 2 & 1\end{smallmatrix}\right)$, $D=\left(\begin{smallmatrix}-1 & 0\\ 0 & 3\end{smallmatrix}\right)$
\item $T=\left(\begin{smallmatrix}0 & 1\\ 1 & 0\end{smallmatrix}\right)$, $D=\left(\begin{smallmatrix}0 & 0\\ 0 & 4\end{smallmatrix}\right)$
\item $T=\left(\begin{smallmatrix}0 & 1\\ 1 & 2\end{smallmatrix}\right)$, $D=\left(\begin{smallmatrix}-3 & 0\\ 0 & 1\end{smallmatrix}\right)$
\item $T=\left(\begin{smallmatrix}1 & 1\\ 1 & 0\end{smallmatrix}\right)$, $D=\left(\begin{smallmatrix}-3 & 0\\ 0 & 0\end{smallmatrix}\right)$
\item $T=\left(\begin{smallmatrix}1 & 0\\ 0 & 1\end{smallmatrix}\right)$, $D=\left(\begin{smallmatrix}0 & 0\\ 0 & 4\end{smallmatrix}\right)$
\item $T=\left(\begin{smallmatrix}1 & 1\\ -1 & 0\end{smallmatrix}\right)$, $D=\left(\begin{smallmatrix}-3 & 0\\ 0 & -1\end{smallmatrix}\right)$
\item $T=\left(\begin{smallmatrix}2 & 1\\ -1 & 0\end{smallmatrix}\right)$, $D=\left(\begin{smallmatrix}-2 & 0\\ 0 & 4\end{smallmatrix}\right)$
\item $T=\left(\begin{smallmatrix}1 & 0\\ 1 & 1\end{smallmatrix}\right)$, $D=\left(\begin{smallmatrix}0 & 0\\ 0 & 4\end{smallmatrix}\right)$
\item $T=\left(\begin{smallmatrix}2 & 1\\ -1 & -1\end{smallmatrix}\right)$, $D=\left(\begin{smallmatrix}-1 & 0\\ 0 & 0\end{smallmatrix}\right)$
\item $T=\left(\begin{smallmatrix}0 & 1\\ 1 & 0\end{smallmatrix}\right)$, $D=\left(\begin{smallmatrix}-1 & 0\\ 0 & 4\end{smallmatrix}\right)$
\item $T=\left(\begin{smallmatrix}0 & 1\\ 1 & 0\end{smallmatrix}\right)$, $D=\left(\begin{smallmatrix}-3 & 0\\ 0 & -2\end{smallmatrix}\right)$
\item $T=\left(\begin{smallmatrix}2 & 1\\ 1 & 1\end{smallmatrix}\right)$, $D=\left(\begin{smallmatrix}-4 & 0\\ 0 & -3\end{smallmatrix}\right)$
\item $T=\left(\begin{smallmatrix}0 & 1\\ 1 & 1\end{smallmatrix}\right)$, $D=\left(\begin{smallmatrix}-3 & 0\\ 0 & 4\end{smallmatrix}\right)$
\item $T=\left(\begin{smallmatrix}1 & 2\\ -1 & -1\end{smallmatrix}\right)$, $D=\left(\begin{smallmatrix}-4 & 0\\ 0 & 0\end{smallmatrix}\right)$
\item $T=\left(\begin{smallmatrix}1 & 0\\ -2 & 1\end{smallmatrix}\right)$, $D=\left(\begin{smallmatrix}0 & 0\\ 0 & 4\end{smallmatrix}\right)$
\item $T=\left(\begin{smallmatrix}1 & 0\\ -1 & 1\end{smallmatrix}\right)$, $D=\left(\begin{smallmatrix}-1 & 0\\ 0 & 3\end{smallmatrix}\right)$
\item $T=\left(\begin{smallmatrix}1 & 2\\ 0 & 1\end{smallmatrix}\right)$, $D=\left(\begin{smallmatrix}2 & 0\\ 0 & 4\end{smallmatrix}\right)$
\item $T=\left(\begin{smallmatrix}1 & 1\\ 0 & -1\end{smallmatrix}\right)$, $D=\left(\begin{smallmatrix}1 & 0\\ 0 & 2\end{smallmatrix}\right)$
\item $T=\left(\begin{smallmatrix}2 & 1\\ 1 & 1\end{smallmatrix}\right)$, $D=\left(\begin{smallmatrix}-3 & 0\\ 0 & -2\end{smallmatrix}\right)$
\item $T=\left(\begin{smallmatrix}1 & 2\\ 0 & -1\end{smallmatrix}\right)$, $D=\left(\begin{smallmatrix}0 & 0\\ 0 & 4\end{smallmatrix}\right)$
\item $T=\left(\begin{smallmatrix}1 & 1\\ 1 & 2\end{smallmatrix}\right)$, $D=\left(\begin{smallmatrix}-3 & 0\\ 0 & 2\end{smallmatrix}\right)$
\item $T=\left(\begin{smallmatrix}0 & 1\\ 1 & 0\end{smallmatrix}\right)$, $D=\left(\begin{smallmatrix}-3 & 0\\ 0 & -1\end{smallmatrix}\right)$
\item $T=\left(\begin{smallmatrix}1 & 0\\ 0 & 1\end{smallmatrix}\right)$, $D=\left(\begin{smallmatrix}-4 & 0\\ 0 & -1\end{smallmatrix}\right)$
\end{enumerate}
\par\noindent\emph{Exercise set 2.}\nopagebreak
\begin{multicols}{2}\raggedright
\begin{enumerate}[label=\arabic*., leftmargin=*, itemsep=1pt]
\item $\lambda$: $-2$, $1$, $2$; diagonalizable
\item $\lambda$: $-3$, $1$, $4$; diagonalizable
\item $\lambda$: $0$ (double), $2$; not diagonalizable
\item $\lambda$: $-3$, $0$, $1$; diagonalizable
\item $\lambda$: $2$ (double), $3$; diagonalizable
\item $\lambda$: $-3$, $3$ (double); not diagonalizable
\item $\lambda$: $-3$, $-1$, $4$; diagonalizable
\item $\lambda$: $-1$, $4$ (double); diagonalizable
\item $\lambda$: $-2$ (double), $-1$; not diagonalizable
\item $\lambda$: $0$, $3$ (double); diagonalizable
\item $\lambda$: $-2$, $0$, $2$; diagonalizable
\item $\lambda$: $-3$ (double), $-2$; not diagonalizable
\item $\lambda$: $-3$, $-2$, $2$; diagonalizable
\item $\lambda$: $-2$, $-1$, $3$; diagonalizable
\item $\lambda$: $0$ (double), $1$; not diagonalizable
\item $\lambda$: $-1$ (double), $4$; diagonalizable
\item $\lambda$: $-1$ (double), $0$; diagonalizable
\item $\lambda$: $-3$, $-2$ (double); not diagonalizable
\item $\lambda$: $-2$, $-1$, $0$; diagonalizable
\item $\lambda$: $-3$ (double), $-1$; diagonalizable
\item $\lambda$: $-2$, $2$ (double); not diagonalizable
\item $\lambda$: $-3$, $-2$, $3$; diagonalizable
\item $\lambda$: $-3$, $0$, $3$; diagonalizable
\item $\lambda$: $-3$ (double), $-2$; not diagonalizable
\item $\lambda$: $-2$, $-1$, $4$; diagonalizable
\item $\lambda$: $-3$, $0$, $3$; diagonalizable
\item $\lambda$: $-2$ (double), $-1$; not diagonalizable
\item $\lambda$: $0$ (double), $1$; diagonalizable
\item $\lambda$: $-1$, $0$, $2$; diagonalizable
\item $\lambda$: $3$ (triple); not diagonalizable
\end{enumerate}
\end{multicols}
\par\noindent\emph{Exercise set 3.}\nopagebreak
\begin{enumerate}[label=\arabic*., leftmargin=*, itemsep=1pt]
\item $A^{4}=\left(\begin{smallmatrix}-14 & 30\\ -15 & 31\end{smallmatrix}\right)$
\item $A^{3}=\left(\begin{smallmatrix}-1 & 0\\ 0 & 8\end{smallmatrix}\right)$
\item $A^{4}=\left(\begin{smallmatrix}16 & 0\\ -15 & 1\end{smallmatrix}\right)$
\item $A^{4}=\left(\begin{smallmatrix}81 & -160\\ 0 & 1\end{smallmatrix}\right)$
\item $A^{4}=\left(\begin{smallmatrix}16 & 0\\ 0 & 1\end{smallmatrix}\right)$
\item $A^{5}=\left(\begin{smallmatrix}487 & -244\\ 488 & -245\end{smallmatrix}\right)$
\item $\lambda_{1,2}=\frac{1\pm\sqrt5}{2}$; $F^n=\left(\begin{smallmatrix}F_{n+1}&F_n\\F_n&F_{n-1}\end{smallmatrix}\right)$
\item $\lambda=2$ is a double eigenvalue, but the eigenspace is spanned by the single vector $(1;0)$
\item $Q=\frac{1}{\sqrt2}\left(\begin{smallmatrix}1&1\\-1&1\end{smallmatrix}\right)$, $D=\mathrm{diag}(2,4)$
\item $\lambda=2;5$: $2+5=7=\operatorname{tr}A$, $2\cdot5=10=\det A$
\end{enumerate}


\section*{Chapter 5. Quadratic Forms, Second-Order Curves and Iterative Methods}
% !Mode:: "TeX:UTF-8"
\javobmashgulot{21}{Quadratic forms and their reduction to canonical form}
\par\noindent\emph{Exercise set 1.}\nopagebreak
\begin{enumerate}[label=\arabic*., leftmargin=*, itemsep=1pt]
\item $\left(\begin{smallmatrix}-2 & 0 & 2\\ 0 & 2 & 2\\ 2 & 2 & 0\end{smallmatrix}\right)$
\item $\left(\begin{smallmatrix}3 & 3 & -2\\ 3 & 3 & 3\\ -2 & 3 & 4\end{smallmatrix}\right)$
\item $\left(\begin{smallmatrix}-2 & 1 & 1\\ 1 & -1 & 3\\ 1 & 3 & -4\end{smallmatrix}\right)$
\item $\left(\begin{smallmatrix}-4 & -1 & 1\\ -1 & 2 & -3\\ 1 & -3 & -2\end{smallmatrix}\right)$
\item $\left(\begin{smallmatrix}-1 & -2 & 2\\ -2 & -4 & 0\\ 2 & 0 & 2\end{smallmatrix}\right)$
\item $\left(\begin{smallmatrix}3 & -3 & -1\\ -3 & 4 & 0\\ -1 & 0 & -3\end{smallmatrix}\right)$
\item $\left(\begin{smallmatrix}1 & -3 & -2\\ -3 & -3 & -3\\ -2 & -3 & -4\end{smallmatrix}\right)$
\item $\left(\begin{smallmatrix}3 & -2 & 1\\ -2 & 1 & 0\\ 1 & 0 & 3\end{smallmatrix}\right)$
\item $\left(\begin{smallmatrix}-2 & -2 & 2\\ -2 & -4 & 0\\ 2 & 0 & -1\end{smallmatrix}\right)$
\item $\left(\begin{smallmatrix}-3 & 3 & -2\\ 3 & -1 & 2\\ -2 & 2 & -3\end{smallmatrix}\right)$
\item $\left(\begin{smallmatrix}-3 & 1 & 0\\ 1 & 4 & 3\\ 0 & 3 & 2\end{smallmatrix}\right)$
\item $\left(\begin{smallmatrix}-1 & -1 & -2\\ -1 & 4 & 0\\ -2 & 0 & 0\end{smallmatrix}\right)$
\item $\left(\begin{smallmatrix}1 & 0 & 1\\ 0 & 2 & -1\\ 1 & -1 & 3\end{smallmatrix}\right)$
\item $\left(\begin{smallmatrix}2 & -1 & -3\\ -1 & -3 & -2\\ -3 & -2 & 3\end{smallmatrix}\right)$
\item $\left(\begin{smallmatrix}-4 & 0 & -3\\ 0 & 4 & 1\\ -3 & 1 & 4\end{smallmatrix}\right)$
\item $\left(\begin{smallmatrix}-1 & 1 & 1\\ 1 & 1 & 1\\ 1 & 1 & -3\end{smallmatrix}\right)$
\item $\left(\begin{smallmatrix}-2 & 2 & -2\\ 2 & 1 & 3\\ -2 & 3 & 1\end{smallmatrix}\right)$
\item $\left(\begin{smallmatrix}1 & 3 & 1\\ 3 & 0 & -2\\ 1 & -2 & -3\end{smallmatrix}\right)$
\item $\left(\begin{smallmatrix}-2 & -3 & 3\\ -3 & -2 & 3\\ 3 & 3 & 2\end{smallmatrix}\right)$
\item $\left(\begin{smallmatrix}1 & -2 & 2\\ -2 & -4 & 2\\ 2 & 2 & 4\end{smallmatrix}\right)$
\item $\left(\begin{smallmatrix}-1 & 0 & -2\\ 0 & -4 & -1\\ -2 & -1 & 3\end{smallmatrix}\right)$
\item $\left(\begin{smallmatrix}4 & 0 & 2\\ 0 & -1 & -2\\ 2 & -2 & -2\end{smallmatrix}\right)$
\item $\left(\begin{smallmatrix}4 & -3 & 0\\ -3 & 2 & 3\\ 0 & 3 & 3\end{smallmatrix}\right)$
\item $\left(\begin{smallmatrix}-4 & 3 & -1\\ 3 & -4 & 3\\ -1 & 3 & 0\end{smallmatrix}\right)$
\item $\left(\begin{smallmatrix}0 & -3 & 0\\ -3 & -1 & -2\\ 0 & -2 & 2\end{smallmatrix}\right)$
\item $\left(\begin{smallmatrix}-3 & 3 & -1\\ 3 & -2 & -2\\ -1 & -2 & 3\end{smallmatrix}\right)$
\item $\left(\begin{smallmatrix}-4 & -1 & -2\\ -1 & 4 & 2\\ -2 & 2 & 3\end{smallmatrix}\right)$
\item $\left(\begin{smallmatrix}-4 & -2 & 3\\ -2 & 1 & 3\\ 3 & 3 & -2\end{smallmatrix}\right)$
\item $\left(\begin{smallmatrix}1 & 1 & 3\\ 1 & 2 & 0\\ 3 & 0 & 3\end{smallmatrix}\right)$
\item $\left(\begin{smallmatrix}0 & -2 & 1\\ -2 & 0 & -2\\ 1 & -2 & -4\end{smallmatrix}\right)$
\end{enumerate}
\par\noindent\emph{Exercise set 2.}\nopagebreak
\begin{multicols}{2}\raggedright
\begin{enumerate}[label=\arabic*., leftmargin=*, itemsep=1pt]
\item $y_1^2-4y_2^2$; indefinite
\item $-5y_1^2+5y_2^2$; indefinite
\item $4y_1^2+6y_2^2$; positive definite
\item $-4y_1^2+y_2^2$; indefinite
\item $4y_1^2-4y_2^2$; indefinite
\item $5y_1^2+y_2^2$; positive definite
\item $y_1^2-4y_2^2$; indefinite
\item $2y_1^2-2y_2^2$; indefinite
\item $y_1^2+6y_2^2$; positive definite
\item $6y_1^2+y_2^2$; positive definite
\item $3y_1^2-2y_2^2$; indefinite
\item $-5y_1^2+5y_2^2$; indefinite
\item $-4y_1^2+y_2^2$; indefinite
\item $-3y_1^2+3y_2^2$; indefinite
\item $-3y_1^2-5y_2^2$; negative definite
\item $-2y_1^2+3y_2^2$; indefinite
\item $5y_1^2-5y_2^2$; indefinite
\item $6y_1^2-4y_2^2$; indefinite
\item $3y_1^2-y_2^2$; indefinite
\item $5y_1^2-5y_2^2$; indefinite
\item $-3y_1^2+2y_2^2$; indefinite
\item $4y_1^2+6y_2^2$; positive definite
\item $-4y_1^2+6y_2^2$; indefinite
\item $-3y_1^2-y_2^2$; negative definite
\item $-4y_1^2+4y_2^2$; indefinite
\item $y_1^2+5y_2^2$; positive definite
\item $2y_1^2-2y_2^2$; indefinite
\item $2y_1^2-3y_2^2$; indefinite
\item $-3y_1^2+3y_2^2$; indefinite
\item $5y_1^2+y_2^2$; positive definite
\end{enumerate}
\end{multicols}
\par\noindent\emph{Exercise set 3.}\nopagebreak
\begin{multicols}{2}\raggedright
\begin{enumerate}[label=\arabic*., leftmargin=*, itemsep=1pt]
\item $\Delta_1=1,\ \Delta_2=4,\ \Delta_3=-9$; not positive definite
\item $\Delta_1=2,\ \Delta_2=-2,\ \Delta_3=-18$; not positive definite
\item $\Delta_1=5,\ \Delta_2=11,\ \Delta_3=30$; positive definite
\item $\Delta_1=4,\ \Delta_2=12,\ \Delta_3=32$; positive definite
\item $\Delta_1=5,\ \Delta_2=4,\ \Delta_3=-5$; not positive definite
\item $\Delta_1=3,\ \Delta_2=5,\ \Delta_3=12$; positive definite
\item $\Delta_1=3,\ \Delta_2=2,\ \Delta_3=-6$; not positive definite
\item $\Delta_1=4,\ \Delta_2=8,\ \Delta_3=22$; positive definite
\item $\Delta_1=3,\ \Delta_2=2,\ \Delta_3=7$; positive definite
\item $\Delta_1=1,\ \Delta_2=4,\ \Delta_3=-13$; not positive definite
\end{enumerate}
\end{multicols}

% !Mode:: "TeX:UTF-8"
\javobmashgulot{22}{Reducing a second-order equation to canonical form}
\par\noindent\emph{Exercise set 1.}\nopagebreak
\begin{multicols}{2}\raggedright
\begin{enumerate}[label=\arabic*., leftmargin=*, itemsep=1pt]
\item ellipse; $O^{\prime}(-2;-1)$, $\frac{X^2}{4}+\frac{Y^2}{16}=1$
\item hyperbola; $O^{\prime}(-2;-3)$, $\frac{X^2}{16}-\frac{Y^2}{1}=1$
\item parabola; $O^{\prime}(-1;0)$, $Y^2=-4X$
\item ellipse; $O^{\prime}(-2;4)$, $\frac{X^2}{25}+\frac{Y^2}{1}=1$
\item hyperbola; $O^{\prime}(-4;0)$, $\frac{X^2}{16}-\frac{Y^2}{16}=1$
\item parabola; $O^{\prime}(-1;-2)$, $Y^2=6X$
\item ellipse; $O^{\prime}(1;0)$, $\frac{X^2}{9}+\frac{Y^2}{16}=1$
\item hyperbola; $O^{\prime}(0;4)$, $\frac{X^2}{9}-\frac{Y^2}{9}=1$
\item parabola; $O^{\prime}(0;-2)$, $Y^2=6X$
\item ellipse; $O^{\prime}(4;-4)$, $\frac{X^2}{25}+\frac{Y^2}{9}=1$
\item hyperbola; $O^{\prime}(4;0)$, $\frac{X^2}{25}-\frac{Y^2}{16}=1$
\item parabola; $O^{\prime}(-2;0)$, $Y^2=-4X$
\item ellipse; $O^{\prime}(-4;0)$, $\frac{X^2}{25}+\frac{Y^2}{16}=1$
\item hyperbola; $O^{\prime}(-3;-3)$, $\frac{X^2}{25}-\frac{Y^2}{9}=1$
\item parabola; $O^{\prime}(-3;2)$, $Y^2=4X$
\item ellipse; $O^{\prime}(4;-1)$, $\frac{X^2}{25}+\frac{Y^2}{1}=1$
\item hyperbola; $O^{\prime}(-3;3)$, $\frac{X^2}{9}-\frac{Y^2}{4}=1$
\item parabola; $O^{\prime}(-4;4)$, $Y^2=4X$
\item ellipse; $O^{\prime}(1;1)$, $\frac{X^2}{9}+\frac{Y^2}{16}=1$
\item hyperbola; $O^{\prime}(-4;-3)$, $\frac{X^2}{9}-\frac{Y^2}{4}=1$
\item parabola; $O^{\prime}(-3;-4)$, $Y^2=2X$
\item ellipse; $O^{\prime}(-2;3)$, $\frac{X^2}{4}+\frac{Y^2}{1}=1$
\item hyperbola; $O^{\prime}(3;4)$, $\frac{X^2}{9}-\frac{Y^2}{9}=1$
\item parabola; $O^{\prime}(3;3)$, $Y^2=2X$
\item ellipse; $O^{\prime}(4;1)$, $\frac{X^2}{4}+\frac{Y^2}{9}=1$
\item hyperbola; $O^{\prime}(-3;-2)$, $\frac{X^2}{25}-\frac{Y^2}{9}=1$
\item parabola; $O^{\prime}(-2;-4)$, $Y^2=6X$
\item ellipse; $O^{\prime}(-3;-2)$, $\frac{X^2}{4}+\frac{Y^2}{9}=1$
\item hyperbola; $O^{\prime}(-3;3)$, $\frac{X^2}{9}-\frac{Y^2}{16}=1$
\item parabola; $O^{\prime}(1;-3)$, $Y^2=4X$
\end{enumerate}
\end{multicols}
\par\noindent\emph{Exercise set 2.}\nopagebreak
\begin{multicols}{2}\raggedright
\begin{enumerate}[label=\arabic*., leftmargin=*, itemsep=1pt]
\item $5x'^2+y'^2=20$; ellipse
\item $6x'^2+y'^2=8$; ellipse
\item $5x'^2-5y'^2=10$; hyperbola
\item $-x'^2-5y'^2=6$; empty set (imaginary ellipse)
\item $6x'^2+y'^2=20$; ellipse
\item $2x'^2-3y'^2=6$; hyperbola
\item $x'^2-4y'^2=20$; hyperbola
\item $3x'^2-2y'^2=6$; hyperbola
\item $3x'^2-3y'^2=12$; hyperbola
\item $5x'^2-5y'^2=4$; hyperbola
\item $4x'^2-4y'^2=6$; hyperbola
\item $3x'^2-2y'^2=8$; hyperbola
\item $5x'^2-y'^2=6$; hyperbola
\item $-2x'^2-4y'^2=6$; empty set (imaginary ellipse)
\item $3x'^2-y'^2=10$; hyperbola
\item $-x'^2-5y'^2=6$; empty set (imaginary ellipse)
\item $-3x'^2-5y'^2=12$; empty set (imaginary ellipse)
\item $3x'^2-y'^2=12$; hyperbola
\item $6x'^2+y'^2=12$; ellipse
\item $x'^2-y'^2=6$; hyperbola
\item $x'^2-3y'^2=8$; hyperbola
\item $-x'^2-5y'^2=6$; empty set (imaginary ellipse)
\item $2x'^2-3y'^2=6$; hyperbola
\item $6x'^2-4y'^2=8$; hyperbola
\item $2x'^2-4y'^2=6$; hyperbola
\item $6x'^2-4y'^2=12$; hyperbola
\item $5x'^2-5y'^2=4$; hyperbola
\item $3x'^2+y'^2=4$; ellipse
\item $x'^2-4y'^2=12$; hyperbola
\item $5x'^2-y'^2=4$; hyperbola
\end{enumerate}
\end{multicols}
\par\noindent\emph{Exercise set 3.}\nopagebreak
\begin{multicols}{2}\raggedright
\begin{enumerate}[label=\arabic*., leftmargin=*, itemsep=1pt]
\item $\delta=36>0$, $\Delta=-1296\ne0$, $I\Delta<0$ – ellipse
\item $O^{\prime}(1;0)$
\item $\operatorname{tg}2\alpha=\sqrt3$, $\alpha=30^\circ$
\item $C(3;-2)$, $R=5$
\item $F(\pm4;0)$, $e=\frac45$
\item $y=\pm\frac34x$, $e=\frac54$
\item $F(2;0)$, $x=-2$
\item a pair of intersecting lines $y=\pm x$
\item $y^2\equiv 8+4+3\equiv1$: $y=1$ or $y=6$
\item $\delta=0$, $\Delta=-16\ne0$ – parabola
\end{enumerate}
\end{multicols}

% !Mode:: "TeX:UTF-8"
\javobmashgulot{23}{The iterative method for the approximate solution of systems of linear algebraic equations}
\par\noindent\emph{Exercise set 1.}\nopagebreak
\begin{enumerate}[label=\arabic*., leftmargin=*, itemsep=1pt]
\item $x^{(1)}=(-1.875;\,1.6667;\,-0.5)$, $x^{(2)}=(-2.2708;\,1.0972;\,-0.0312)$; exact: $(-2;\,1;\,0)$
\item $x^{(1)}=(1.7273;\,-1.3333;\,0.1111)$, $x^{(2)}=(2.0606;\,-1.9091;\,0.835)$; exact: $(2;\,-2;\,1)$
\item $x^{(1)}=(1;\,2.5;\,-2)$, $x^{(2)}=(2.1667;\,1.875;\,-1.8125)$; exact: $(2;\,2;\,-2)$
\item $x^{(1)}=(-0.125;\,-0.2222;\,3.2222)$, $x^{(2)}=(0.7639;\,-0.9244;\,3.2886)$; exact: $(1;\,-1;\,3)$
\item $x^{(1)}=(-2.2222;\,-0.4;\,0.9)$, $x^{(2)}=(-1.9667;\,-1.0244;\,0.7978)$; exact: $(-2;\,-1;\,1)$
\item $x^{(1)}=(-2.8889;\,3.2222;\,0.5)$, $x^{(2)}=(-2.1173;\,2.3704;\,2.1111)$; exact: $(-2;\,3;\,2)$
\item $x^{(1)}=(1.4;\,-2.1667;\,3.75)$, $x^{(2)}=(0.9667;\,-2.0083;\,2.8583)$; exact: $(1;\,-2;\,3)$
\item $x^{(1)}=(-0.125;\,2.9;\,0)$, $x^{(2)}=(0.2375;\,2.8625;\,-0.9458)$; exact: $(0;\,3;\,-1)$
\item $x^{(1)}=(1.875;\,0.2;\,-4.5)$, $x^{(2)}=(3.5625;\,-1.075;\,-3.5625)$; exact: $(3;\,-1;\,-3)$
\item $x^{(1)}=(2.4286;\,-0.1667;\,3)$, $x^{(2)}=(3.3333;\,0.4286;\,2.1905)$; exact: $(3;\,0;\,2)$
\item $x^{(1)}=(-2.75;\,0.8571;\,1.4)$, $x^{(2)}=(-3.2464;\,0.2571;\,2.1214)$; exact: $(-3;\,0;\,2)$
\item $x^{(1)}=(0.5;\,2.25;\,-3.6667)$, $x^{(2)}=(-0.1111;\,3.0417;\,-3.2222)$; exact: $(0;\,3;\,-3)$
\item $x^{(1)}=(-0.3333;\,2.3333;\,2)$, $x^{(2)}=(0;\,3;\,2.0833)$; exact: $(0;\,3;\,2)$
\item $x^{(1)}=(-3.2222;\,-0.7143;\,0.7273)$, $x^{(2)}=(-3.0649;\,-1.0707;\,1.1183)$; exact: $(-3;\,-1;\,1)$
\item $x^{(1)}=(-2.3333;\,-1.7;\,0.1818)$, $x^{(2)}=(-2.2727;\,-0.9818;\,1.2818)$; exact: $(-2;\,-1;\,1)$
\item $x^{(1)}=(3.2222;\,-0.6667;\,-2.6667)$, $x^{(2)}=(2.2593;\,-2.3333;\,-3.6296)$; exact: $(2;\,-2;\,-3)$
\item $x^{(1)}=(-2.1667;\,-2.5;\,-3)$, $x^{(2)}=(-0.75;\,-3.4167;\,-2.1667)$; exact: $(-1;\,-3;\,-2)$
\item $x^{(1)}=(-0.3333;\,1.125;\,3.1429)$, $x^{(2)}=(-1.2817;\,-0.0536;\,2.7738)$; exact: $(-1;\,0;\,3)$
\item $x^{(1)}=(-3.9;\,-0.7143;\,2.5714)$, $x^{(2)}=(-2.9143;\,-1.0816;\,1.6612)$; exact: $(-3;\,-1;\,2)$
\item $x^{(1)}=(-0.125;\,0.3333;\,-2.875)$, $x^{(2)}=(0.9531;\,0.375;\,-2.9427)$; exact: $(1;\,0;\,-3)$
\item $x^{(1)}=(-1.1667;\,-3.5;\,1.75)$, $x^{(2)}=(-1.1667;\,-3.0625;\,2.0417)$; exact: $(-1;\,-3;\,2)$
\item $x^{(1)}=(-3;\,0.5;\,0.8571)$, $x^{(2)}=(-3;\,-0.2857;\,0.1429)$; exact: $(-3;\,0;\,0)$
\item $x^{(1)}=(0.25;\,0.2727;\,1)$, $x^{(2)}=(0.0682;\,0.0682;\,0.9369)$; exact: $(0;\,0;\,1)$
\item $x^{(1)}=(-2.6;\,3;\,-0.5)$, $x^{(2)}=(-3.3;\,2.4143;\,0.92)$; exact: $(-3;\,3;\,1)$
\item $x^{(1)}=(0.6364;\,2.75;\,-0.8)$, $x^{(2)}=(1.6045;\,2.7886;\,-1.9)$; exact: $(2;\,3;\,-2)$
\item $x^{(1)}=(-3.0909;\,-1.25;\,3)$, $x^{(2)}=(-2.2045;\,-2;\,3.526)$; exact: $(-2;\,-2;\,3)$
\item $x^{(1)}=(3.3333;\,-2.1429;\,-3)$, $x^{(2)}=(2.6905;\,-2.0952;\,-1.9286)$; exact: $(3;\,-2;\,-2)$
\item $x^{(1)}=(-0.25;\,1.4444;\,0.7143)$, $x^{(2)}=(-0.7897;\,1.6865;\,1.0556)$; exact: $(-1;\,2;\,1)$
\item $x^{(1)}=(2.3333;\,2.6667;\,-3.3)$, $x^{(2)}=(2.8556;\,3.2481;\,-3.1333)$; exact: $(3;\,3;\,-3)$
\item $x^{(1)}=(-3;\,1.75;\,-1.3333)$, $x^{(2)}=(-2.9074;\,1;\,-0.75)$; exact: $(-3;\,1;\,-1)$
\end{enumerate}
\par\noindent\emph{Exercise set 2.}\nopagebreak
\begin{multicols}{2}\raggedright
\begin{enumerate}[label=\arabic*., leftmargin=*, itemsep=1pt]
\item $\|B\|_\infty=\frac{5}{8}$; the method converges
\item $\|B\|_\infty=\frac{3}{5}$; the method converges
\item $\|B\|_\infty=2$; the sufficient condition fails
\item $\|B\|_\infty=\frac{4}{7}$; the method converges
\item $\|B\|_\infty=\frac{6}{11}$; the method converges
\item $\|B\|_\infty=3$; the sufficient condition fails
\item $\|B\|_\infty=\frac{6}{11}$; the method converges
\item $\|B\|_\infty=\frac{5}{9}$; the method converges
\item $\|B\|_\infty=\frac{7}{2}$; the sufficient condition fails
\item $\|B\|_\infty=\frac{1}{2}$; the method converges
\item $\|B\|_\infty=\frac{1}{2}$; the method converges
\item $\|B\|_\infty=\frac{8}{3}$; the sufficient condition fails
\item $\|B\|_\infty=\frac{4}{7}$; the method converges
\item $\|B\|_\infty=\frac{2}{3}$; the method converges
\item $\|B\|_\infty=7$; the sufficient condition fails
\item $\|B\|_\infty=\frac{4}{9}$; the method converges
\item $\|B\|_\infty=\frac{4}{7}$; the method converges
\item $\|B\|_\infty=6$; the sufficient condition fails
\item $\|B\|_\infty=\frac{4}{7}$; the method converges
\item $\|B\|_\infty=\frac{1}{2}$; the method converges
\item $\|B\|_\infty=\frac{7}{2}$; the sufficient condition fails
\item $\|B\|_\infty=\frac{2}{3}$; the method converges
\item $\|B\|_\infty=\frac{3}{5}$; the method converges
\item $\|B\|_\infty=6$; the sufficient condition fails
\item $\|B\|_\infty=\frac{4}{9}$; the method converges
\item $\|B\|_\infty=\frac{1}{2}$; the method converges
\item $\|B\|_\infty=7$; the sufficient condition fails
\item $\|B\|_\infty=\frac{2}{5}$; the method converges
\item $\|B\|_\infty=\frac{2}{5}$; the method converges
\item $\|B\|_\infty=4$; the sufficient condition fails
\end{enumerate}
\end{multicols}
\par\noindent\emph{Exercise set 3.}\nopagebreak
\begin{multicols}{2}\raggedright
\begin{enumerate}[label=\arabic*., leftmargin=*, itemsep=1pt]
\item $x^{(3)}=(2.2109;\,3.1875)$; exact $(2;\,3)$
\item $x^{(3)}=(2.9438;\,1.18)$; exact $(3;\,1)$
\item $x^{(3)}=(2.1125;\,-2.12)$; exact $(2;\,-2)$
\item $x^{(3)}=(-4.0514;\,-3.1469)$; exact $(-4;\,-3)$
\item $x^{(3)}=(3.6694;\,-2.3829)$; exact $(4;\,-3)$
\item $x^{(3)}=(-1.0062;\,2.9969)$; exact $(-1;\,3)$
\item $x^{(3)}=(-1.9986;\,-0.9973)$; exact $(-2;\,-1)$
\item $x^{(3)}=(-0.0816;\,2)$; exact $(0;\,2)$
\item $x^{(3)}=(3.94;\,-0.775)$; exact $(4;\,-1)$
\item $x^{(3)}=(3.7333;\,4.48)$; exact $(4;\,4)$
\end{enumerate}
\end{multicols}

% !Mode:: "TeX:UTF-8"
\javobmashgulot{24}{The Seidel method for the approximate solution of systems of linear algebraic equations}
\par\noindent\emph{Exercise set 1.}\nopagebreak
\begin{enumerate}[label=\arabic*., leftmargin=*, itemsep=1pt]
\item $x^{(1)}=(3.8571;\,2.8095;\,1.6667)$, $x^{(2)}=(3.1429;\,2.1349;\,1.9444)$; exact: $(3;\,2;\,2)$
\item $x^{(1)}=(-2;\,-3;\,1.375)$, $x^{(2)}=(-0.9625;\,-3.075;\,0.9859)$; exact: $(-1;\,-3;\,1)$
\item $x^{(1)}=(-1.5;\,0.5625;\,1.9531)$, $x^{(2)}=(-0.9427;\,1.0156;\,1.9967)$; exact: $(-1;\,1;\,2)$
\item $x^{(1)}=(1;\,2.25;\,-2.1667)$, $x^{(2)}=(1.1818;\,3.0398;\,-2.0316)$; exact: $(1;\,3;\,-2)$
\item $x^{(1)}=(-0.8182;\,-1.8455;\,3.3327)$, $x^{(2)}=(0.2911;\,-0.8461;\,2.911)$; exact: $(0;\,-1;\,3)$
\item $x^{(1)}=(1;\,0.8333;\,3.375)$, $x^{(2)}=(-0.0324;\,-0.0733;\,2.9878)$; exact: $(0;\,0;\,3)$
\item $x^{(1)}=(-1.75;\,-0.3056;\,-1.7396)$, $x^{(2)}=(-1.2387;\,-0.8626;\,-1.9485)$; exact: $(-1;\,-1;\,-2)$
\item $x^{(1)}=(-0.5455;\,-0.5124;\,1.991)$, $x^{(2)}=(0.0907;\,0.0231;\,2.0184)$; exact: $(0;\,0;\,2)$
\item $x^{(1)}=(-0.8;\,-1.16;\,-1.64)$, $x^{(2)}=(-0.24;\,-1.82;\,-1.9229)$; exact: $(0;\,-2;\,-2)$
\item $x^{(1)}=(1.3333;\,-0.7273;\,-3.0606)$, $x^{(2)}=(2.596;\,-1.9063;\,-2.9341)$; exact: $(3;\,-2;\,-3)$
\item $x^{(1)}=(0.7778;\,2.4333;\,1.5767)$, $x^{(2)}=(1.7641;\,2.8869;\,1.9179)$; exact: $(2;\,3;\,2)$
\item $x^{(1)}=(-3;\,-1.25;\,1.9375)$, $x^{(2)}=(-2.9297;\,-1.0342;\,2.009)$; exact: $(-3;\,-1;\,2)$
\item $x^{(1)}=(-2.125;\,-1.9688;\,2.9609)$, $x^{(2)}=(-2.001;\,-1.9998;\,2.9997)$; exact: $(-2;\,-2;\,3)$
\item $x^{(1)}=(-0.1;\,-1.475;\,-2.8031)$, $x^{(2)}=(1.0359;\,-1.0582;\,-2.9972)$; exact: $(1;\,-1;\,-3)$
\item $x^{(1)}=(-0.3333;\,2.9333;\,-0.0917)$, $x^{(2)}=(0.013;\,3.0209;\,0.0059)$; exact: $(0;\,3;\,0)$
\item $x^{(1)}=(-1;\,-1.25;\,-2.775)$, $x^{(2)}=(-1;\,-1.9438;\,-2.9831)$; exact: $(-1;\,-2;\,-3)$
\item $x^{(1)}=(-2.25;\,3.3889;\,0.9325)$, $x^{(2)}=(-0.8626;\,2.977;\,1.013)$; exact: $(-1;\,3;\,1)$
\item $x^{(1)}=(-2.7143;\,1.5143;\,2.0816)$, $x^{(2)}=(-3.0117;\,0.979;\,1.9967)$; exact: $(-3;\,1;\,2)$
\item $x^{(1)}=(3.375;\,-0.9062;\,0.1205)$, $x^{(2)}=(3.0502;\,-1.0176;\,0.0118)$; exact: $(3;\,-1;\,0)$
\item $x^{(1)}=(0.75;\,0.8333;\,0.9028)$, $x^{(2)}=(0.9844;\,0.9838;\,0.992)$; exact: $(1;\,1;\,1)$
\item $x^{(1)}=(-1.1429;\,1.7551;\,-1.7398)$, $x^{(2)}=(-1.8907;\,1.9688;\,-1.9668)$; exact: $(-2;\,2;\,-2)$
\item $x^{(1)}=(0.7143;\,-2.3333;\,1.1476)$, $x^{(2)}=(-0.1163;\,-1.9508;\,0.9749)$; exact: $(0;\,-2;\,1)$
\item $x^{(1)}=(-1.125;\,0.0179;\,-2.8357)$, $x^{(2)}=(-2.1839;\,0.9793;\,-3.0572)$; exact: $(-2;\,1;\,-3)$
\item $x^{(1)}=(-2.5714;\,1.25;\,-1.1825)$, $x^{(2)}=(-1.8764;\,0.9544;\,-0.9624)$; exact: $(-2;\,1;\,-1)$
\item $x^{(1)}=(0.3333;\,1.0952;\,0.0736)$, $x^{(2)}=(-0.0072;\,1.019;\,-0.0054)$; exact: $(0;\,1;\,0)$
\item $x^{(1)}=(2;\,2.1667;\,3)$, $x^{(2)}=(2.8611;\,2.9537;\,3)$; exact: $(3;\,3;\,3)$
\item $x^{(1)}=(-2;\,-1.6667;\,-3)$, $x^{(2)}=(-2.2222;\,-0.9259;\,-3.0741)$; exact: $(-2;\,-1;\,-3)$
\item $x^{(1)}=(-2.7778;\,3.5972;\,-2.2817)$, $x^{(2)}=(-1.8322;\,2.9086;\,-1.9499)$; exact: $(-2;\,3;\,-2)$
\item $x^{(1)}=(-0.75;\,-3.375;\,-0.9531)$, $x^{(2)}=(0.0938;\,-3.0176;\,-0.9817)$; exact: $(0;\,-3;\,-1)$
\item $x^{(1)}=(-3.3333;\,1.3;\,1.0381)$, $x^{(2)}=(-2.9;\,0.9624;\,1.0035)$; exact: $(-3;\,1;\,1)$
\end{enumerate}
\par\noindent\emph{Exercise set 2.}\nopagebreak
\begin{multicols}{2}\raggedright
\begin{enumerate}[label=\arabic*., leftmargin=*, itemsep=1pt]
\item $x^{(3)}=(-0.9975;\,1.001)$; exact $(-1;\,1)$
\item $x^{(3)}=(3;\,-1)$; exact $(3;\,-1)$
\item $x^{(3)}=(-0.9974;\,3.9996)$; exact $(-1;\,4)$
\item $x^{(3)}=(1.0394;\,-2.0197)$; exact $(1;\,-2)$
\item $x^{(3)}=(-2;\,-1)$; exact $(-2;\,-1)$
\item $x^{(3)}=(3;\,0)$; exact $(3;\,0)$
\item $x^{(3)}=(0.875;\,-4.0625)$; exact $(1;\,-4)$
\item $x^{(3)}=(3;\,1)$; exact $(3;\,1)$
\item $x^{(3)}=(0.036;\,-4.0135)$; exact $(0;\,-4)$
\item $x^{(3)}=(2;\,2)$; exact $(2;\,2)$
\item $x^{(3)}=(-4;\,0)$; exact $(-4;\,0)$
\item $x^{(3)}=(-1.94;\,-3.024)$; exact $(-2;\,-3)$
\item $x^{(3)}=(-0.838;\,-3.081)$; exact $(-1;\,-3)$
\item $x^{(3)}=(-1.0098;\,2.0025)$; exact $(-1;\,2)$
\item $x^{(3)}=(0;\,1)$; exact $(0;\,1)$
\item $x^{(3)}=(-2.0208;\,0.9844)$; exact $(-2;\,1)$
\item $x^{(3)}=(-2;\,-3)$; exact $(-2;\,-3)$
\item $x^{(3)}=(-2;\,0)$; exact $(-2;\,0)$
\item $x^{(3)}=(2.9985;\,2.0002)$; exact $(3;\,2)$
\item $x^{(3)}=(-4;\,0)$; exact $(-4;\,0)$
\item $x^{(3)}=(-2;\,-1)$; exact $(-2;\,-1)$
\item $x^{(3)}=(-1.9927;\,-3.9984)$; exact $(-2;\,-4)$
\item $x^{(3)}=(2.9089;\,-3.0342)$; exact $(3;\,-3)$
\item $x^{(3)}=(0.9999;\,-4)$; exact $(1;\,-4)$
\item $x^{(3)}=(-2.689;\,-3.8134)$; exact $(-3;\,-4)$
\item $x^{(3)}=(1.0086;\,-3.0037)$; exact $(1;\,-3)$
\item $x^{(3)}=(-4;\,-3)$; exact $(-4;\,-3)$
\item $x^{(3)}=(-4;\,0)$; exact $(-4;\,0)$
\item $x^{(3)}=(1.0703;\,3.9736)$; exact $(1;\,4)$
\item $x^{(3)}=(-1;\,-1)$; exact $(-1;\,-1)$
\end{enumerate}
\end{multicols}
\par\noindent\emph{Exercise set 3.}\nopagebreak
\begin{multicols}{2}\raggedright
\begin{enumerate}[label=\arabic*., leftmargin=*, itemsep=1pt]
\item Jacobi: $0.0366$; Seidel: $0.0048$
\item Jacobi: $0.0508$; Seidel: $0.001$
\item Jacobi: $0.0248$; Seidel: $0.0001$
\item Jacobi: $0.1332$; Seidel: $0.0017$
\item Jacobi: $0.0508$; Seidel: $0.0063$
\item Jacobi: $0.0317$; Seidel: $0.0027$
\item Jacobi: $0.0402$; Seidel: $0.0015$
\item Jacobi: $0.048$; Seidel: $0.0039$
\item Jacobi: $0.0918$; Seidel: $0.0098$
\item Jacobi: $0.0069$; Seidel: $0.0007$
\end{enumerate}
\end{multicols}

\normalsize
